\input{section_shapes}
\plugin{cube}{Cube intersection primitive}
\order{0}
\parameters{
    \parameter{toWorld}{\Transform\Or\Animation}{
       Specifies an optional linear object-to-world transformation.
       \default{none (i.e. object space $=$ world space)}
    }
    \parameter{flipNormals}{\Boolean}{
       Is the cube inverted, i.e. should the normal vectors
       be flipped? \default{\code{false}, i.e. the normals point outside}
    }
}

\renderings{
    \rendering{Basic example}
        {shape_cube_basic}
    \rendering{A textured and stretched cube with the default parameterization
     (Listing~\ref{lst:cube-example})}
        {shape_cube_parameterization}
}

This shape plugin describes a simple cube/cuboid intersection primitive. By
default, it creates a cube between the world-space positions $(-1, -1, -1)$ and $(1,1,1)$.
However, an arbitrary linear transformation may be specified to translate, rotate, scale
or skew it as desired. The parameterization of this shape maps every face onto the
rectangle $[0,1]^2$ in $uv$ space.
\vspace{5mm}
\begin{xml}[caption={Example of a textured and stretched cube}, label=lst:cube-example]
<shape type="cube">
  <transform name="toWorld">
    <scale z="2"/>
  </transform>

  <bsdf type="diffuse">
    <texture type="checkerboard" name="reflectance">
      <float name="uvscale" value="6"/>
    </texture>
  </bsdf>
</shape>
\end{xml}
\plugin{sphere}{Sphere intersection primitive}
\order{1}
\parameters{
    \parameter{center}{\Point}{
      Center of the sphere in object-space \default{(0, 0, 0)}
    }
    \parameter{radius}{\Float}{
      Radius of the sphere in object-space units \default{1}
    }
    \parameter{toWorld}{\Transform\Or\Animation}{
       Specifies an optional linear object-to-world transformation.
       Note that non-uniform scales are not permitted!
       \default{none (i.e. object space $=$ world space)}
    }
    \parameter{flipNormals}{\Boolean}{
       Is the sphere inverted, i.e. should the normal vectors
       be flipped? \default{\code{false}, i.e. the normals point outside}
    }
}

\renderings{
    \rendering{Basic example, see \lstref{sphere-basic}}
        {shape_sphere_basic}
    \rendering{A textured sphere with the default parameterization}
        {shape_sphere_parameterization}
}

This shape plugin describes a simple sphere intersection primitive. It should
always be preferred over sphere approximations modeled using triangles.

\begin{xml}[caption={A sphere can either be configured using a linear
\code{toWorld} transformation or the \code{center} and \code{radius} parameters (or both).
   The above two declarations are equivalent.}, label=lst:sphere-basic]
<shape type="sphere">
    <transform name="toWorld">
        <scale value="2"/>
        <translate x="1" y="0" z="0"/>
    </transform>
    <bsdf type="diffuse"/>
</shape>

<shape type="sphere">
    <point name="center" x="1" y="0" z="0"/>
    <float name="radius" value="2"/>
    <bsdf type="diffuse"/>
</shape>
\end{xml}
When a \pluginref{sphere} shape is turned into an \pluginref{area} light source,
Mitsuba switches to an efficient sampling strategy \cite{Shirley91Direct} that
has particularly low variance. This makes it a good default choice for lighting
new scenes (\figref{spherelight}).
\renderings{
    \rendering{Spherical area light modeled using triangles}
        {shape_sphere_arealum_tri}
    \rendering{Spherical area light modeled using the \pluginref{sphere} plugin}
        {shape_sphere_arealum_analytic}

    \caption{
        \label{fig:spherelight}
        Area lights built from the combination of the \pluginref{area}
        and \pluginref{sphere} plugins produce renderings that have an
        overall lower variance.
    }
}
\begin{xml}[caption=Instantiation of a sphere emitter]
<shape type="sphere">
    <point name="center" x="0" y="1" z="0"/>
    <float name="radius" value="1"/>

    <emitter type="area">
        <blackbody name="intensity" temperature="7000K"/>
    </emitter>
</shape>
\end{xml}
\plugin{cylinder}{Cylinder intersection primitive}
\order{2}
\parameters{
    \parameter{p0}{\Point}{
      Object-space starting point of the cylinder's centerline \default{(0, 0, 0)}
    }
    \parameter{p1}{\Point}{
      Object-space endpoint of the cylinder's centerline \default{(0, 0, 1)}
    }
    \parameter{radius}{\Float}{
      Radius of the cylinder in object-space units \default{1}
    }
    \parameter{flipNormals}{\Boolean}{
       Is the cylinder inverted, i.e. should the normal vectors
       be flipped? \default{\code{false}, i.e. the normals point outside}
    }
    \parameter{toWorld}{\Transform\Or\Animation}{
       Specifies an optional linear object-to-world transformation.
       Note that non-uniform scales are not permitted!
       \default{none (i.e. object space $=$ world space)}
    }
}
\renderings{
    \rendering{Cylinder with the default one-sided shading}
        {shape_cylinder_onesided}
    \rendering{Cylinder with two-sided shading, see \lstref{cylinder-twosided}}
        {shape_cylinder_twosided}
}
This shape plugin describes a simple cylinder intersection primitive.
It should always be preferred over approximations modeled using
triangles. Note that the cylinder does not have endcaps -- also,
it's interior has inward-facing normals, which most scattering
models in Mitsuba will treat as fully absorbing. If this is not
desirable, consider using the \pluginref{twosided} plugin.

\begin{xml}[caption={A simple example for instantiating a
cylinder, whose interior is visible}, label=lst:cylinder-twosided]
<shape type="cylinder">
    <float name="radius" value="0.3"/>
    <bsdf type="twosided">
        <bsdf type="diffuse"/>
    </bsdf>
</shape>
\end{xml}
\plugin{rectangle}{Rectangle intersection primitive}
\order{3}
\parameters{
    \parameter{toWorld}{\Transform\Or\Animation}{
       Specifies a linear object-to-world transformation.
       It is allowed to use non-uniform scaling, but no shear.
       \default{none (i.e. object space $=$ world space)}
    }
    \parameter{flipNormals}{\Boolean}{
       Is the rectangle inverted, i.e. should the normal vectors
       be flipped? \default{\code{false}}
    }
}
\renderings{
    \rendering{Two rectangles configured as a reflective surface and an
    emitter (\lstref{rectangle})}{shape_rectangle}
}

This shape plugin describes a simple rectangular intersection primitive.
It is mainly provided as a convenience for those cases when creating and
loading an external mesh with two triangles is simply too tedious, e.g.
when an area light source or a simple ground plane are needed.

By default, the rectangle covers the XY-range $[-1,1]\times[-1,1]$
and has a surface normal that points into the positive $Z$ direction.
To change the rectangle scale, rotation, or translation, use the
\code{toWorld} parameter.

\vspace{2mm}
\begin{xml}[caption={A simple example involving two rectangle instances}, label=lst:rectangle]
<scene version=$\MtsVer$>
    <shape type="rectangle">
        <bsdf type="diffuse"/>
    </shape>
    <shape type="rectangle">
        <transform name="toWorld">
            <rotate x="1" angle="90"/>
            <scale x="0.4" y="0.3" z="0.2"/>
            <translate y="1" z="0.2"/>
        </transform>
        <emitter type="area">
            <spectrum name="intensity" value="3"/>
        </emitter>
    </shape>
    <!-- ... other definitions ... -->
</scene>
\end{xml}
\plugin{disk}{Disk intersection primitive}
\order{4}
\parameters{
    \parameter{toWorld}{\Transform\Or\Animation}{
       Specifies a linear object-to-world transformation.
       Note that non-uniform scales are not permitted!
       \default{none (i.e. object space $=$ world space)}
    }
    \parameter{flipNormals}{\Boolean}{
       Is the disk inverted, i.e. should the normal vectors
       be flipped? \default{\code{false}}
    }
    \vspace{-8mm}
}
\renderings{
    \rendering{Rendering with an disk emitter and a textured disk, showing
    the default parameterization. (\lstref{disk})}{shape_disk}
}

\vspace{-1mm}
This shape plugin describes a simple disk intersection primitive. It is
usually preferable over discrete approximations made from triangles.

By default, the disk has unit radius and is located at the origin. Its
surface normal points into the positive $Z$ direction.
To change the disk scale, rotation, or translation, use the
\code{toWorld} parameter.

\begin{xml}[caption={A simple example involving two disk instances}, label=lst:disk]
<scene version=$\MtsVer$>
    <shape type="disk">
        <bsdf type="diffuse">
            <texture name="reflectance" type="checkerboard">
                <float name="uvscale" value="5"/>
            </texture>
        </bsdf>
    </shape>
    <shape type="disk">
        <transform name="toWorld">
            <rotate x="1" angle="90"/>
            <scale value="0.3"/>
            <translate y="1" z="0.3"/>
        </transform>
        <emitter type="area">
            <spectrum name="intensity" value="4"/>
        </emitter>
    </shape>
</scene>
\end{xml}
\plugin{obj}{Wavefront OBJ mesh loader}
\order{5}
\parameters{
    \parameter{filename}{\String}{
      Filename of the OBJ file that should be loaded
    }
    \parameter{faceNormals}{\Boolean}{
      When set to \code{true}, any existing or computed vertex normals are
      discarded and \emph{face normals} will instead be used during rendering.
      This gives the rendered object a faceted appearance.\default{\code{false}}
    }
    \parameter{maxSmoothAngle}{\Float}{
      When specified, Mitsuba will discard all vertex normals in the input mesh and rebuild
      them in a way that is sensitive to the presence of creases and corners. For more
      details on this parameter, see below. Disabled by default.
    }
    \parameter{flipNormals}{\Boolean}{
      Optional flag to flip all normals. \default{\code{false}, i.e.
      the normals are left unchanged}.
    }
    \parameter{flipTexCoords}{\Boolean}{
      Treat the vertical component of the texture as inverted? Most OBJ files use
      this convention. \default{\code{true}}
    }
    \parameter{toWorld}{\Transform\Or\Animation}{
       Specifies an optional linear object-to-world transformation.
       \default{none (i.e. object space $=$ world space)}
    }
    \parameter{shapeIndex}{\Integer}{
      When the file contains multiple meshes, this parameter can
      be used to select a single one. \default{\code{-1}, \mbox{i.e. load all}}
    }
    \parameter{collapse}{\Boolean}{
      Collapse all meshes into a single shape \default{\code{false}}
    }
    \parameter{loadMaterials}{\Boolean}{
      \mbox{Import materials from a \code{mtl} file, if it exists?\default{\code{true}}}
    }
}
\renderings{
    \label{fig:rungholt}
    \bigrendering{An example scene with both geometry and materials imported using the Wavefront OBJ mesh loader
    (Neu Rungholt model courtesy of \texttt{kescha}, converted from Minecraft to OBJ by Morgan McGuire)}{shape_obj}
}

This plugin implements a simple loader for Wavefront OBJ files. It handles
meshes containing triangles and quadrilaterals, and it also imports vertex normals
and texture coordinates.

Loading an ordinary OBJ file is as simple as writing:
\begin{xml}
<shape type="obj">
    <string name="filename" value="myShape.obj"/>
</shape>
\end{xml}
\paragraph{Material import:}
When the OBJ file references a Wavefront material description (a \code{.mtl} file),
Mitsuba attempts to reproduce the material within and associate it with the shape.
This is restricted to fairly basic materials and textures, hence in most cases
it will be preferable to override this behavior by specifying an
explicit Mitsuba BSDF that should be used instead.
This can be done by passing it as a child argument, e.g.
\begin{xml}
<shape type="obj">
    <string name="filename" value="myShape.obj"/>
    <bsdf type="roughplastic">
        <rgb name="diffuseReflectance" value="0.2, 0.6, 0.3"/>
    </bsdf>
</shape>
\end{xml}
The \code{mtl} material attributes that are automatically handled by Mitsuba include:
\begin{itemize}
\item Diffuse and glossy materials (optionally textured)
\item Smooth glass and metal
\item Textured transparency
\item Bump maps
\end{itemize}

In some cases, OBJ files contain \emph{multiple} objects with different associated
materials. In this case, the materials can be overwritten individually, by specifying
the corresponding names. For instance, if the OBJ file contains two materials named
\code{Glass} and \code{Water}, these can be overwritten as follows
\begin{xml}
<shape type="obj">
    <string name="filename" value="myShape.obj"/>
    <bsdf name="Glass" type="dielectric">
        <float name="intIOR" value="1.5"/>
    </bsdf>
    <bsdf name="Water" type="dielectric">
        <float name="intIOR" value="1.333"/>
    </bsdf>
</shape>
\end{xml}
\paragraph{The \code{maxSmoothAngle} parameter:}
\label{sec:maxSmoothAngle}
When given a mesh without vertex normals, Mitsuba will by default
create a smoothly varying normal field over the entire shape. This can produce
undesirable output when the input mesh contains regions that are intentionally
not smooth (i.e. corners, creases). Meshes that do include vertex
normals sometimes incorrectly interpolate normals over such regions, leading to
much the same problem.

The \code{maxSmoothAngle} parameter can be issued to force inspection of the dihedral angle associated with
each edge in the input mesh and disable normal interpolation locally where this angle exceeds
a certain threshold value. A reasonable value might be something like \code{30} (degrees).
The underlying analysis is somewhat costly and hence this parameter should only be used when it is
actually needed (i.e. when the mesh contains creases or edges and does not come with
valid vertex normals).

\remarks{
\item Importing geometry via OBJ files should only be used as an absolutely
last resort. Due to inherent limitations of this format, the files tend to be unreasonably
large, and parsing them requires significant amounts of memory and processing power. What's worse
is that the internally stored data is often truncated, causing a loss of precision.

If possible, use the \pluginref{ply} or \pluginref{serialized} plugins instead. For convenience, it
is also possible to convert legacy OBJ files into \code{.serialized} files using the \code{mtsimport}
utility. Using the resulting output will significantly accelerate the scene loading time.
}
\plugin{ply}{PLY (Stanford Triangle Format) mesh loader}
\order{6}
\parameters{
    \parameter{filename}{\String}{
      Filename of the PLY file that should be loaded
    }
    \parameter{faceNormals}{\Boolean}{
      When set to \code{true}, Mitsuba will use face normals when rendering
      the object, which will give it a faceted appearance. \default{\code{false}}
    }
    \parameter{maxSmoothAngle}{\Float}{
      When specified, Mitsuba will discard all vertex normals in the input mesh and rebuild
      them in a way that is sensitive to the presence of creases and corners. For more
      details on this parameter, see page~\pageref{sec:maxSmoothAngle}. Disabled by default.
    }
    \parameter{flipNormals}{\Boolean}{
      Optional flag to flip all normals. \default{\code{false}, i.e.
      the normals are left unchanged}.
    }
    \parameter{toWorld}{\Transform\Or\Animation}{
       Specifies an optional linear object-to-world transformation.
       \default{none (i.e. object space $=$ world space)}
    }
    \parameter{srgb}{\Boolean}{
      When set to \code{true}, any vertex colors will be interpreted as sRGB,
      instead of linear RGB \default{\code{true}}.
    }
}
\renderings{
    \rendering{The PLY plugin is useful for loading large geometry. (Dragon
        statue courtesy of XYZ RGB)}{shape_ply_dragon}
    \rendering{The Stanford bunny loaded with \code{faceNormals=true}. Note
        the faceted appearance.}{shape_ply_bunny}
}
This plugin implements a fast loader for the Stanford PLY format (both
the ASCII and binary format). It is based on the \code{libply} library by
Ares Lagae (\url{http://people.cs.kuleuven.be/~ares.lagae/libply}).
The current plugin implementation supports triangle meshes with optional
UV coordinates, vertex normals, and vertex colors.

When loading meshes that contain vertex colors, note that they need to be
explicitly referenced in a BSDF using a special texture named
\pluginref{vertexcolors}.
\plugin{serialized}{Serialized mesh loader}
\order{7}
\parameters{
    \parameter{filename}{\String}{
      Filename of the geometry file that should be loaded
    }
    \parameter{shapeIndex}{\Integer}{
        A \code{.serialized} file may contain several separate meshes. This parameter
        specifies which one should be loaded. \default{\code{0}, i.e. the first one}
    }
    \parameter{faceNormals}{\Boolean}{
      When set to \code{true}, any existing or computed vertex normals are
      discarded and \emph{face normals} will instead be used during rendering.
      This gives the rendered object a faceted appearance.\default{\code{false}}
    }
    \parameter{maxSmoothAngle}{\Float}{
      When specified, Mitsuba will discard all vertex normals in the input mesh and rebuild
      them in a way that is sensitive to the presence of creases and corners. For more
      details on this parameter, see page~\pageref{sec:maxSmoothAngle}. Disabled by default.
    }
    \parameter{flipNormals}{\Boolean}{
      Optional flag to flip all normals. \default{\code{false}, i.e.
      the normals are left unchanged}.
    }
    \parameter{toWorld}{\Transform\Or\Animation}{
       Specifies an optional linear object-to-world transformation.
       \default{none (i.e. object space $=$ world space)}
    }
}
The serialized mesh format represents the most space and time-efficient way
of getting geometry information into Mitsuba. It stores indexed triangle meshes
in a lossless gzip-based encoding that (after decompression) nicely matches up
with the internally used data structures. Loading such files is considerably
faster than the \pluginref{ply} plugin and orders of magnitude faster than
the \pluginref{obj} plugin. \vspace{-3mm}

\paragraph{Format description:}
The \code{serialized} file format uses the little endian encoding, hence
all fields below should be interpreted accordingly. The contents are structured as
follows:
\begin{center}
\begin{longtable}{>{\bfseries}p{2cm}p{11cm}}
\toprule
Type & Content\\
\midrule
\code{uint16}&   File format identifier: \ \  \code{0x041C}\\
\code{uint16}&   File version identifier. Currently set to \ \  \code{0x0004}\\
\midrule
\multicolumn{2}{|c|}{\emph{From this point on, the stream is
compressed by the \code{DEFLATE} algorithm.}}\\
\multicolumn{2}{|c|}{\emph{The used encoding is that of
the \code{zlib} library.}}\\
\midrule
\code{uint32}&An 32-bit integer whose bits can be used
to specify the following flags:\\[-4mm]
&\begin{description}
 \item[\code{0x0001}] The mesh data includes per-vertex normals\vspace{-2mm}
 \item[\code{0x0002}] The mesh data includes texture coordinates\vspace{-2mm}
 \item[\code{0x0008}] The mesh data includes vertex colors\vspace{-2mm}
 \item[\code{0x0010}] Use face normals instead of smothly interpolated vertex
 normals. Equivalent to specifying \code{faceNormals=true} to the plugin. \vspace{-6mm}
 \item[\code{0x1000}] The subsequent content is represented in single precision\vspace{-2mm}
 \item[\code{0x2000}] The subsequent content is represented in double precision
\end{description}\\[-4mm]
\code{string}&A null-terminated string (utf-8), which denotes the name of the shape.\\
\code{uint64}&Number of vertices in the mesh\\
\code{uint64}&Number of triangles in the mesh\\
\code{array}&Array of all vertex positions (X, Y, Z, X, Y, Z, ...)
specified in binary single or double precision format (as denoted by the
flags)\\
\code{array}&Array of all vertex normal directions (X, Y, Z, X, Y, Z, ...)
specified in binary single or double precision format. When the mesh has
no vertex normals, this field is omitted.\\
\code{array}&Array of all vertex texture coordinates (U, V, U, V, ...)
specified in binary single or double precision format. When the mesh has
no texture coordinates, this field is omitted.\\
\code{array}&Array of all vertex colors (R, G, B, R, G, B, ...)
specified in binary single or double precision format. When the mesh has
no vertex colors, this field is omitted.\\
\code{array}&Indexed triangle data (\code{[i1, i2, i3]}, \code{[i1, i2, i3]}, ..)
specified in \code{uint32} or in \code{uint64} format (the latter
is used when the number of vertices exceeds \code{0xFFFFFFFF}).\\
\bottomrule
\end{longtable}
\end{center}
\paragraph{Multiple shapes:}
It is possible to store multiple meshes in a single \code{.serialized}
file. This is done by simply concatenating their data streams,
where every one is structured according to the above description.
Hence, after each mesh, the stream briefly reverts back to an
uncompressed format, followed by an uncompressed header, and so on.
This is neccessary for efficient read access to arbitrary sub-meshes.

\paragraph{End-of-file dictionary:} In addition to the previous table,
a \code{.serialized} file also concludes with a brief summary at the end of
the file, which specifies the starting position of each sub-mesh:
\begin{center}
\begin{longtable}{>{\bfseries}p{2cm}p{11cm}}
\toprule
Type & Content\\
\midrule
\code{uint64}& File offset of the first mesh (in bytes)---this is always zero.\\
\code{uint64}& File offset of the second mesh\\
$\cdots$ & $\cdots$\\
\code{uint64}& File offset of the last sub-shape\\
\code{uint32}& Total number of meshes in the \code{.serialized} file\\
\bottomrule
\end{longtable}
\end{center}
\plugin{shapegroup}{Shape group for geometry instancing}
\order{8}
\parameters{
    \parameter{\Unnamed}{\Shape}{One or more shapes that should be
        made available for geometry instancing}
}

This plugin implements a container for shapes that should be
made available for geometry instancing. Any shapes placed in a
\pluginref{shapegroup} will not be visible on their own---instead, the
renderer will precompute ray intersection acceleration data structures
so that they can efficiently be referenced many times using the
\pluginref{instance} plugin. This is useful for rendering things like
forests, where only a few distinct types of trees have to be kept
in memory. An example is given below:

\vspace{5mm}
\begin{xml}[caption={An example of geometry instancing}, label=lst:instancing]
<!-- Declare a named shape group containing two objects -->
<shape type="shapegroup" id="myShapeGroup">
    <shape type="ply">
        <string name="filename" value="data.ply"/>
        <bsdf type="roughconductor"/>
    </shape>
    <shape type="sphere">
        <transform name="toWorld">
            <scale value="5"/>
            <translate y="20"/>
        </transform>
        <bsdf type="diffuse"/>
    </shape>
</shape>

<!-- Instantiate the shape group without any kind of transformation -->
<shape type="instance">
    <ref id="myShapeGroup"/>
</shape>

<!-- Create instance of the shape group, but rotated, scaled, and translated -->
<shape type="instance">
    <ref id="myShapeGroup"/>
    <transform name="toWorld">
        <rotate x="1" angle="45"/>
        <scale value="1.5"/>
        <translate z="10"/>
    </transform>
</shape>
\end{xml}
\plugin{instance}{Geometry instance}
\order{9}
\parameters{
    \parameter{\Unnamed}{\ShapeGroup}{A reference to a
    shape group that should be instantiated}
    \parameter{toWorld}{\Transform\Or\Animation}{
       Specifies an optional linear instance-to-world transformation.
       \default{none (i.e. instance space $=$ world space)}
    }
}
\renderings{
   \rendering{Surface viewed from the top}{shape_instance_fractal_top}
   \rendering{Surface viewed from the bottom}{shape_instance_fractal_bot}
   \caption{
      A visualization of a fractal surface by Irving and Segerman.
      (a 2D Gospel curve developed up to level 5 along the third
      dimension). This scene makes use of instancing to replicate
      similar structures to cheaply render a shape that effectively
      consists of several hundred millions of triangles.
   }
}

This plugin implements a geometry instance used to efficiently replicate
geometry many times. For details on how to create instances, refer to
the \pluginref{shapegroup} plugin.
\remarks{
  \item Note that it is \emph{not} possible to assign a different
   material to each instance --- the material assignment specified within
   the shape group is the one that matters.
  \item Shape groups cannot be used to replicate shapes with
  attached emitters, sensors, or subsurface scattering models.
}
\plugin{hair}{Hair intersection shape}
\order{11}
\parameters{
    \parameter{filename}{\String}{
      Filename of the hair data file that should be loaded
    }
    \parameter{radius}{\Float}{
      Radius of the hair segments in world-space units
      \default{0.025, which assumes that the scene
      is modeled in millimeters.}.
    }
    \parameter{angleThreshold}{\Float}{
      For performance reasons, the plugin will merge adjacent hair
      segments when the angle of their tangent directions is below
      than this value (in degrees). \default{1}.
    }
    \parameter{reduction}{\Float}{
      When the reduction ratio is set to a value between zero and one, the hair
      plugin stochastically culls this portion of the input data (where
      1 corresponds to removing all hairs). To approximately preserve the
      appearance in renderings, the hair radius is enlarged (see Cook et al.
      \cite{Cook2007Stochastic}). This parameter is convenient for fast
      previews. \default{0, i.e. all geometry is rendered}
    }
    \parameter{toWorld}{\Transform}{
       Specifies an optional linear object-to-world transformation.
       Note that non-uniform scales are not permitted!
       \default{none, i.e. object space $=$ world space}
    }
}
\renderings{
    \centering
    \fbox{\includegraphics[width=6cm]{images/shape_hair}}\hspace{4.5cm}
    \caption{A close-up of the hair shape rendered with a diffuse
    scattering model (an actual hair scattering model will
    be needed for realistic appearance)}
}
The plugin implements a space-efficient acceleration structure for
hairs made from many straight cylindrical hair segments with miter
joints. The underlying idea is that intersections with straight cylindrical
hairs can be found quite efficiently, and curved hairs are easily
approximated using a series of such segments.

The plugin supports two different input formats: a simple (but not
particularly efficient) ASCII format containing the coordinates of a
hair vertex on every line. An empty line marks the beginning of a
new hair. The following snippet is an example of this format:\newpage
\begin{xml}
.....
-18.5498 -21.7669 22.8138
-18.6358 -21.3581 22.9262
-18.7359 -20.9494 23.0256

-30.6367 -21.8369 6.78397
-30.7289 -21.4145 6.76688
-30.8226 -20.9933 6.73948
.....
\end{xml}

There is also a binary format, which starts with the identifier
``\texttt{BINARY\_HAIR}'' (11 bytes), followed by the number of
vertices as a 32-bit little endian integer.
The remainder of the file consists of the vertex positions stored as
single-precision XYZ coordinates (again in little-endian byte ordering).
To mark the beginning of a new hair strand, a single $+\infty$ floating
point value can be inserted between the vertex data.
\plugin{heightfield}{Height field intersection shape}
\order{11}
\parameters{
    \parameter{shadingNormals}{\Boolean}{
      Use linearly interpolated shading normals over the height
      field as opposed to discontinuous normals from the underlying
      bilinear patches? \default{\code{true}, i.e. interpolate smoothly varying normals}
    }
    \parameter{flipNormals}{\Boolean}{
      Optional flag to flip all normals. \default{\code{false}, i.e.
      the normals are left unchanged}.
    }
    \parameter{toWorld}{\Transform}{
       Specifies an optional linear object-to-world transformation.
       \default{none, i.e. object space $=$ world space}
    }
    \parameter{width, height}{\Integer}{
      When the nested texture is procedural (see below),
      this parameter specifies the resolution at which it should
      be rasterized to create a height field made of bilinear patches.
    }
    \parameter{scale}{\Float}{Scale factor that is applied to the height field
      values\default{No scaling, i.e. \code{1}}
    }
    \parameter{filename}{\String}{
      File name of an image file containing height field values. Alternatively,
      a nested texture can be provided (see below).
    }
    \parameter{\Unnamed}{\Texture}{
      A nested texture that specifies the height field values. This
      could be a bitmap-backed texture or one that is procedurally defined.
      In the latter case, it will be rasterized using the resolution specified
      by the \code{width} and \code{height} arguments.
    }
}
\vspace{-2mm}
\renderings{
    \rendering{Heigh field rendering of a mountain, see \lstref{heightfield-bitmap}}
        {shape_heightfield}
}
This plugin implements an efficient height field intersection shape, i.e.
a two-dimensional plane that is vertically displaced using height values
loaded from a texture.
Internally, the height field is represented as a min-max mipmap
\cite{Tevs2008Maximum}, allowing cheap storage and efficient ray
intersection queries. It is generally preferable to represent
height fields using this specialized plugin rather than converting
them into triangle meshes.

\begin{xml}[caption={Declaring a height field from a monochromatic scaled bitmap texture}, label=lst:heightfield-bitmap]
<shape type="heightfield">
    <string name="filename" value="mountain_profile.exr"/>
    <float name="scale" value="0.5"/>
</shape>
\end{xml}
\input{section_bsdf}
\plugin{diffuse}{Smooth diffuse material}
\order{1}
\icon{bsdf_diffuse}
\parameters{
    \parameter{reflectance}{\Spectrum\Or\Texture}{
      Specifies the diffuse albedo of the
      material \default{0.5}
    }
}

\renderings{
    \rendering{Homogeneous reflectance, see \lstref{diffuse-uniform}}
        {bsdf_diffuse_plain}
    \rendering{Textured reflectance, see \lstref{diffuse-textured}}
        {bsdf_diffuse_textured}
}

The smooth diffuse material (also referred to as ``Lambertian'')
represents an ideally diffuse material with a user-specified amount of
reflectance. Any received illumination is scattered so that the surface
looks the same independently of the direction of observation.

Apart from a  homogeneous reflectance value, the plugin can also accept
a nested or referenced texture map to be used as the source of reflectance
information, which is then mapped onto the shape based on its UV
parameterization. When no parameters are specified, the model uses the default
of 50\% reflectance.

Note that this material is one-sided---that is, observed from the
back side, it will be completely black. If this is undesirable,
consider using the \pluginref{twosided} BRDF adapter plugin.
\vspace{4mm}

\begin{xml}[caption={A diffuse material, whose reflectance is specified
    as an sRGB color}, label=lst:diffuse-uniform]
<bsdf type="diffuse">
    <srgb name="reflectance" value="#6d7185"/>
</bsdf>
\end{xml}

\begin{xml}[caption=A diffuse material with a texture map,
    label=lst:diffuse-textured]
<bsdf type="diffuse">
    <texture type="bitmap" name="reflectance">
        <string name="filename" value="wood.jpg"/>
    </texture>
</bsdf>
\end{xml}
\plugin{roughdiffuse}{Rough diffuse material}
\order{2}
\icon{bsdf_roughdiffuse}
\parameters{
    \parameter{reflectance}{\Spectrum\Or\Texture}{
        Specifies the diffuse albedo of the
        material. \default{0.5}
    }
    \parameter{alpha}{\Spectrum\Or\Texture}{
        Specifies the roughness of the unresolved surface micro-geometry
        using the \emph{root mean square} (RMS) slope of the
        microfacets. \default{0.2}
    }
    \parameter{useFastApprox}{\Boolean}{
        This parameter selects between the full version of the model
        or a fast approximation that still retains most qualitative features.
        \default{\texttt{false}, i.e. use the high-quality version}
    }
}

\renderings{
    \rendering{Smooth diffuse surface ($\alpha=0$)}
       {bsdf_roughdiffuse_0}
    \rendering{Very rough diffuse surface ($\alpha=0.7$)}
       {bsdf_roughdiffuse_0_7}
  \vspace{-3mm}
  \caption{The effect of switching from smooth to rough diffuse scattering
  is fairly subtle on this model---generally, there will be higher
  reflectance at grazing angles, as well as an overall reduced contrast.}\vspace{3mm}
}

This reflectance model describes the interaction of light with a \emph{rough}
diffuse material, such as plaster, sand, clay, or concrete, or ``powdery''
surfaces. The underlying theory was developed by Oren and Nayar
\cite{Oren1994Generalization}, who model the microscopic surface structure as
unresolved planar facets arranged in V-shaped grooves, where each facet
is an ideal diffuse reflector. The model takes into account shadowing,
masking, as well as interreflections between the facets.

Since the original publication, this approach has been shown to
be a good match for many real-world materials, particularly compared
to Lambertian scattering, which does not take surface roughness into account.

The implementation in Mitsuba uses a surface roughness parameter $\alpha$ that
is slightly different from the slope-area variance in the original 1994 paper.
The reason for this change is to make the parameter $\alpha$ portable
across different models (i.e. \pluginref{roughdielectric}, \pluginref{roughplastic},
\pluginref{roughconductor}).

To get an intuition about the effect of the
parameter $\alpha$, consider the following approximate classification:
a value of $\alpha=0.001-0.01$ corresponds to a material
with slight imperfections on an otherwise smooth surface (for such small
values, the model will behave identically to \pluginref{diffuse}), $\alpha=0.1$
is relatively rough, and $\alpha=0.3-0.7$ is \emph{extremely} rough
(e.g. an etched or ground surface).

Note that this material is one-sided---that is, observed from the
back side, it will be completely black. If this is undesirable,
consider using the \pluginref{twosided} BRDF adapter plugin.
\plugin{dielectric}{Smooth dielectric material}
\order{3}
\icon{bsdf_dielectric}
\parameters{
    \parameter{intIOR}{\Float\Or\String}{Interior index of refraction specified
     numerically or using a known material name. \default{\texttt{bk7} / 1.5046}}
    \parameter{extIOR}{\Float\Or\String}{Exterior index of refraction specified
     numerically or using a known material name. \default{\texttt{air} / 1.000277}}
    \parameter{specular\showbreak Reflectance}{\Spectrum\Or\Texture}{Optional
        factor that can be used to modulate the specular reflection component. Note
        that for physical realism, this parameter should never be touched. \default{1.0}}
    \parameter{specular\showbreak Transmittance}{\Spectrum\Or\Texture}{Optional
        factor that can be used to modulate the specular transmission component. Note
        that for physical realism, this parameter should never be touched. \default{1.0}}
}

\renderings{
    \medrendering{Air$\leftrightarrow$Water (IOR: 1.33) interface.
        See \lstref{dielectric-water}.}{bsdf_dielectric_water}
    \medrendering{Air$\leftrightarrow$Diamond (IOR: 2.419)}{bsdf_dielectric_diamond}
    \medrendering{Air$\leftrightarrow$Glass (IOR: 1.504) interface with absorption.
        See \lstref{dielectric-glass}.}{bsdf_dielectric_glass}
}

This plugin models an interface between two dielectric materials having mismatched
indices of refraction (for instance, water and air). Exterior and interior IOR values
can be specified independently, where ``exterior'' refers to the side that contains
the surface normal. When no parameters are given, the plugin activates the defaults, which
describe a borosilicate glass BK7/air interface.

In this model, the microscopic structure of the surface is assumed to be perfectly
smooth, resulting in a degenerate\footnote{Meaning that for any given incoming ray of light,
the model always scatters into a discrete set of directions, as opposed to a continuum.}
BSDF described by a Dirac delta distribution. For a similar model that instead describes a
rough surface microstructure, take a look at the \pluginref{roughdielectric} plugin.

\begin{xml}[caption=A simple air-to-water interface, label=lst:dielectric-water]
<shape type="...">
    <bsdf type="dielectric">
        <string name="intIOR" value="water"/>
        <string name="extIOR" value="air"/>
    </bsdf>
<shape>
\end{xml}

When using this model, it is crucial that the scene contains
meaningful and mutually compatible indices of refraction changes---see
\figref{glass-explanation} for a description of what this entails.

In many cases, we will want to additionally describe the \emph{medium} within a
dielectric material. This requires the use of a rendering technique that is
aware of media (e.g. the volumetric path tracer). An example of how one might
describe a slightly absorbing piece of glass is shown below:
\begin{xml}[caption=A glass material with absorption (based on the
   Beer-Lambert law). This material can only be used by an integrator
   that is aware of participating media., label=lst:dielectric-glass]
<shape type="...">
    <bsdf type="dielectric">
        <float name="intIOR" value="1.504"/>
        <float name="extIOR" value="1.0"/>
    </bsdf>

    <medium type="homogeneous" name="interior">
        <rgb name="sigmaS" value="0, 0, 0"/>
        <rgb name="sigmaA" value="4, 4, 2"/>
    </medium>
<shape>
\end{xml}
\vspace{1cm}

\begin{table}[h!]
    \centering
    \begin{tabular}{>{\ttfamily}p{5cm}r@{.}lp{.8cm}>{\ttfamily}p{5cm}r@{.}l}
        \toprule
        \rmfamily \textbf{Name} & \multicolumn{2}{l}{\textbf{Value}}& &
        \rmfamily \textbf{Name} & \multicolumn{2}{l}{\textbf{Value}}\\
        \cmidrule{1-3} \cmidrule{5-7}
        vacuum               & 1 & 0 &  &
        bromine              & 1 & 661\\
        helium               & 1 & 00004 & &
        water ice            & 1 & 31\\
        hydrogen             & 1 & 00013& &
        fused quartz         & 1 & 458\\[-.8mm]
        \cmidrule{5-7}\\[-5.5mm]
        air                  & 1 & 00028& &
        pyrex                & 1 & 470\\
        carbon dioxide       & 1 & 00045& &
        acrylic glass        & 1 & 49\\[-.8mm]
        \cmidrule{1-3}\\[-5.5mm]
        water                & 1 & 3330& &
        polypropylene        & 1 & 49\\
        acetone              & 1 & 36 & &
        bk7                  & 1 & 5046\\
        ethanol              & 1 & 361& &
        sodium chloride      & 1 & 544\\
        carbon tetrachloride & 1 & 461& &
        amber                & 1 & 55\\
        glycerol             & 1 & 4729& &
        pet                  & 1 & 575\\
        benzene              & 1 & 501& &
        diamond              & 2 & 419\\
        silicone oil         & 1 & 52045\\
        \bottomrule
    \end{tabular}
    \caption{
        \label{tbl:dielectric-iors}
         This table lists all supported material names along with
         along with their associated index of refraction at standard
         conditions. These material names can be used with the plugins
         \pluginref{dielectric},\
         \pluginref{roughdielectric},\
         \pluginref{plastic}, \
         \pluginref{roughplastic}, as well as
         \pluginref{coating}.
    }
\end{table}
\remarks{
 \item Dispersion is currently unsupported but will be enabled in a future release.
}
\plugin{thindielectric}{Thin dielectric material}
\order{4}
\parameters{
    \parameter{intIOR}{\Float\Or\String}{Interior index of refraction specified
     numerically or using a known material name. \default{\texttt{bk7} / 1.5046}}
    \parameter{extIOR}{\Float\Or\String}{Exterior index of refraction specified
     numerically or using a known material name. \default{\texttt{air} / 1.000277}}
    \parameter{specular\showbreak Reflectance}{\Spectrum\Or\Texture}{Optional
        factor that can be used to modulate the specular reflection component. Note
        that for physical realism, this parameter should never be touched. \default{1.0}}
    \parameter{specular\showbreak Transmittance}{\Spectrum\Or\Texture}{Optional
        factor that can be used to modulate the specular transmission component. Note
        that for physical realism, this parameter should never be touched. \default{1.0}}
}

This plugin models a \emph{thin} dielectric material that is embedded inside another
dielectric---for instance, glass surrounded by air. The interior of the material
is assumed to be so thin that its effect on transmitted rays is negligible,
Hence, light exits such a material without any form of angular deflection
(though there is still specular reflection).

This model should be used for things like glass windows that were modeled using only a
single sheet of triangles or quads. On the other hand, when the window consists of
proper closed geometry, \pluginref{dielectric} is the right choice. This is illustrated below:

\begin{figure}[h]
\setcounter{subfigure}{0}
\centering
\hfill
\subfloat[The \pluginref{dielectric} plugin models a single transition from one index of refraction to another]
    {\includegraphics[width=4.5cm]{images/bsdf_dielectric_figure.pdf}}\hfill
\subfloat[The \pluginref{thindielectric} plugin models a pair of interfaces causing a transient index of refraction change]
     {\includegraphics[width=4.5cm]{images/bsdf_thindielectric_figure.pdf}}\hfill
\subfloat[Windows modeled using a single sheet of geometry are the most frequent application of this BSDF]
     {\fbox{\includegraphics[width=4.5cm]{images/bsdf_thindielectric_window.jpg}}}\hspace*\fill
\caption{
    \label{fig:thindielectric-diff}
    An illustration of the difference between the \pluginref{dielectric} and \pluginref{thindielectric} plugins}
\end{figure}

The implementation correctly accounts for multiple internal reflections
inside the thin dielectric at \emph{no significant extra cost}, i.e. paths
of the type $R, TRT, TR^3T, ..$ for reflection and $TT, TR^2, TR^4T, ..$ for
refraction, where $T$ and $R$ denote individual reflection and refraction
events, respectively.
\plugin{roughdielectric}{Rough dielectric material}
\order{5}
\icon{bsdf_roughdielectric}
\parameters{
    \parameter{distribution}{\String}{
         Specifies the type of microfacet normal distribution
         used to model the surface roughness.
         \vspace{-1mm}
      \begin{enumerate}[(i)]
          \item \code{beckmann}: Physically-based distribution derived from
              Gaussian random surfaces. This is the default.\vspace{-1.5mm}
          \item \code{ggx}: The GGX \cite{Walter07Microfacet} distribution (also known as
              Trowbridge-Reitz \cite{Trowbridge19975Average} distribution)
              was designed to better approximate the long tails observed in measurements
              of ground surfaces, which are not modeled by the Beckmann distribution.
          \vspace{-1.5mm}
          \item \code{phong}: Anisotropic Phong distribution by
             Ashikhmin and Shirley \cite{Ashikhmin2005Anisotropic}.
             In most cases, the \code{ggx} and \code{beckmann} distributions
             should be preferred, since they provide better importance sampling
             and accurate shadowing/masking computations.
             \vspace{-4mm}
      \end{enumerate}
    }
    \parameter{alpha, alphaU, alphaV}{\Float\Or\Texture}{
        Specifies the roughness of the unresolved surface micro-geometry
        along the tangent and bitangent directions. When the Beckmann
        distribution is used, this parameter is equal to the
        \emph{root mean square} (RMS) slope of the microfacets.
        \code{alpha} is a convenience parameter to initialize both
        \code{alphaU} and \code{alphaV} to the same value. \default{0.1}.
    }
    \parameter{intIOR}{\Float\Or\String}{Interior index of refraction specified
        numerically or using a known material name. \default{\texttt{bk7} / 1.5046}}
    \parameter{extIOR}{\Float\Or\String}{Exterior index of refraction specified
        numerically or using a known material name. \default{\texttt{air} / 1.000277}}
    \parameter{sampleVisible}{\Boolean}{
        Enables a sampling technique proposed by Heitz and D'Eon~\cite{Heitz1014Importance},
        which focuses computation on the visible parts of the microfacet normal
        distribution, considerably reducing variance in some cases.
        \default{\code{true}, i.e. use visible normal sampling}
    }
    \parameter{specular\showbreak Reflectance,\newline
        specular\showbreak Transmittance}{\Spectrum\Or\Texture}{Optional
        factor that can be used to modulate the specular reflection/transmission component. Note
        that for physical realism, this parameter should never be touched. \default{1.0}}
}\vspace{4mm}

This plugin implements a realistic microfacet scattering model for rendering
rough interfaces between dielectric materials, such as a transition from air to
ground glass. Microfacet theory describes rough surfaces as an arrangement of
unresolved and ideally specular facets, whose normal directions are given by
a specially chosen \emph{microfacet distribution}. By accounting for shadowing
and masking effects between these facets, it is possible to reproduce the important
off-specular reflections peaks observed in real-world measurements of such
materials.
\renderings{
    \rendering{Anti-glare glass (Beckmann, $\alpha=0.02$)}
        {bsdf_roughdielectric_beckmann_0_0_2.jpg}
    \rendering{Rough glass (Beckmann, $\alpha=0.1$)}
        {bsdf_roughdielectric_beckmann_0_1.jpg}
}

This plugin is essentially the ``roughened'' equivalent of the (smooth) plugin
\pluginref{dielectric}. For very low values of $\alpha$, the two will
be identical, though scenes using this plugin will take longer to render
due to the additional computational burden of tracking surface roughness.

The implementation is based on the paper ``Microfacet Models
for Refraction through Rough Surfaces'' by Walter et al.
\cite{Walter07Microfacet}. It supports three different types of microfacet
distributions and has a texturable roughness parameter. Exterior and
interior IOR values can be specified independently, where ``exterior''
refers to the side that contains the surface normal. Similar to the
\pluginref{dielectric} plugin, IOR values can either be specified
numerically, or based on a list of known materials (see
\tblref{dielectric-iors} for an overview). When no parameters are given,
the plugin activates the default settings, which describe a borosilicate
glass BK7/air interface with a light amount of roughness modeled using a
Beckmann distribution.

To get an intuition about the effect of the surface roughness parameter
$\alpha$, consider the following approximate classification: a value of
$\alpha=0.001-0.01$ corresponds to a material with slight imperfections
on an otherwise smooth surface finish, $\alpha=0.1$ is relatively rough,
and $\alpha=0.3-0.7$ is \emph{extremely} rough (e.g. an etched or ground
finish). Values significantly above that are probably not too realistic.

Please note that when using this plugin, it is crucial that the scene contains
meaningful and mutually compatible index of refraction changes---see
\figref{glass-explanation} for an example of what this entails. Also, note that
the importance sampling implementation of this model is close, but
not always a perfect a perfect match to the underlying scattering distribution,
particularly for high roughness values and when the \texttt{ggx}
microfacet distribution is used. Hence, such renderings may
converge slowly.

\subsubsection*{Technical details}
All microfacet distributions allow the specification of two distinct
roughness values along the tangent and bitangent directions. This can be
used to provide a material with a ``brushed'' appearance. The alignment
of the anisotropy will follow the UV parameterization of the underlying
mesh. This means that such an anisotropic material cannot be applied to
triangle meshes that are missing texture coordinates.

Since Mitsuba 0.5.1, this plugin uses a new importance sampling technique
contributed by Eric Heitz and Eugene D'Eon, which restricts the sampling
domain to the set of visible (unmasked) microfacet normals. The previous
approach of sampling all normals is still available and can be enabled
by setting \code{sampleVisible} to \code{false}.
Note that this new method is only available for the \code{beckmann} and
\code{ggx} microfacet distributions. When the \code{phong} distribution
is selected, the parameter has no effect.

When rendering with the Phong microfacet distribution, a conversion is
used to turn the specified Beckmann-equivalent $\alpha$ roughness value
into the exponent parameter of this distribution. This is done in a way,
such that the same value $\alpha$ will produce a similar appearance across
different microfacet distributions.

When rendering with the Phong microfacet distribution, a conversion is
used to turn the specified Beckmann-equivalent $\alpha$ roughness value
into the exponents of the distribution. This is done in a way, such that
the different distributions all produce a similar appearance for the
same value of $\alpha$.

\renderings{
    \rendering{Ground glass (GGX, $\alpha$=0.304,
        \lstref{roughdielectric-roughglass})}{bsdf_roughdielectric_ggx_0_304.jpg}
    \rendering{Textured roughness (\lstref{roughdielectric-textured})}
        {bsdf_roughdielectric_textured.jpg}
}

\begin{xml}[caption=A material definition for ground glass, label=lst:roughdielectric-roughglass]
<bsdf type="roughdielectric">
    <string name="distribution" value="ggx"/>
    <float name="alpha" value="0.304"/>
    <string name="intIOR" value="bk7"/>
    <string name="extIOR" value="air"/>
</bsdf>
\end{xml}

\begin{xml}[caption=A texture can be attached to the roughness parameter, label=lst:roughdielectric-textured]
<bsdf type="roughdielectric">
    <string name="distribution" value="beckmann"/>
    <float name="intIOR" value="1.5046"/>
    <float name="extIOR" value="1.0"/>

    <texture name="alpha" type="bitmap">
        <string name="filename" value="roughness.exr"/>
    </texture>
</bsdf>
\end{xml}
\plugin{conductor}{Smooth conductor}
\order{6}
\icon{bsdf_conductor}
\parameters{
    \parameter{material}{\String}{Name of a material preset, see
          \tblref{conductor-iors}.\!\default{\texttt{Cu} / copper}}
    \parameter{eta, k}{\Spectrum}{Real and imaginary components of the material's index of
            refraction \default{based on the value of \texttt{material}}}
    \parameter{extEta}{\Float\Or\String}{
          Real-valued index of refraction of the surrounding dielectric,
          or a material name of a dielectric \default{\code{air}}
    }
    \parameter{specular\showbreak Reflectance}{\Spectrum\Or\Texture}{Optional
        factor that can be used to modulate the specular reflection component. Note
        that for physical realism, this parameter should never be touched. \default{1.0}}
}
\renderings{
    \rendering{Measured copper material (the default), rendered using 30
    spectral samples between 360 and 830$nm$}
        {bsdf_conductor_copper.jpg}
    \rendering{Measured gold material (\lstref{conductor-gold})}
        {bsdf_conductor_gold.jpg}
}

This plugin implements a perfectly smooth interface to a conducting material,
such as a metal. For a similar model that instead describes a rough surface
microstructure, take a look at the separately available
\pluginref{roughconductor} plugin.

In contrast to dielectric materials, conductors do not transmit
any light. Their index of refraction is complex-valued and tends to undergo
considerable changes throughout the visible color spectrum.

To facilitate the tedious task of specifying spectrally-varying index of
refraction information, Mitsuba ships with a set of measured data for
several materials, where visible-spectrum information was publicly
available\footnote{
  These index of refraction values are identical to the data distributed
  with PBRT. They are originally from the Luxpop database
  (\url{www.luxpop.com}) and are based on data by Palik et al.
  \cite{Palik1998Handbook} and measurements of atomic scattering factors
  made by the Center For X-Ray Optics (CXRO) at Berkeley and the
  Lawrence Livermore National Laboratory (LLNL).
}.

Note that \tblref{conductor-iors} also includes several popular optical
coatings, which are not actually conductors. These materials can also
be used with this plugin, though note that the plugin will ignore any
refraction component that the actual material might have had.
There is also a special material profile named \code{none}, which disables
the computation of Fresnel reflectances and produces an idealized
100\% reflecting mirror.

When using this plugin, you should ideally compile Mitsuba with support for
spectral rendering to get the most accurate results. While it also works
in RGB mode, the computations will be more approximate in nature.
Also note that this material is one-sided---that is, observed from the
back side, it will be completely black. If this is undesirable,
consider using the \pluginref{twosided} BRDF adapter plugin.\vspace{4mm}

\begin{xml}[caption=A material configuration for a smooth conductor with
   measured gold data, label=lst:conductor-gold]
<shape type="...">
    <bsdf type="conductor">
        <string name="material" value="Au"/>
    </bsdf>
<shape>
\end{xml}
\vspace{5mm}
It is also possible to load spectrally varying index of refraction data from
two external files containing the real and imaginary components,
respectively (see \secref{format-spectra} for details on the file
format):
\begin{xml}[caption=Rendering a smooth conductor with custom data]
<shape type="...">
    <bsdf type="conductor">
        <spectrum name="eta" filename="conductorIOR.eta.spd"/>
        <spectrum name="k" filename="conductorIOR.k.spd"/>
    </bsdf>
<shape>
\end{xml}
\vspace{1.5cm}
\begin{table}[hb!]
\centering
\scriptsize
\begin{tabular}{>{\ttfamily}llp{1mm}>{\ttfamily}ll}
\toprule
\rmfamily\textbf{Preset(s)} & \textbf{Description} &&
\rmfamily\textbf{Preset(s)} & \textbf{Description}\\
\cmidrule{1-2} \cmidrule{4-5}
a-C                   & Amorphous carbon &&            Na\_palik             & Sodium \\
Ag                    & Silver &&                      Nb, Nb\_palik         & Niobium \\
Al                    & Aluminium &&                   Ni\_palik             & Nickel \\
AlAs, AlAs\_palik     & Cubic aluminium arsenide &&    Rh, Rh\_palik         & Rhodium \\
AlSb, AlSb\_palik     & Cubic aluminium antimonide &&  Se, Se\_palik         & Selenium \\
Au                    & Gold &&                        SiC, SiC\_palik       & Hexagonal silicon carbide \\
Be, Be\_palik         & Polycrystalline beryllium &&   SnTe, SnTe\_palik     & Tin telluride\\
Cr                    & Chromium &&                    Ta, Ta\_palik         & Tantalum \\
CsI, CsI\_palik       & Cubic caesium iodide &&        Te, Te\_palik         & Trigonal tellurium \\
Cu, Cu\_palik         & Copper &&                      ThF4, ThF4\_palik     & Polycryst. thorium (IV) fluoride \\
Cu2O, Cu2O\_palik     & Copper (I) oxide &&            TiC, TiC\_palik       & Polycrystalline titanium carbide \\
CuO, CuO\_palik       & Copper (II) oxide &&           TiN, TiN\_palik       & Titanium nitride \\
d-C, d-C\_palik       & Cubic diamond &&               TiO2, TiO2\_palik     & Tetragonal titan. dioxide\\
Hg, Hg\_palik         & Mercury &&                     VC, VC\_palik         & Vanadium carbide \\
HgTe, HgTe\_palik     & Mercury telluride &&           V\_palik              & Vanadium \\
Ir, Ir\_palik         & Iridium &&                     VN, VN\_palik         & Vanadium nitride \\
K, K\_palik           & Polycrystalline potassium &&   W                     & Tungsten \\
Li, Li\_palik         & Lithium &&                     \\
MgO, MgO\_palik       & Magnesium oxide &&             \\
Mo, Mo\_palik         & Molybdenum &&                  none & No mat. profile ($\to$ 100\% reflecting mirror)\\
\bottomrule
\end{tabular}
\caption{
    \label{tbl:conductor-iors}
     This table lists all supported materials that can be passed into the
     \pluginref{conductor} and \pluginref{roughconductor} plugins. Note that
     some of them are not actually conductors---this is not a problem,
     they can be used regardless (though only the reflection component and
     no transmission will be simulated). In most cases, there are
     multiple entries for each material, which represent measurements by
     different authors.
}
\end{table}
\plugin{roughconductor}{Rough conductor material}
\order{7}
\icon{bsdf_roughconductor}
\parameters{
    \parameter{distribution}{\String}{
         Specifies the type of microfacet normal distribution
         used to model the surface roughness.
         \vspace{-1mm}
      \begin{enumerate}[(i)]
          \item \code{beckmann}: Physically-based distribution derived from
              Gaussian random surfaces. This is the default.\vspace{-1.5mm}
          \item \code{ggx}: The GGX \cite{Walter07Microfacet} distribution (also known as
              Trowbridge-Reitz \cite{Trowbridge19975Average} distribution)
              was designed to better approximate the long tails observed in measurements
              of ground surfaces, which are not modeled by the Beckmann distribution.
          \vspace{-1.5mm}
          \item \code{phong}: Anisotropic Phong distribution by
             Ashikhmin and Shirley \cite{Ashikhmin2005Anisotropic}.
             In most cases, the \code{ggx} and \code{beckmann} distributions
             should be preferred, since they provide better importance sampling
             and accurate shadowing/masking computations.
             \vspace{-4mm}
      \end{enumerate}
    }
    \parameter{alpha, alphaU, alphaV}{\Float\Or\Texture}{
        Specifies the roughness of the unresolved surface micro-geometry
        along the tangent and bitangent directions. When the Beckmann
        distribution is used, this parameter is equal to the
        \emph{root mean square} (RMS) slope of the microfacets.
        \code{alpha} is a convenience parameter to initialize both
        \code{alphaU} and \code{alphaV} to the same value. \default{0.1}.
    }
    \parameter{material}{\String}{Name of a material preset, see
          \tblref{conductor-iors}.\!\default{\texttt{Cu} / copper}}
    \parameter{eta, k}{\Spectrum}{Real and imaginary components of the material's index of
            refraction \default{based on the value of \texttt{material}}}
    \parameter{extEta}{\Float\Or\String}{
          Real-valued index of refraction of the surrounding dielectric,
          or a material name of a dielectric \default{\code{air}}
    }
    \parameter{sampleVisible}{\Boolean}{
        Enables a sampling technique proposed by Heitz and D'Eon~\cite{Heitz1014Importance},
        which focuses computation on the visible parts of the microfacet normal
        distribution, considerably reducing variance in some cases.
        \default{\code{true}, i.e. use visible normal sampling}
    }
    \parameter{specular\showbreak Reflectance}{\Spectrum\Or\Texture}{Optional
        factor that can be used to modulate the specular reflection component. Note
        that for physical realism, this parameter should never be touched. \default{1.0}}
}
\vspace{3mm}
This plugin implements a realistic microfacet scattering model for rendering
rough conducting materials, such as metals. It can be interpreted as a fancy
version of the Cook-Torrance model and should be preferred over
heuristic models like \pluginref{phong} and \pluginref{ward} if possible.
\renderings{
    \rendering{Rough copper (Beckmann, $\alpha=0.1$)}
        {bsdf_roughconductor_copper.jpg}
    \rendering{Vertically brushed aluminium (Anisotropic Phong,
        $\alpha_u=0.05,\ \alpha_v=0.3$), see
        \lstref{roughconductor-aluminium}}
        {bsdf_roughconductor_anisotropic_aluminium.jpg}
}

Microfacet theory describes rough surfaces as an arrangement of unresolved
and ideally specular facets, whose normal directions are given by a
specially chosen \emph{microfacet distribution}. By accounting for shadowing
and masking effects between these facets, it is possible to reproduce the
important off-specular reflections peaks observed in real-world measurements
of such materials.

This plugin is essentially the ``roughened'' equivalent of the (smooth) plugin
\pluginref{conductor}. For very low values of $\alpha$, the two will
be identical, though scenes using this plugin will take longer to render
due to the additional computational burden of tracking surface roughness.

The implementation is based on the paper ``Microfacet Models
for Refraction through Rough Surfaces'' by Walter et al.
\cite{Walter07Microfacet}. It supports three different types of microfacet
distributions and has a texturable roughness parameter.
To facilitate the tedious task of specifying spectrally-varying index of
refraction information, this plugin can access a set of measured materials
for which visible-spectrum information was publicly available
(see \tblref{conductor-iors} for the full list).
There is also a special material profile named \code{none}, which disables
the computation of Fresnel reflectances and produces an idealized
100\% reflecting mirror.

When no parameters are given, the plugin activates the default settings,
which describe copper with a medium amount of roughness modeled using a
Beckmann distribution.

To get an intuition about the effect of the surface roughness parameter
$\alpha$, consider the following approximate classification: a value of
$\alpha=0.001-0.01$ corresponds to a material with slight imperfections
on an otherwise smooth surface finish, $\alpha=0.1$ is relatively rough,
and $\alpha=0.3-0.7$ is \emph{extremely} rough (e.g. an etched or ground
finish). Values significantly above that are probably not too realistic.
\vspace{4mm}
\begin{xml}[caption={A material definition for brushed aluminium}, label=lst:roughconductor-aluminium]
<bsdf type="roughconductor">
    <string name="material" value="Al"/>
    <string name="distribution" value="phong"/>
    <float name="alphaU" value="0.05"/>
    <float name="alphaV" value="0.3"/>
</bsdf>
\end{xml}

\subsubsection*{Technical details}
All microfacet distributions allow the specification of two distinct
roughness values along the tangent and bitangent directions. This can be
used to provide a material with a ``brushed'' appearance. The alignment
of the anisotropy will follow the UV parameterization of the underlying
mesh. This means that such an anisotropic material cannot be applied to
triangle meshes that are missing texture coordinates.

\label{sec:visiblenormal-sampling}
Since Mitsuba 0.5.1, this plugin uses a new importance sampling technique
contributed by Eric Heitz and Eugene D'Eon, which restricts the sampling
domain to the set of visible (unmasked) microfacet normals. The previous
approach of sampling all normals is still available and can be enabled
by setting \code{sampleVisible} to \code{false}.
Note that this new method is only available for the \code{beckmann} and
\code{ggx} microfacet distributions. When the \code{phong} distribution
is selected, the parameter has no effect.

When rendering with the Phong microfacet distribution, a conversion is
used to turn the specified Beckmann-equivalent $\alpha$ roughness value
into the exponent parameter of this distribution. This is done in a way,
such that the same value $\alpha$ will produce a similar appearance across
different microfacet distributions.

When using this plugin, you should ideally compile Mitsuba with support for
spectral rendering to get the most accurate results. While it also works
in RGB mode, the computations will be more approximate in nature.
Also note that this material is one-sided---that is, observed from the
back side, it will be completely black. If this is undesirable,
consider using the \pluginref{twosided} BRDF adapter.
\plugin{plastic}{Smooth plastic material}
\order{8}
\icon{bsdf_plastic}
\parameters{
    \parameter{intIOR}{\Float\Or\String}{Interior index of refraction specified
     numerically or using a known material name. \default{\texttt{polypropylene} / 1.49}}
    \parameter{extIOR}{\Float\Or\String}{Exterior index of refraction specified
     numerically or using a known material name. \default{\texttt{air} / 1.000277}}
    \parameter{specular\showbreak Reflectance}{\Spectrum\Or\Texture}{Optional
        factor that can be used to modulate the specular reflection component. Note that
        for physical realism, this parameter should never be touched. \default{1.0}}
    \parameter{diffuse\showbreak Reflectance}{\Spectrum\Or\Texture}{Optional
        factor used to modulate the diffuse reflection component\default{0.5}}
    \parameter{nonlinear}{\Boolean}{
        Account for nonlinear color shifts due to internal scattering? See the
        main text for details.\default{Don't account for them and
        preserve the texture colors, i.e. \code{false}}
    }
}

\renderings{
    \rendering{A rendering with the default parameters}{bsdf_plastic_default}
    \rendering{A rendering with custom parameters (\lstref{plastic-shiny})}
        {bsdf_plastic_shiny}
}

\vspace{3mm}
This plugin describes a smooth plastic-like material with internal scattering.
It uses the Fresnel reflection and transmission coefficients to provide
direction-dependent specular and diffuse components.
Since it is simple, realistic, and fast, this model is often a better choice
than the \pluginref{phong}, \pluginref{ward}, and \pluginref{roughplastic}
plugins when rendering smooth plastic-like materials.

For convenience, this model allows to specify IOR values either numerically,
or based on a list of known materials (see \tblref{dielectric-iors} for
an overview).

Note that this plugin is quite similar to what one would get by applying the
\pluginref{coating} plugin to the \pluginref{diffuse} material. The main
difference is that this plugin is significantly faster, while at the same
time causing less variance. Furthermore, it accounts for multiple
interreflections inside the material (read on for details), which avoids
a serious energy loss problem of the aforementioned plugin
combination.
\newpage

\begin{xml}[caption=A shiny material whose diffuse reflectance is
    specified using sRGB, label=lst:plastic-shiny]
<bsdf type="plastic">
    <srgb name="diffuseReflectance" value="#18455c"/>
    <float name="intIOR" value="1.9"/>
</bsdf>
\end{xml}

\renderings{
    \medrendering{Diffuse textured rendering}{bsdf_plastic_diffuse}
    \medrendering{Plastic model, \code{nonlinear=false}}{bsdf_plastic_preserve}
    \medrendering{Plastic model, \code{nonlinear=true}}{bsdf_plastic_nopreserve}
    \caption{
       \label{fig:plastic-nonlinear}
       When asked to do so, this model can account for subtle nonlinear color shifts due
       to internal scattering processes. The above images show a textured
       object first rendered using \pluginref{diffuse}, then
       \pluginref{plastic} with the default parameters, and finally using
       \pluginref{plastic} and support for nonlinear color shifts.
    }
}

\subsubsection*{Internal scattering}
Internally, this is model simulates the interaction of light with a diffuse
base surface coated by a thin dielectric layer. This is a convenient
abstraction rather than a restriction. In other words, there are many
materials that can be rendered with this model, even if they might not
fit this description perfectly well.

\begin{figure}[h]
\setcounter{subfigure}{0}
\centering
\subfloat[At the boundary, incident illumination is partly \mbox{reflected} and refracted]
     {\includegraphics[width=4.9cm]{images/bsdf_plastic_intscat_1.pdf}}\hfill
\subfloat[The refracted portion scatters diffusely at the base layer]
    {\includegraphics[width=4.9cm]{images/bsdf_plastic_intscat_2.pdf}}\hfill
\subfloat[Some of the illumination undergoes further internal scattering events]
    {\includegraphics[width=4.9cm]{images/bsdf_plastic_intscat_3.pdf}}
\caption{
    \label{fig:plastic-intscat}
    An illustration of the scattering events that are internally
    handled by this plugin}
\end{figure}

Given illumination that is incident upon such a material, a portion
of the illumination is specularly reflected at the material
boundary, which results in a sharp reflection in the mirror direction
(\subfigref{plastic-intscat}{a}).
The remaining illumination refracts into the material, where it
scatters from the diffuse base layer. (\subfigref{plastic-intscat}{b}).
While some of the diffusely scattered illumination is able to
directly refract outwards again, the remainder is reflected from the
interior side of the dielectric boundary and will in fact remain
trapped inside the material for some number of internal scattering
events until it is finally able to escape (\subfigref{plastic-intscat}{c}).

Due to the mathematical simplicity of this setup, it is possible to work
out the correct form of the model without actually having to simulate
the potentially large number of internal scattering events.

Note that due to the internal scattering, the diffuse color of the
material is in practice slightly different from the color of the
base layer on its own---in particular, the material color will tend to shift towards
darker colors with higher saturation. Since this can be counter-intuitive when
using bitmap textures, these color shifts are disabled by default. Specify
the parameter \code{nonlinear=true} to enable them. \figref{plastic-nonlinear}
illustrates the resulting change. This effect is also seen in real life,
for instance a piece of wood will look slightly darker after coating it
with a layer of varnish.
\plugin{roughplastic}{Rough plastic material}
\order{9}
\icon{bsdf_roughplastic}
\parameters{
    \parameter{distribution}{\String}{
         Specifies the type of microfacet normal distribution
         used to model the surface roughness.
         \vspace{-1mm}
      \begin{enumerate}[(i)]
          \item \code{beckmann}: Physically-based distribution derived from
              Gaussian random surfaces. This is the default.\vspace{-1.5mm}
          \item \code{ggx}: The GGX \cite{Walter07Microfacet} distribution (also known as
              Trowbridge-Reitz \cite{Trowbridge19975Average} distribution)
              was designed to better approximate the long tails observed in measurements
              of ground surfaces, which are not modeled by the Beckmann distribution.
          \vspace{-1.5mm}
          \item \code{phong}: Classical Phong distribution.
             In most cases, the \code{ggx} and \code{beckmann} distributions
             should be preferred, since they provide better importance sampling
             and accurate shadowing/masking computations.
             \vspace{-4mm}
      \end{enumerate}
    }
    \parameter{alpha}{\Float\Or\Texture}{
        Specifies the roughness of the unresolved surface micro-geometry.
        When the Beckmann distribution is used, this parameter is equal to the
        \emph{root mean square} (RMS) slope of the microfacets.
        \default{0.1}.
    }

    \parameter{intIOR}{\Float\Or\String}{Interior index of refraction specified
     numerically or using a known material name. \default{\texttt{polypropylene} / 1.49}}
    \parameter{extIOR}{\Float\Or\String}{Exterior index of refraction specified
     numerically or using a known material name. \default{\texttt{air} / 1.000277}}
    \parameter{sampleVisible}{\Boolean}{
        Enables an improved importance sampling technique. Refer to
        pages \pageref{plg:roughconductor} and \pageref{sec:visiblenormal-sampling}
        for details. \default{\code{true}}
    }
    \parameter{specular\showbreak Reflectance}{\Spectrum\Or\Texture}{Optional
        factor that can be used to modulate the specular reflection component. Note
        that for physical realism, this parameter should never be touched. \default{1.0}}
    \parameter{diffuse\showbreak Reflectance}{\Spectrum\Or\Texture}{Optional
        factor used to modulate the diffuse reflection component\default{0.5}}
    \parameter{nonlinear}{\Boolean}{
        Account for nonlinear color shifts due to internal scattering? See the
        \pluginref{plastic} plugin for details.\default{Don't account for them and
        preserve the texture colors, i.e. \code{false}}
    }
}

\vspace{3mm}
This plugin implements a realistic microfacet scattering model for rendering
rough dielectric materials with internal scattering, such as plastic. It can
be interpreted as a fancy version of the Cook-Torrance model and should be
preferred over heuristic models like \pluginref{phong} and \pluginref{ward}
when possible.

Microfacet theory describes rough surfaces as an arrangement of
unresolved and ideally specular facets, whose normal directions are given by
a specially chosen \emph{microfacet distribution}. By accounting for shadowing
and masking effects between these facets, it is possible to reproduce the important
off-specular reflections peaks observed in real-world measurements of such
materials.

\renderings{
    \rendering{Beckmann, $\alpha=0.1$}{bsdf_roughplastic_beckmann}
    \rendering{GGX, $\alpha=0.3$}{bsdf_roughplastic_ggx}
}

This plugin is essentially the ``roughened'' equivalent of the (smooth) plugin
\pluginref{plastic}. For very low values of $\alpha$, the two will
be identical, though scenes using this plugin will take longer to render
due to the additional computational burden of tracking surface roughness.

For convenience, this model allows to specify IOR values either numerically,
or based on a list of known materials (see \tblref{dielectric-iors} on
\tblpage{dielectric-iors} for an overview).
When no parameters are given, the plugin activates the defaults,
which describe a white polypropylene plastic material with a light amount
of roughness modeled using the Beckmann distribution.

Like the \pluginref{plastic} material, this model internally simulates the
interaction of light with a diffuse base surface coated by a thin dielectric
layer (where the coating layer is now \emph{rough}). This is a convenient
abstraction rather than a restriction. In other words, there are many
materials that can be rendered with this model, even if they might not
fit this description perfectly well.

The simplicity of this setup makes it possible to account for interesting
nonlinear effects due to internal scattering, which is controlled by
the \texttt{nonlinear} parameter. For more details, please refer to the description
of this parameter given in the \pluginref{plastic} plugin section
on \pluginpage{plastic}.

To get an intuition about the effect of the surface roughness parameter
$\alpha$, consider the following approximate classification: a value of
$\alpha=0.001-0.01$ corresponds to a material with slight imperfections
on an otherwise smooth surface finish, $\alpha=0.1$ is relatively rough,
and $\alpha=0.3-0.7$ is \emph{extremely} rough (e.g. an etched or ground
finish). Values significantly above that are probably not too realistic.

\renderings{
    \medrendering{Diffuse textured rendering}{bsdf_plastic_diffuse}
    \medrendering{Textured rough plastic model and \code{nonlinear=false}}{bsdf_roughplastic_preserve}
    \medrendering{Textured rough plastic model and \code{nonlinear=true}}{bsdf_roughplastic_nopreserve}
    \caption{
       When asked to do so, this model can account for subtle nonlinear color shifts due
       to internal scattering processes. The above images show a textured
       object first rendered using \pluginref{diffuse}, then
       \pluginref{roughplastic} with the default parameters, and finally using
       \pluginref{roughplastic} and support for nonlinear color shifts.
    }
}
\renderings{
    \rendering{Wood material with smooth horizontal stripes}{bsdf_roughplastic_roughtex1}
    \rendering{A material with imperfections at a much smaller scale than what
      is modeled e.g. using a bump map.}{bsdf_roughplastic_roughtex2}\vspace{-3mm}
    \caption{
        The ability to texture the roughness parameter makes it possible
        to render materials with a structured finish, as well as
        ``smudgy'' objects.
    }
}
\vspace{2mm}
\begin{xml}[caption={A material definition for black plastic material with
   a spatially varying roughness.},
   label=lst:roughplastic-varyingalpha]
<bsdf type="roughplastic">
    <string name="distribution" value="beckmann"/>
    <float name="intIOR" value="1.61"/>
    <spectrum name="diffuseReflectance" value="0"/>
    <!-- Fetch roughness values from a texture and slightly reduce them -->
    <texture type="scale" name="alpha">
        <texture name="alpha" type="bitmap">
            <string name="filename" value="roughness.png"/>
        </texture>
        <float name="scale" value="0.6"/>
    </texture>
</bsdf>
\end{xml}

\subsubsection*{Technical details}
The implementation of this model is partly based on the paper ``Microfacet
Models for Refraction through Rough Surfaces'' by Walter et al.
\cite{Walter07Microfacet}. Several different types of microfacet
distributions are supported. Note that the choices are slightly more
restricted here---in comparison to other rough scattering models in
Mitsuba, anisotropic distributions are not allowed.

The implementation of this model makes heavy use of a \emph{rough
Fresnel transmittance} function, which is a generalization of the
usual Fresnel transmittion coefficient to microfacet surfaces. Unfortunately,
this function is normally prohibitively expensive, since each
evaluation involves a numerical integration over the sphere.

To avoid this performance issue, Mitsuba ships with data files
(contained in the \code{data/microfacet} directory) containing precomputed
values of this function over a large range of parameter values. At runtime,
the relevant parts are extracted using tricubic interpolation.

When rendering with the Phong microfacet distribution, a conversion is
used to turn the specified Beckmann-equivalent $\alpha$ roughness value
into the exponent parameter of this distribution. This is done in a way,
such that the same value $\alpha$ will produce a similar appearance across
different microfacet distributions.
\plugin{coating}{Smooth dielectric coating}
\order{10}
\icon{bsdf_coating}

\parameters{
    \parameter{intIOR}{\Float\Or\String}{Interior index of refraction specified
     numerically or using a known material name. \default{\texttt{bk7} / 1.5046}}
    \parameter{extIOR}{\Float\Or\String}{Exterior index of refraction specified
     numerically or using a known material name. \default{\texttt{air} / 1.000277}}
    \parameter{thickness}{\Float}{Denotes the thickness of the layer (to
     model absorption --- should be specified in inverse units of \code{sigmaA})\default{1}}
    \parameter{sigmaA}{\Spectrum\Or\Texture}{The absorption coefficient of the
     coating layer. \default{0, i.e. there is no absorption}}
    \parameter{specular\showbreak Reflectance}{\Spectrum\Or\Texture}{Optional
        factor that can be used to modulate the specular reflection component. Note
        that for physical realism, this parameter should never be touched. \default{1.0}}
    \parameter{\Unnamed}{\BSDF}{A nested BSDF model that should be coated.}
}

\renderings{
    \rendering{Rough copper}
        {bsdf_coating_uncoated}
    \rendering{The same material coated with a single layer of
        clear varnish (see \lstref{coating-roughcopper})}
        {bsdf_coating_roughconductor}
}

This plugin implements a smooth dielectric coating (e.g. a layer of varnish)
in the style of the paper ``Arbitrarily Layered Micro-Facet Surfaces'' by
Weidlich and Wilkie \cite{Weidlich2007Arbitrarily}. Any BSDF in Mitsuba
can be coated using this plugin, and multiple coating layers can even
be applied in sequence. This allows designing interesting custom materials
like car paint or glazed metal foil. The coating layer can optionally be
tinted (i.e. filled with an absorbing medium), in which case this model also
accounts for the directionally dependent absorption within the layer.

Note that the plugin discards illumination that undergoes internal
reflection within the coating. This can lead to a noticeable energy
loss for materials that reflect much of their energy near or below the critical
angle (i.e. diffuse or very rough materials).
Therefore, users are discouraged to use this plugin to coat smooth
diffuse materials, since there is a separately available plugin
named \pluginref{plastic}, which covers the same case and does not
suffer from energy loss.\newpage

\renderings{
    \smallrendering{$\code{thickness}=0$}{bsdf_coating_0}
    \smallrendering{$\code{thickness}=1$}{bsdf_coating_1}
    \smallrendering{$\code{thickness}=5$}{bsdf_coating_5}
    \smallrendering{$\code{thickness}=15$}{bsdf_coating_15}
    \caption{The effect of the layer thickness parameter on
       a tinted coating ($\code{sigmaT}=(0.1, 0.2, 0.5)$)}
}

\vspace{4mm}

\begin{xml}[caption=Rough copper coated with a transparent layer of
    varnish, label=lst:coating-roughcopper]
<bsdf type="coating">
    <float name="intIOR" value="1.7"/>

    <bsdf type="roughconductor">
        <string name="material" value="Cu"/>
        <float name="alpha" value="0.1"/>
    </bsdf>
</bsdf>
\end{xml}

\renderings{
    \rendering{Coated rough copper with a bump map applied on top}{bsdf_coating_coatedbump}
    \rendering{Bump mapped rough copper with a coating on top}{bsdf_coating_bumpcoating}
    \caption{Some interesting materials can be created simply by applying
    Mitsuba's material modifiers in different orders.}
}

\subsubsection*{Technical details}
Evaluating the internal component of this model entails refracting the
incident and exitant rays through the dielectric interface, followed by
querying the nested material with this modified direction pair. The result
is attenuated by the two Fresnel transmittances and the absorption, if
any.
\plugin{roughcoating}{Rough dielectric coating}
\order{11}
\icon{bsdf_roughcoating}
\parameters{
    \parameter{distribution}{\String}{
         Specifies the type of microfacet normal distribution
         used to model the surface roughness.
         \vspace{-1mm}
      \begin{enumerate}[(i)]
          \item \code{beckmann}: Physically-based distribution derived from
              Gaussian random surfaces. This is the default.\vspace{-1.5mm}
          \item \code{ggx}: The GGX \cite{Walter07Microfacet} distribution (also known as
              Trowbridge-Reitz \cite{Trowbridge19975Average} distribution)
              was designed to better approximate the long tails observed in measurements
              of ground surfaces, which are not modeled by the Beckmann distribution.
          \vspace{-1.5mm}
          \item \code{phong}: Classical Phong distribution.
             In most cases, the \code{ggx} and \code{beckmann} distributions
             should be preferred, since they provide better importance sampling
             and accurate shadowing/masking computations.
             \vspace{-4mm}
      \end{enumerate}
    }
    \parameter{alpha}{\Float\Or\Texture}{
        Specifies the roughness of the unresolved surface micro-geometry.
        When the Beckmann distribution is used, this parameter is equal to the
        \emph{root mean square} (RMS) slope of the microfacets.
        \default{0.1}.
    }
    \parameter{sampleVisible}{\Boolean}{
        Enables an improved importance sampling technique. Refer to
        pages \pageref{plg:roughconductor} and \pageref{sec:visiblenormal-sampling}
        for details. \default{\code{true}}
    }
    \parameter{intIOR}{\Float\Or\String}{Interior index of refraction specified
     numerically or using a known material name. \default{\texttt{bk7} / 1.5046}}
    \parameter{extIOR}{\Float\Or\String}{Exterior index of refraction specified
     numerically or using a known material name. \default{\texttt{air} / 1.000277}}
    \parameter{thickness}{\Float}{Denotes the thickness of the layer (to
     model absorption --- should be specified in inverse units of \code{sigmaA})\default{1}}
    \parameter{sigmaA}{\Spectrum\Or\Texture}{The absorption coefficient of the
     coating layer. \default{0, i.e. there is no absorption}}
    \parameter{specular\showbreak Reflectance}{\Spectrum\Or\Texture}{Optional
        factor that can be used to modulate the specular reflection component. Note
        that for physical realism, this parameter should never be touched. \default{1.0}}
    \parameter{\Unnamed}{\BSDF}{A nested BSDF model that should be coated.}
}
\renderings{
    \rendering{Rough gold coated with a \emph{smooth} varnish layer}
        {bsdf_roughcoating_gold_smooth}
    \rendering{Rough gold coated with a \emph{rough} ($\alpha\!=\!0.03$) varnish layer}
        {bsdf_roughcoating_gold_rough}
}

This plugin implements a \emph{very} approximate\footnote{
The model only accounts for roughness
in the specular reflection and Fresnel transmittance through the interface.
The interior model receives incident illumination
that is transformed \emph{as if} the coating was smooth. While
that's not quite correct, it is a convenient workaround when the
\pluginref{coating} plugin produces specular highlights that are too sharp.}
model that simulates a rough dielectric coating. It is essentially the
roughened version of \pluginref{coating}.
Any BSDF in Mitsuba can be coated using this plugin and multiple coating
layers can even be applied in sequence, which allows designing interesting
custom materials. The coating layer can optionally be tinted (i.e. filled
with an absorbing medium), in which case this model also accounts for the
directionally dependent absorption within the layer.

Note that the plugin discards illumination that undergoes internal
reflection within the coating. This can lead to a noticeable energy
loss for materials that reflect much of their energy near or below the critical
angle (i.e. diffuse or very rough materials).

The implementation here is motivated by the paper
``Arbitrarily Layered Micro-Facet Surfaces'' by Weidlich and
Wilkie \cite{Weidlich2007Arbitrarily}, though the implementation
works differently.
\plugin{bumpmap}{Bump map modifier}
\order{12}
\icon{bsdf_bumpmap}

\parameters{
    \parameter{\Unnamed}{\Texture}{
      The luminance of this texture specifies the amount of
      displacement. The implementation ignores any constant
      offset---only changes in the luminance matter.
    }
    \parameter{\Unnamed}{\BSDF}{A BSDF model that should
    be affected by the bump map}
}
\renderings{
    \rendering{Bump map based on tileable diagonal lines}{bsdf_bumpmap_1}
    \rendering{An irregular bump map}{bsdf_bumpmap_2}
}

Bump mapping \cite{Blinn1978Simulation} is a simple technique for cheaply
adding surface detail to a rendering. This is done by perturbing the
shading coordinate frame based on a displacement height field provided
as a texture. This method can lend objects a highly realistic and detailed
appearance (e.g. wrinkled or covered by scratches and other imperfections)
without requiring any changes to the input geometry.

The implementation in Mitsuba uses the common approach of ignoring
the usually negligible texture-space derivative of the base mesh
surface normal. As side effect of this decision, it is invariant
to constant offsets in the height field texture---only variations in
its luminance cause changes to the shading frame.

Note that the magnitude of the height field variations influences
the strength of the displacement. If desired, the \pluginref{scale}
texture plugin can be used to magnify or reduce the effect of a
bump map texture.
\begin{xml}[caption=A rough metal model with a scaled image-based bump map]
<bsdf type="bumpmap">
    <!-- The bump map is applied to a rough metal BRDF -->
    <bsdf type="roughconductor"/>

    <texture type="scale">
        <!-- The scale of the displacement gets multiplied by 10x -->
        <float name="scale" value="10"/>

        <texture type="bitmap">
            <string name="filename" value="bumpmap.png"/>
        </texture>
    </texture>
</bsdf>
\end{xml}
\plugin{normalmap}{Normal map modifier}
\order{13}
\icon{bsdf_bumpmap}

\parameters{
    \parameter{\Unnamed}{\Texture}{
      The color values of this texture specify the perturbed
      normals relative in the local surface coordinate system.
    }
    \parameter{\Unnamed}{\BSDF}{A BSDF model that should
    be affected by the normal map}
}

This plugin is conceptually similar to the \pluginref{bumpmap} plugin
but uses a normal map instead of a bump map. A normal map is a RGB texture, whose color channels
encode the XYZ coordinates of the desired surface normals.
These are specified \emph{relative} to the local shading frame,
which means that a normal map with a value of $(0,0,1)$ everywhere
causes no changes to the surface.
To turn the 3D normal directions into (nonnegative) color values
suitable for this plugin, the
mapping $x \mapsto (x+1)/2$ must be applied to each component.
\plugin{phong}{Modified Phong BRDF}
\order{14}
\parameters{
    \parameter{exponent}{\Float\Or\Texture}{
        Specifies the Phong exponent \default{30}.
    }
    \parameter{specular\showbreak Reflectance}{\Spectrum\Or\Texture}{
        Specifies the weight of the specular reflectance component.\default{0.2}}
    \parameter{diffuse\showbreak Reflectance}{\Spectrum\Or\Texture}{
        Specifies the weight of the diffuse reflectance component\default{0.5}}
}
\renderings{
    \rendering{Exponent$\,=60$}{bsdf_phong_60}
    \rendering{Exponent$\,=300$}{bsdf_phong_300}
}

This plugin implements the modified Phong reflectance model as described in
\cite{Phong1975Illumination} and \cite{Lafortune1994Using}. This heuristic
model is mainly included for historical reasons---its use in new scenes is
discouraged, since significantly more realistic models have been developed
since 1975.

If possible, it is recommended to switch to a BRDF that is based on
microfacet theory and includes knowledge about the material's index of
refraction. In Mitsuba, two good alternatives to \pluginref{phong} are
the plugins \pluginref{roughconductor} and \pluginref{roughplastic}
(depending on the material type).

When using this plugin, note that the diffuse and specular reflectance
components should add up to a value less than or equal to one (for each
color channel). Otherwise, they will automatically be scaled appropriately
to ensure energy conservation.
\plugin{ward}{Anisotropic Ward BRDF}
\order{15}
\parameters{
    \parameter{variant}{\String}{
        Determines the variant of the Ward model to use:
        \begin{enumerate}[(i)]
            \item \code{ward}: The original model by Ward \cite{Ward1992Measuring}
            --- suffers from energy loss at grazing angles.
            \item \code{ward-duer}: Corrected Ward model with lower energy loss
            at grazing angles \cite{Dur2006Improved}.
            Does not always conserve energy.
            \item \code{balanced}: Improved version of the \code{ward-duer}
            model with energy balance at all angles \cite{Geisler2010New}.
        \end{enumerate}
        Default: \texttt{balanced}
    }
    \parameter{alphaU, alphaV}{\Float\Or\Texture}{
        Specifies the anisotropic roughness values along the tangent and
        bitangent directions.
        \default{0.1}.
    }
    \parameter{specular\showbreak Reflectance}{\Spectrum\Or\Texture}{
        Specifies the weight of the specular reflectance component.\default{0.2}}
    \parameter{diffuse\showbreak Reflectance}{\Spectrum\Or\Texture}{
        Specifies the weight of the diffuse reflectance component\default{0.5}}
}
\renderings{
    \rendering{$\alpha_u=0.1,\ \alpha_v=0.3$}{bsdf_ward_01_03}
    \rendering{$\alpha_u=0.3,\ \alpha_v=0.1$}{bsdf_ward_03_01}
}

This plugin implements the anisotropic Ward reflectance model and
several extensions. They are described in the papers
\begin{enumerate}[(i)]
   \item ``Measuring and Modeling Anisotropic Reflection''
     by Greg Ward \cite{Ward1992Measuring}
   \item ``Notes on the Ward BRDF'' by Bruce Walter \cite{Walter2005Notes}
   \item ``An Improved Normalization for the Ward Reflectance Model''
     by Arne D\"ur \cite{Dur2006Improved}
   \item ``A New Ward BRDF Model with Bounded Albedo'' by
     Geisler-Moroder et al. \cite{Geisler2010New}
\end{enumerate}

Like the Phong BRDF, the Ward model does not take the Fresnel reflectance
of the material into account. In an experimental study by Ngan et al.
\cite{Ngan2005Experimental}, the Ward model performed noticeably worse than
models based on microfacet theory.

For this reason, it is usually preferable to switch to a microfacet model
that incorporates knowledge about the material's index of refraction. In Mitsuba,
two such alternatives to \pluginref{ward} are given by the plugins
\pluginref{roughconductor} and \pluginref{roughplastic} (depending on the
material type).

When using this plugin, note that the diffuse and specular reflectance
components should add up to a value less than or equal to one (for each
color channel). Otherwise, they will automatically be scaled appropriately
to ensure energy conservation.
\plugin{mixturebsdf}{Mixture material}
\order{16}
\parameters{
    \parameter{weights}{\String}{A comma-separated list of BSDF weights}
    \parameter{\Unnamed}{\BSDF}{Multiple BSDF instances that should be
    mixed according to the specified weights}
}
\renderings{
    \medrendering{Smooth glass}{bsdf_mixturebsdf_smooth}
    \medrendering{Rough glass}{bsdf_mixturebsdf_rough}
    \medrendering{An mixture of 70\% smooth glass and 30\% rough glass
    results in a more realistic smooth material with imperfections
    (\lstref{mixture-example})}{bsdf_mixturebsdf_result}
}

This plugin implements a ``mixture'' material, which represents
linear combinations of multiple BSDF instances. Any surface scattering
model in Mitsuba (be it smooth, rough, reflecting, or transmitting) can
be mixed with others in this manner to synthesize new models. There
is no limit on how many models can be mixed, but their combination
weights must be non-negative and sum to a value of one or less to ensure
energy balance. When they sum to less than one, the material will
absorb a proportional amount of the incident illlumination.

\vspace{4mm}
\begin{xml}[caption={A material definition for a mixture of 70\% smooth
    and 30\% rough glass},
    label=lst:mixture-example]
<bsdf type="mixturebsdf">
    <string name="weights" value="0.7, 0.3"/>

    <bsdf type="dielectric"/>

    <bsdf type="roughdielectric">
        <float name="alpha" value="0.3"/>
    </bsdf>
</bsdf>
\end{xml}
\plugin{blendbsdf}{Blended material}
\order{17}
\parameters{
    \parameter{weight}{\Float\Or\Texture}{A floating point value or texture
     with values between zero and one. The extreme values zero and one activate the
     first and second nested BSDF respectively, and inbetween values
     interpolate accordingly. \default{0.5}}
    \parameter{\Unnamed}{\BSDF}{Two nested BSDF instances that should be
    mixed according to the specified blending weight}
}

\renderings{
    \rendering{A material created by blending between dark rough plastic and
    smooth gold based on a binary bitmap texture (\lstref{blendbsdf})}{bsdf_blendbsdf}
}

This plugin implements a ``blend'' material, which represents
linear combinations of two BSDF instances. It is conceptually very similar
to the \pluginref{mixturebsdf} plugin. The main difference is that
\pluginref{blendbsdf} can interpolate based on a texture rather than a set
of constants.

Any surface scattering model in Mitsuba (be it smooth, rough, reflecting, or
transmitting) can be mixed with others in this manner to synthesize new models.

\begin{xml}[caption=Description of the material shown above,
    label=lst:blendbsdf]
<bsdf type="blendbsdf">
    <texture name="weight" type="bitmap">
        <string name="wrapMode" value="repeat"/>
        <string name="filename" value="pattern.png"/>
    </texture>

    <bsdf type="conductor">
        <string name="material" value="Au"/>
    </bsdf>

    <bsdf type="roughplastic">
        <spectrum name="diffuseReflectance" value="0"/>
    </bsdf>
</bsdf>
\end{xml}
\plugin{mask}{Opacity mask}
\order{19}
\parameters{
    \parameter{opacity}{\Spectrum\Or\Texture}{
         Specifies the per-channel opacity (where $1=$ completely opaque)\default{0.5}.
    }
    \parameter{\Unnamed}{\BSDF}{A base BSDF model that represents the
        non-transparent portion of the scattering}
}
\renderings{
    \rendering{Rendering without an opacity mask}
        {bsdf_mask_before.jpg}
    \rendering{Rendering \emph{with} an opacity mask (\lstref{mask-leaf})}
        {bsdf_mask_after.jpg}
}
This plugin applies an opacity mask to add nested BSDF instance. It interpolates
between perfectly transparent and completely opaque based on the \code{opacity}
parameter.

The transparency is internally implemented as a forward-facing Dirac delta distribution.
Note that the standard \pluginref{path} tracer does not have a good
sampling strategy to deal with this, but the volumetric path tracer
(\pluginref{volpath}) does. It may thus be preferable when rendering scenes
that contain the \pluginref{mask} plugin, even if there is nothing
``volumetric'' in the scene.
\vspace{5mm}

\begin{xml}[caption=Material configuration for a transparent leaf,
   label=lst:mask-leaf]
<bsdf type="mask">
    <!-- Base material: a two-sided textured diffuse BSDF -->
    <bsdf type="twosided">
        <bsdf type="diffuse">
            <texture name="reflectance" type="bitmap">
                <string name="filename" value="leaf.png"/>
            </texture>
        </bsdf>
    </bsdf>

    <!-- Fetch the opacity mask from the alpha channel -->
    <texture name="opacity" type="bitmap">
        <string name="filename" value="leaf.png"/>
        <string name="channel" value="a"/>
    </texture>
</bsdf>
\end{xml}
\plugin{twosided}{Two-sided BRDF adapter}
\order{19}
\parameters{
    \parameter{\Unnamed}{\BSDF}{A nested BRDF that should
    be turned into a two-sided scattering model. If two BRDFs
    are specified, they will be placed on the front and back side, respectively.}
}

\renderings{
    \unframedrendering{From this angle, the Cornell box scene shows visible back-facing geometry}
        {bsdf_twosided_before}
    \unframedrendering{Applying the \pluginref{twosided} plugin fixes the rendering}
        {bsdf_twosided_after}
}

By default, all non-transmissive scattering models in Mitsuba
are \emph{one-sided} --- in other words, they absorb all light
that is received on the interior-facing side of any associated
surfaces. Holes and visible back-facing parts are thus exposed
as black regions.

Usually, this is a good idea, since it will reveal modeling
issues early on. But sometimes one is forced to deal with
improperly closed geometry, where the one-sided behavior is
bothersome. In that case, this plugin can be used to turn
one-sided scattering models into proper two-sided versions of
themselves. The plugin has no parameters other than a required
nested BSDF specification. It is also possible to supply two
different BRDFs that should be placed on the front and back
side, respectively.
\vspace{4mm}

\begin{xml}[caption=A two-sided diffuse material]
<bsdf type="twosided">
    <bsdf type="diffuse">
         <spectrum name="reflectance" value="0.4"/>
    </bsdf>
</bsdf>
\end{xml}
\plugin{difftrans}{Diffuse transmitter}
\icon{bsdf_difftrans}

\parameters{
    \parameter{transmittance}{\Spectrum\Or\Texture}{
      Specifies the diffuse transmittance of the material
      \default{0.5}
    }
}

\renderings{
    \rendering{The model with default parameters}{bsdf_difftrans.jpg}
}

This BSDF models a non-reflective material, where any entering light loses
its directionality and is diffusely scattered from the other side. This
model can be combined\footnote{For instance using the
\pluginref{mixturebsdf} plugin.} with a surface reflection model to
describe translucent substances that have internal multiple scattering
processes (e.g. plant leaves).
\plugin{hk}{Hanrahan-Krueger BSDF}
\icon{bsdf_hk}

\parameters{
    \parameter{material}{\String}{Name of a material preset, see
          \tblref{medium-coefficients}. \default{\texttt{skin1}}}
    \parameter{sigmaS}{\Spectrum\Or\Texture}{Specifies the scattering coefficient
     of the internal layer. \default{based on \code{material}}}
    \parameter{sigmaA}{\Spectrum\Or\Texture}{Specifies the absorption coefficient
     of the internal layer. \default{based on \code{material}}}
    \parameter{sigmaT \& albedo}{\Spectrum\Or\Texture}{
     Optional: Alternatively, the scattering and absorption coefficients may also be
     specified using the extinction coefficient \code{sigmaT} and the
     single-scattering albedo. Note that only one of the parameter passing
     conventions can be used at a time (i.e. use either \code{sigmaS\&sigmaA}
     \emph{or} \code{sigmaT\&albedo})}
    \parameter{thickness}{\Float}{Denotes the thickness of the layer.
     (should be specified in inverse units of \code{sigmaA} and \code{sigmaS})\default{1}}
    \parameter{\Unnamed}{\Phase}{A nested phase function instance that represents
     the type of scattering interactions occurring within the layer}
}

\renderings{
    \rendering{An index-matched scattering layer with parameters $\sigma_s=2$, $\sigma_a=0.1$, thickness$=0.1$}{bsdf_hk_1}
    \rendering{Example of the HK model with a dielectric coating (and the \code{ketchup} material preset, see \lstref{hk-coated})}{bsdf_hk_2}
    \caption{
    \label{fig:hk-example}
    Renderings using the uncoated and coated form of the Hanrahan-Krueger model.
    }
    \vspace{3mm}
}

This plugin provides an implementation of the Hanrahan-Krueger BSDF
\cite{Hanrahan1993Reflection} for simulating single scattering in thin
index-matched layers filled with a random scattering medium.
In addition, the implementation also accounts for attenuated
light that passes through the medium without undergoing any scattering events.

This BSDF requires a phase function to model scattering interactions within the
random medium. When no phase function is explicitly specified, it uses an
isotropic one ($g=0$) by default. A sample usage for instantiating the
plugin is given on the next page:\newpage
\begin{xml}
<bsdf type="hk">
    <spectrum name="sigmaS" value="2"/>
    <spectrum name="sigmaA" value="0.1"/>
    <float name="thickness" value="0.1"/>

    <phase type="hg">
        <float name="g" value="0.8"/>
    </phase>
</bsdf>
\end{xml}

When used in conjuction with the \pluginref{coating} plugin, it is possible
to model refraction and reflection at the layer boundaries when the indices
of refraction are mismatched. The combination of these two plugins then
reproduces the full model as it was originally proposed by Hanrahan and
Krueger \cite{Hanrahan1993Reflection}.

Note that this model does not account for light that undergoes multiple
scattering events within the layer. This leads to energy loss,
particularly at grazing angles, which can be seen in the left-hand image of
\figref{hk-example}.

\begin{xml}[caption=A thin dielectric layer with measured ketchup scattering parameters, label=lst:hk-coated]
<bsdf type="coating">
    <float name="extIOR" value="1.0"/>
    <float name="intIOR" value="1.5"/>

    <bsdf type="hk">
        <string name="material" value="ketchup"/>
        <float name="thickness" value="0.01"/>
    </bsdf>
</bsdf>
\end{xml}

Note that when \texttt{sigmaS} = \texttt{sigmaA}$\ = 0$, or when \texttt{thickness=0},
any geometry associated with this BSDF becomes invisible, as light will pass through
unchanged.

The implementation in Mitsuba is based on code by Tom Kazimiers and Marios Papas.
Marios Papas has kindly verified the implementation of the coated and uncoated variants
against both a path tracer and a separate reference implementation.
\plugin{irawan}{Irawan \& Marschner woven cloth BRDF}

\parameters{
    \parameter{filename}{\String}{Path to a weave pattern description}
    \parameter{repeatU, repeatV}{\Float}{Specifies the number
        of weave pattern repetitions over a $[0,1]^2$ region of the UV
        parameterization}
    \parameter{(\emph{Additional parameters})}{\Spectrum\Or\Float}{
        Weave pattern files may define their own custom parameters; this is
        useful for instance to support changing the color of a weave
        without having to create a new file every time.
        These parameters must be specified directly to the plugin
        so that they can be appropriately resolved when the pattern file is loaded.
    }
}

This plugin implements the Irawan \& Marschner BRDF,
a realistic model for rendering woven materials.
This spatially-varying reflectance model uses an explicit description
of the underlying weave pattern to create fine-scale texture and
realistic reflections across a wide range of different weave types.
To use the model, you must provide a special weave pattern
file---for an example of what these look like, see the
examples scenes available on the Mitsuba website.

A detailed explanation of the model is beyond the scope of this manual.
For reference, it is described in detail in the PhD thesis of
Piti Irawan (``The Appearance of Woven Cloth'' \cite{IrawanThesis}).
The code in Mitsuba a modified port of a previous Java implementation
by Piti, which has been extended with a simple domain-specific weave
pattern description language. \vspace{8mm}

\renderings{
    \unframedmedrendering{Silk charmeuse}{bsdf_irawan_charmeuse}
    \unframedmedrendering{Cotton denim}{bsdf_irawan_denim}
    \unframedmedrendering{Wool gabardine}{bsdf_irawan_gabardine}
}
\renderings{
    \setcounter{subfigure}{3}
    \unframedmedrendering{Polyester lining cloth}{bsdf_irawan_polyester}
    \unframedmedrendering{Silk shantung}{bsdf_irawan_shantung}
    \unframedmedrendering{Cotton twill}{bsdf_irawan_twill}
}
\input{section_textures}
\plugin{bitmap}{Bitmap texture}
\order{1}
\parameters{
    \parameter{filename}{\String}{
      Filename of the bitmap to be loaded
    }
    \parameter{wrapMode, wrapModeU, wrapModeV}{\String}{
      Behavior of texture lookups outside of the $[0,1]$ $uv$ range.\vspace{-1mm}
      \begin{enumerate}[(i)]
          \item \code{repeat}: Repeat the texture indefinitely\vspace{-1mm}
          \item \code{mirror}: Mirror the texture along its boundaries\vspace{-1mm}
          \item \code{clamp}: Clamp $uv$ coordinates to $[0,1]$ before a lookup\vspace{-1mm}
          \item \code{zero}: Switch to a zero-valued texture \vspace{-1mm}
          \item \code{one}: Switch to a one-valued texture \vspace{-1mm}
      \end{enumerate}
      Default: \code{repeat}. The parameter \code{wrapMode} is a shortcut for
      setting both \code{wrapModeU} and \code{wrapModeV} at the same time.
    }
    \parameter{gamma}{\Float}{
      Optional parameter to override the gamma value of the source bitmap,
      where 1 indicates a linear color space and the special value -1
      corresponds to sRGB. \default{automatically detect based on the
      image type and metadata}
    }
    \parameter{filterType}{\String}{
      Specifies the texture filturing that should be used for lookups\vspace{-1mm}
      \begin{enumerate}[(i)]
          \item \code{ewa}: Elliptically weighted average (a.k.a.
          anisotropic filtering). This produces the best quality\vspace{-1mm}
          \item \code{trilinear}: Simple trilinear (isotropic) filtering.\vspace{-1mm}
          \item \code{nearest}: No filtering, do nearest neighbor lookups.\vspace{-1mm}
      \end{enumerate}
      Default: \code{ewa}.
    }
    \parameter{maxAnisotropy}{\Float}{
       Specific to \code{ewa} filtering, this parameter limits the
       anisotropy (and thus the computational cost) of filtured texture lookups. The
       default of 20 is a good compromise.
    }
    \parameter{cache}{\Boolean}{
       Preserve generated MIP map data in a cache file? This will cause a file named
       \emph{filename}\code{.mip} to be created.
       \default{automatic---use caching for textures larger than 1M pixels.}
    }
    \parameter{uoffset, voffset}{\Float}{
      Numerical offset that should be applied to UV lookups
    }
    \parameter{uscale, vscale}{\Float}{
      Multiplicative factors that should be applied to UV lookups
    }
    \parameter{channel}{\String}{
      Create a monochromatic texture based on one of the image channels
      (e.g. \texttt{r}, \texttt{g}, \texttt{b}, \texttt{a}, \texttt{x},
      \texttt{y}, \texttt{z} etc.). \default{use all channels}
    }
}
This plugin provides a bitmap-backed texture source that supports \emph{filtered}
texture lookups on\footnote{Some of these may not be available depending on how
Mitsuba was compiled.} JPEG, PNG, OpenEXR, RGBE, TGA, and BMP files. Filtered
lookups are useful to avoid aliasing when rendering textures that contain high
frequencies (see the next page for an example).

The plugin operates as follows: when loading a bitmap file, it is first converted
into a linear color space. Following this, a MIP map is constructed that is necessary
to perform filtered lookups during rendering. A \emph{MIP map} is a hierarchy of
progressively lower resolution versions of the input image, where the resolution of
adjacent levels differs by a factor of two. Mitsuba creates this hierarchy using
Lanczos resampling to obtain very high quality results.
Note that textures may have an arbitrary resolution and are not limited to powers of two.
Three different filtering modes are supported:

\begin{enumerate}[(i)]
\item Nearest neighbor lookups effectively disable filtering and always query the
highest-resolution version of the texture without any kind of interpolation. This is
fast and requires little memory (no MIP map is created), but results in visible aliasing.
Only a single pixel value is accessed.

\item The trilinear filter performs bilinear interpolation on two adjacent MIP levels
and blends the results. Because it cannot do anisotropic (i.e. slanted) lookups in texture space,
it must compromise either on the side of blurring or aliasing. The implementation in Mitsuba
chooses blurring over aliasing (though note that (\textbf{b}) is an extreme case).
Only 8 pixel values are accessed.

\item The EWA filter performs anisotropicically filtered lookups on two adjacent MIP map levels
and blends them. This produces the best quality, but at the expense of computation time.
Generally, 20-40 pixel values must be read for a single EWA texture lookup. To limit
the number of pixel accesses, the \code{maxAnisotropy} parameter can be used to bound
the amount of anisotropy that a texture lookup is allowed to have.
\end{enumerate}
\renderings{
    \rendering{Nearest-neighbor filter. Note the aliasing}{tex_bitmap_nearest}
    \rendering{Trilinear filter. Note the blurring}{tex_bitmap_trilinear}
    \vspace{-5mm}
}
\renderings{
    \setcounter{subfigure}{2}
    \rendering{EWA filter}{tex_bitmap_ewa}
    \rendering{Ground truth (512 samples per pixel)}{tex_bitmap_gt}
    \caption{A somewhat contrived comparison of the different filters when rendering a high-frequency
    checkerboard pattern using four samples per pixel. The EWA method (the default)
    pre-filters the texture anisotropically to limit blurring and aliasing, but has a
    higher computational cost than the other filters.}
}
\paragraph{Caching and memory requirements:}
When a texture is read, Mitsuba internally converts it into an uncompressed linear format
using a half precision (\code{float16})-based representation. This is convenient for
rendering but means that textures require copious amounts of memory (in particular, the
size of the occupied memory region might be orders of magnitude greater than that of the
original input file).

For instance, a basic 10 megapixel image requires as much as 76 MiB of memory! Loading,
color space transformation, and MIP map construction require up to several seconds in this case.
To reduce these overheads, Mitsuba 0.4.0 introduced MIP map caches. When a large
texture is loaded for the first time, a MIP map cache file with the name \emph{filename}\code{.mip}
is generated. This is essentially a verbatim copy of the in-memory representation created
during normal rendering. Storing this information as a separate file has two advantages:

\begin{enumerate}[(i)]
   \item MIP maps do not have to be regenerated in subsequent Mitsuba runs,
    which substantially reduces scene loading times.
   \item Because the texture storage is entirely disk-backed and can be \emph{memory-mapped},
   Mitsuba is able to work with truly massive textures that would otherwise exhaust the main system memory.
\end{enumerate}

The texture caches are automatically regenerated when the input texture is modified.
Of course, the cache files can be cumbersome when they are not needed anymore. On Linux
or Mac OS, they can safely be deleted by executing the following command within a scene directory.

\begin{shell}
$\code{\$}$ find . -name "*.mip" -delete
\end{shell}
\plugin{checkerboard}{Checkerboard}
\order{2}
\parameters{
    \parameter{color0, color1}{\Spectrum}{
      Color values for the two differently-colored patches
      \default{0.4 and 0.2}
    }
    \parameter{uoffset, voffset}{\Float}{
      Numerical offset that should be applied to UV values before a lookup
    }
    \parameter{uscale, vscale}{\Float}{
      Multiplicative factors that should be applied to UV values before a lookup
    }
}
\renderings{
    \rendering{Checkerboard applied to the material test object
               as well as the ground plane}{tex_checkerboard}
}
This plugin implements a simple procedural checkerboard texture with
customizable colors.
\plugin{gridtexture}{Procedural grid texture}
\order{2}
\parameters{
    \parameter{color0}{\Spectrum}{
      Color values of the background
      \default{0.2}
    }
    \parameter{color1}{\Spectrum}{
      Color value of the lines
      \default{0.4}
    }
    \parameter{lineWidth}{\Float}{Width of the grid lines in UV space
       \default{0.01}
    }
    \parameter{uscale, vscale}{\Float}{
      Multiplicative factors that should be applied to UV values before a lookup
    }
    \parameter{uoffset, voffset}{\Float}{
      Numerical offset that should be applied to UV values before a lookup
    }
}
\renderings{
    \rendering{Grid texture applied to the material test object}{tex_gridtexture}
}
This plugin implements a simple procedural grid texture with customizable
colors and line width.
\plugin{scale}{Scaling passthrough texture}
\order{3}
\parameters{
    \parameter{value}{\Spectrum\Or\Texture}{
      Specifies the spectrum or nested texture that should be scaled
    }
    \parameter{value}{\Float}{
      Specifies the scale value
    }
}

This simple plugin wraps a nested texture plugin and multiplies its
contents by a user-specified value. This can be quite useful when a
texture is too dark or too bright. The plugin can also be used to adjust
the height of a bump map when using the \pluginref{bumpmap} plugin.

\begin{xml}[caption=Scaling the contents of a bitmap texture]
<texture type="scale">
    <float name="scale" value="0.5"/>

    <texture type="bitmap">
        <string name="filename" value="wood.jpg"/>
    </texture>
</texture>
\end{xml}
\plugin{vertexcolors}{Vertex color passthrough texture}
\order{5}
When rendering with a mesh that contains vertex colors,
this plugin exposes the underlying color data as a texture.
Currently, this is only supported by the \code{PLY}
file format loader.

Here is an example:
\begin{xml}[caption=Rendering a PLY file with vertex colors]
<shape type="ply">
    <string name="filename" value="mesh.ply"/>

    <bsdf type="diffuse">
        <texture type="vertexcolors" name="reflectance"/>
    </bsdf>
</shape>
\end{xml}
\plugin{wireframe}{Wireframe texture}
\order{6}
\parameters{
    \parameter{interiorColor}{\Spectrum}{
      Color value of the interior of triangles
      \default{0.5}
    }
    \parameter{edgeColor}{\Spectrum}{
      Edge color value
      \default{0.1}
    }
    \parameter{lineWidth}{\Float}{
       World-space width of the mesh edges
       \default{automatic}
    }
    \parameter{stepWidth}{\Float}{
       Controls the width of the step function used for the
       color transition. It is specified as a value between zero
       and one (relative to the \code{lineWidth} parameter)
       \default{0.5}
    }
}
\renderings{
    \rendering{Wireframe texture applied to the material test object}{tex_wireframe}
}

This plugin implements a simple two-color wireframe texture map
that reveals the structure of a triangular mesh.
\plugin{curvature}{Curvature texture}
\order{7}
\parameters{
    \parameter{curvature}{\String}{
      Specifies what should be shown -- must be equal to
      \code{mean} or \code{gaussian}.
    }
    \parameter{scale}{\Float}{
       A scale factor to bring curvature values into the
       displayable range [-1, 1]. Everything outside of this range
       will be clamped.
    }
}

\renderings{
    \rendering{Mean curvature}{tex_curvature_mean}
    \rendering{Gaussian curvature}{tex_curvature_gaussian}
}

This texture can visualize the mean and Gaussian curvature of the underlying
shape for inspection. Red and blue denote positive and negative values,
respectively.
\input{section_subsurface}
\plugin{dipole}{Dipole-based subsurface scattering model}
\parameters{
    \parameter{material}{\String}{
        Name of a material preset, see
        \tblref{medium-coefficients}. \default{\texttt{skin1}}
    }
    \parameter{sigmaA, sigmaS}{\Spectrum}{
        Absorption and scattering
        coefficients of the medium in inverse scene units.
        These parameters are mutually exclusive with \code{sigmaT} and \code{albedo}
        \default{configured based on \code{material}}
    }
    \parameter{sigmaT, albedo}{\Spectrum}{
        Extinction coefficient in inverse scene units
        and a (unitless) single-scattering albedo.
        These parameters are mutually exclusive with \code{sigmaA} and \code{sigmaS}
        \default{configured based on \code{material}}
    }
    \parameter{scale}{\Float}{
        Optional scale factor that will be applied to the \code{sigma*} parameters.
        It is provided for convenience when accomodating data based on different units,
        or to simply tweak the density of the medium. \default{1}}
    \parameter{intIOR}{\Float\Or\String}{Interior index of refraction specified
     numerically or using a known material name. \default{based on \code{material}}}
    \parameter{extIOR}{\Float\Or\String}{Exterior index of refraction specified
     numerically or using a known material name. \default{based on \code{material}}}
    \parameter{irrSamples}{\Integer}{
        Number of samples to use when estimating the
        irradiance at a point on the surface \default{16}
    }
}

\renderings{
   \rendering{The material test ball rendered with the \code{skimmilk}
   material preset}{subsurface_dipole.jpg}
   \rendering{The material test ball rendered with the \code{skin1}
   material preset}{subsurface_dipole_2.jpg}
}
\renderings{
   \rendering{\code{scale=1}}{subsurface_dipole_dragon.jpg}
   \rendering{\code{scale=0.2}}{subsurface_dipole_dragon2.jpg}
   \caption{The dragon model rendered with the \code{skin2}
   material preset (model courtesy of XYZ RGB). The \code{scale}
   parameter is useful to communicate the relative size of
   an object to the viewer.}
}

This plugin implements the classic dipole subsurface scattering model
from radiative transport and medical physics \cite{Eason1978Theory,
Farrell1992Diffusion} in the form proposed by Jensen et al.
\cite{Jensen2001Practical}. It relies on the assumption that light entering
a material will undergo many (i.e. hundreds) of internal scattering
events, such that diffusion theory becomes applicable. In this
case\footnote{and after making several fairly strong simplifications:
the geometry is assumed to be a planar half-space, and the internal
scattering from the material boundary is only considered approximately.}
a simple analytic solution of the subsurface scattering profile is available
that enables simulating this effect without having to account for the vast
numbers of internal scattering events individually.

For each \code{dipole} instance in the scene, the plugin adds a pre-process pass
to the rendering that computes the irradiance on a large set of sample positions
spread uniformly over the surface in question. The locations of these
points are chosen using a technique by Bowers et al. \cite{Bowers2010Parallel}
that creates particularly well-distributed (blue noise) samples. Later during
rendering, these  illumination samples are convolved with the diffusion profile
using a fast hierarchical technique proposed by Jensen and Buhler \cite{Jensen2005Rapid}.

There are two different ways of configuring the medium properties.
One possibility is to load a material preset
using the \code{material} parameter---see \tblref{medium-coefficients}
for details. Alternatively, when specifying parameters by hand, they
can either be provided using the scattering and absorption coefficients,
or by declaring the extinction coefficient and single scattering albedo
(whichever is more convenient). Mixing these parameter initialization
methods is not allowed.

All scattering parameters (named \code{sigma*}) should
be provided in inverse scene units. For instance, when a world-space
distance of 1 unit corresponds to a meter, the scattering coefficents must
be in units of inverse meters. For convenience, the \code{scale}
parameter can be used to correct this. For instance, when the scene is
in meters and the coefficients are in inverse millimeters, set
\code{scale=1000}.

Note that a subsurface integrator can be associated with an \code{id}
and shared by several shapes using the reference mechanism introduced in
\secref{format}. This can be useful when an object is made up of many
separate sub-shapes.

\renderings{
   \medrendering{Rendered using \pluginref{dipole}}{subsurface_dipole_bad1.jpg}
   \medrendering{Rendered using \pluginref{homogeneous}}{subsurface_dipole_bad2.jpg}
   \medrendering{\code{irrSamples} set too low}{subsurface_dipole_bad3.jpg}
   \caption{Two problem cases that may occur when rendering with the \pluginref{dipole}:
    \textbf{(a)-(b)}: These two renderings show a glass ball filled with diluted milk
    rendered using diffusion theory and radiative transport, respectively.
    The former produces an incorrect result, since the assumption of
    many scattering events breaks down.
    \textbf{(c)}: When the number of irradiance samples is too low when rendering
    with the dipole model, the resulting noise becomes visible as ``blotchy'' artifacts
    in the rendering.}
}

\subsubsection*{Typical material setup}
To create a realistic material with subsurface scattering, it is necessary
to associate the underlying shape with an appropriately configured surface
and subsurface scattering model. Both should be aware of the material's
index of refraction.

Because the \pluginref{dipole} plugin is responsible for all internal
scattering, the surface scattering model should only account for specular
reflection due to the index of refraction change. There are two models
in Mitsuba that can do this: \pluginref{plastic} and
\pluginref{roughplastic} (for smooth and rough interfaces, respectively).
An example is given on the next page.
\pagebreak
\begin{xml}
<shape type="...">
    <subsurface type="dipole">
        <string name="intIOR" value="water"/>
        <string name="extIOR" value="air"/>
        <rgb name="sigmaS" value="87.2, 127.2, 143.2"/>
        <rgb name="sigmaA" value="1.04, 5.6, 11.6"/>
        <integer name="irrSamples" value="64"/>
    </subsurface>

    <bsdf type="plastic">
        <string name="intIOR" value="water"/>
        <string name="extIOR" value="air"/>
        <!-- Note: the diffuse component must be disabled! -->
        <spectrum name="diffuseReflectance" value="0"/>
    </bsdf>
<shape>
\end{xml}

\remarks{
   \item This plugin only implements the multiple scattering component of
   the dipole model, i.e. single scattering is omitted. Furthermore, the
   numerous assumptions built into the underlying theory can cause severe
   inaccuracies.

   For this reason, this plugin is the right choice for making pictures
   that ``look nice'', but it should be avoided when the output must hold
   up to real-world measurements. In this case, please use participating media
   (\secref{media}).

  \item It is quite important that the \code{sigma*} parameters have the right units.
  For instance: if the \code{sigmaT} parameter is accidentally set to a value that
  is too small by a factor of 1000, the plugin will attempt to create
  one million times as many irradiance samples, which will likely cause
  the rendering process to crash with an ``out of memory'' failure.
}
\plugin{singlescatter}{Single scattering in participating media}
\parameters{
    \parameter{material}{\String}{
        Name of a material preset, see
        \tblref{medium-coefficients}. \default{\texttt{skin1}}
    }
    \parameter{sigmaA, sigmaS}{\Spectrum}{
        Absorption and scattering
        coefficients of the medium in inverse scene units.
        These parameters are mutually exclusive with \code{sigmaT} and \code{albedo}
        \default{configured based on \code{material}}
    }
    \parameter{sigmaT, albedo}{\Spectrum}{
        Extinction coefficient in inverse scene units
        and a (unitless) single-scattering albedo.
        These parameters are mutually exclusive with \code{sigmaA} and \code{sigmaS}
        \default{configured based on \code{material}}
    }
    \parameter{scale}{\Float}{
        Optional scale factor that will be applied to the \code{sigma*} parameters.
        It is provided for convenience when accomodating data based on different units,
        or to simply tweak the density of the medium. \default{1}}

}

\renderings{
   \rendering{The bumpy sphere test scene rendered with amber material}{bumpy_sphere.jpg}
}

This plugin implements the single scattering model in participating media
\cite{Holzschuch2015}.

There are two different ways of configuring the medium properties.
One possibility is to load a material preset
using the \code{material} parameter---see \tblref{medium-coefficients}
for details. Alternatively, when specifying parameters by hand, they
can either be provided using the scattering and absorption coefficients,
or by declaring the extinction coefficient and single scattering albedo
(whichever is more convenient). Mixing these parameter initialization
methods is not allowed.

All scattering parameters (named \code{sigma*}) should
be provided in inverse scene units. For instance, when a world-space
distance of 1 unit corresponds to a meter, the scattering coefficents must
be in units of inverse meters. For convenience, the \code{scale}
parameter can be used to correct this. For instance, when the scene is
in meters and the coefficients are in inverse millimeters, set
\code{scale=1000}.

Note that a subsurface integrator can be associated with an \code{id}
and shared by several shapes using the reference mechanism introduced in
\secref{format}. This can be useful when an object is made up of many
separate sub-shapes.


\remarks{
   \item This plugin implements only the single scattering
   component.
   Single scattering comes in two flavors: fast, the default straight line
   approximation (with multiple samples along the line), and slow, that is
   the accurate version described in \cite{Holzschuch2015}.

  \item It is quite important that the \code{sigma*} parameters have the right units.
  For instance: if the \code{sigmaT} parameter is accidentally set to a value that
  is too small by a factor of 1000, the plugin will attempt to create
  one million times as many irradiance samples, which will likely cause
  the rendering process to crash with an ``out of memory'' failure.
}
\input{section_media}
\plugin{homogeneous}{Homogeneous participating medium}
\order{1}
\parameters{
    \parameter{material}{\String}{
        Name of a material preset, see
        \tblref{medium-coefficients}. \default{\texttt{skin1}}
    }
    \parameter{sigmaA, sigmaS}{\Spectrum}{
        Absorption and scattering
        coefficients of the medium in inverse scene units.
        These parameters are mutually exclusive with \code{sigmaT} and \code{albedo}
        \default{configured based on \code{material}}
    }
    \parameter{sigmaT, albedo}{\Spectrum}{
        Extinction coefficient in inverse scene units
        and a (unitless) single-scattering albedo.
        These parameters are mutually exclusive with \code{sigmaA} and \code{sigmaS}
        \default{configured based on \code{material}}
    }
    \parameter{\footnotesize{scale}}{\Float}{
        Optional scale factor that will be applied to the \code{sigma*} parameters.
        It is provided for convenience when accomodating data based on different units,
        or to simply tweak the density of the medium. \default{1}
    }
    \parameter{\Unnamed}{\Phase}{
         A nested phase function that describes the directional
         scattering properties of the medium. When none is specified,
         the renderer will automatically use an instance of
         \pluginref{isotropic}.
    }
}

This class implements a flexible homogeneous participating
medium with support for arbitrary phase functions and various
medium sampling methods. It provides two different ways of configuring
the medium properties. One possibility is to load a material preset
using the \code{material} parameter---see \tblref{medium-coefficients}
for details. Alternatively, when specifying parameters by hand, they can either
be provided using the scattering and absorption coefficients, or
by declaring the extinction coefficient and single scattering
albedo (whichever is more convenient). Mixing these parameter
initialization methods is not allowed.

All scattering parameters (named \code{sigma*}) should
be provided in inverse scene units. For instance, when a world-space
distance of 1 unit corresponds to a meter, the scattering coefficents should
have units of inverse meters. For convenience, the \code{scale}
parameter can be used to correct the units. For instance, when the scene is
in meters and the coefficients are in inverse millimeters, set
\code{scale} to \code{1000}.

\renderings{
   \rendering{A squishy ball rendered with subsurface scattering and
   a dielectric BSDF (courtesy of Chanxi Zheng)}{medium_homogeneous_squishy.jpg}
}

\begin{xml}[caption=Declaration of a forward scattering medium with high albedo]
<medium id="myMedium" type="homogeneous">
    <spectrum name="sigmaS" value="1"/>
    <spectrum name="sigmaA" value="0.05"/>

    <phase type="hg">
        <float name="g" value="0.7"/>
    </phase>
</medium>
\end{xml}

\textbf{Note}: Rendering media that have a spectrally
varying extinction coefficient can be tricky, since all
commonly used medium sampling methods suffer from high
variance in that case. Here, it may often make more sense to render
several monochromatic images separately (using only the coefficients for
a single channel) and then merge them back into a RGB image. There
is a \code{mtsutil} (\secref{mtsutil}) plugin named \code{joinrgb}
that will perform this RGB merging process.

\begin{table}[h!]
    \centering
    \vspace{3mm}
    {\footnotesize
    \begin{tabular}{>{\ttfamily}p{3.8cm}p{.4cm}>{\ttfamily}p{3.8cm}p{.4cm}>{\ttfamily}p{3.8cm}}
        \toprule
        \rmfamily \small\textbf{Name} &&
        \rmfamily \small\textbf{Name} &&
        \rmfamily \small\textbf{Name} \\
        \cmidrule{1-1} \cmidrule{3-3} \cmidrule{5-5}
        Apple && Chicken1 && Chicken2 \\
        Cream && Ketchup  && Potato \\
        Skimmilk && Skin1 && Skin2 \\
        Spectralon && Wholemilk && \\
        \cmidrule{1-1} \cmidrule{3-3} \cmidrule{5-5}
        Lowfat Milk              &&  Gatorade                &&    White Grapefruit Juice     \\
        Reduced Milk             &&  Chardonnay              &&    Shampoo                    \\
        Regular Milk             &&  White Zinfandel         &&    Strawberry Shampoo         \\
        Espresso                 &&  Merlot                  &&    \mbox{Head \& Shoulders Shampoo}  \\
        Mint Mocha Coffee        &&  Budweiser Beer          &&    Lemon Tea Powder           \\
        Lowfat Soy Milk          &&  Coors Light Beer        &&    Orange Juice Powder        \\
        Regular Soy Milk         &&  Clorox                  &&    Pink Lemonade Powder       \\
        Lowfat Chocolate Milk    &&  Apple Juice             &&    Cappuccino Powder          \\
        Regular Chocolate Milk   &&  Cranberry Juice         &&    Salt Powder                \\
        Coke                     &&  Grape Juice             &&    Sugar Powder               \\
        Pepsi                    &&  Ruby Grapefruit Juice   &&    Suisse Mocha               \\
        Sprite                   &&                          &&                               \\
        \bottomrule
    \end{tabular}}
    \caption{\label{tbl:medium-coefficients}This
         table lists all supported medium material presets. The
         top entries are from Jensen et al. \cite{Jensen2001Practical}, and the
         bottom ones are from Narasimhan et al. \cite{Narasimhan2006Acquiring}.
         They all use units of $\frac{1}{mm}$, so remember to set
         \code{scale} appropriately when your scene is not
         in units of millimeters.
         These material presets can be used with the plugins
         \pluginref{homogeneous},\
         \pluginref{dipole}, and \
         \pluginref{hk}
    }
\end{table}
\plugin{heterogeneous}{Heterogeneous participating medium}
\order{2}
\parameters{
    \parameter{method}{\String}{
        Specifies the sampling method that is used to generate
        scattering events within the medium.
        \begin{enumerate}[(i)]
            \item \code{simpson}: Sampling is done by inverting a
            deterministic quadrature rule based on composite
            Simpson integration over small ray segments. Benefits
            from the use of good sample generators (e.g. \pluginref{ldsampler}).
            \item \code{woodcock}: Generate samples using
            Woodcock tracking. This is usually faster and
            always unbiased, but has the disadvantages of not benefiting
            from good sample generators and not providing
            information that is required by bidirectional
            rendering techniques.
        \end{enumerate}
        Default: \texttt{woodcock}
    }
    \parameter{density}{\Volume}{
        Volumetric data source that supplies the medium densities
        (in inverse scene units)
    }
    \parameter{albedo}{\Volume}{
        Volumetric data source that supplies the
        single-scattering albedo
    }
    \parameter{orientation}{\Volume}{
        Optional: volumetric data source that supplies the
        local particle orientations throughout the medium
    }
    \parameter{scale}{\Float}{
        Optional scale factor that will be applied to the \code{density} parameter.
        Provided for convenience when accomodating data based on different units,
        or to simply tweak the density of the medium. \default{1}
    }
    \parameter{\Unnamed}{\Phase}{
         A nested phase function that describes the directional
         scattering properties of the medium. When none is specified,
         the renderer will automatically use an instance of
         \pluginref{isotropic}.
    }
}

\renderings{
    \medrendering{40}{medium_heterogeneous_density_40}
    \medrendering{200}{medium_heterogeneous_density_200}
    \medrendering{1000}{medium_heterogeneous_density_1000}
    \vspace{-2mm}
    \caption{Renderings of an index-matched medium using different scale factors (\lstref{hetvolume})}
}

This plugin provides a flexible heterogeneous medium implementation, which
acquires its data from nested \code{volume} instances. These can be
constant, use a procedural function, or fetch data from disk, e.g. using a
memory-mapped density grid. See \secref{volumes} for details on volume data
sources.

Instead of allowing separate volumes to be provided for the scattering
and absorption parameters \code{sigmaS} and \code{sigmaA} (as is done in
\pluginref{homogeneous}), this class instead takes the approach of
enforcing a spectrally uniform value of \code{sigmaT}, which must be
provided using a nested scalar-valued volume named \code{density}.

Another nested spectrum-valued \code{albedo} volume must also be provided, which is
used to compute the scattering coefficient $\sigma_s$ using the expression
$\sigma_s = \code{scale} * \code{density} * \code{albedo}$ (i.e. 'albedo' contains the
single-scattering albedo of the medium.

Optionally, one can also provide an vector-valued \code{orientation} volume,
which contains local particle orientation that will be passed to
scattering models that support this, such as a the Micro-flake or
Kajiya-Kay phase functions.

\vspace{4mm}

\begin{xml}[label=lst:hetvolume,caption=A simple heterogeneous medium backed by a grid volume]
<!-- Declare a heterogeneous participating medium named 'smoke' -->
<medium type="heterogeneous" id="smoke">
    <string name="method" value="simpson"/>

    <!-- Acquire density values from an external data file -->
    <volume name="density" type="gridvolume">
        <string name="filename" value="frame_0150.vol"/>
    </volume>

    <!-- The albedo is constant and set to 0.9 -->
    <volume name="albedo" type="constvolume">
        <spectrum name="value" value="0.9"/>
    </volume>

    <!-- Use an isotropic phase function -->
    <phase type="isotropic"/>

    <!-- Scale the density values as desired -->
    <float name="scale" value="200"/>
 </medium>

<!-- Attach the index-matched medium to a shape in the scene -->
<shape type="obj">
    <!-- Load an OBJ file, which contains a mesh version
         of the axis-aligned box of the volume data file -->
    <string name="filename" value="bounds.obj"/>

    <!-- Reference the medium by ID -->
    <ref name="interior" id="smoke"/>

    <!-- If desired, this shape could also declare
         a BSDF to create an index-mismatched
         transition, e.g.

    <bsdf type="dielectric"/>
    -->
</shape>
\end{xml}
\input{section_phase}
\plugin{isotropic}{Isotropic phase function}
\order{1}

\renderings{
    \rendering{Isotropic}{phase_isotropic}
    \rendering{Anisotropic micro-flakes}{phase_microflakes_005}
    \caption{Heterogeneous volume renderings of a scarf model
    with isotropic and anisotropic phase functions.}
}


This phase function simulates completely uniform scattering,
where all directionality is lost after a single scattering
interaction. It does not have any parameters.
\plugin{hg}{Henyey-Greenstein phase function}
\order{2}
\parameters{
    \parameter{g}{\Float}{
      This parameter must be somewhere in the range $-1$ to $1$
      (but not equal to $-1$ or $1$). It denotes the \emph{mean cosine}
      of scattering interactions. A value greater than zero indicates that
      medium interactions predominantly scatter incident light into a similar
      direction (i.e. the medium is \emph{forward-scattering}), whereas
      values smaller than zero cause the medium to be
      scatter more light in the opposite direction.
    }
}
This plugin implements the phase function model proposed by
Henyey and Greenstein \cite{Henyey1941Diffuse}. It is
parameterizable from backward- ($g<0$) through
isotropic- ($g=0$) to forward ($g>0$) scattering.
\plugin{rayleigh}{Rayleigh phase function}
\order{3}

Scattering by particles that are much smaller than the wavelength
of light (e.g. individual molecules in the atmosphere) is well-approximated
by the Rayleigh phase function. This plugin implements an unpolarized
version of this scattering model (i.e the effects of polarization are ignored).
This plugin is useful for simulating scattering in planetary atmospheres.

This model has no parameters.
\plugin{kkay}{Kajiya-Kay phase function}
This plugin implements the Kajiya-Kay \cite{Kajiya1989Rendering}
phase function for volumetric rendering of fibers, e.g.
hair or cloth.

The function is normalized so that it has no energy loss when
\code{ks}=1 and illumination arrives perpendicularly to the surface.
\plugin{microflake}{Micro-flake phase function}
\parameters{
    \parameter{stddev}{\Float}{
      Standard deviation of the micro-flake normals. This
      specifies the roughness of the fibers in the medium.
    }
}

\renderings{
    \rendering{\code{stddev}=0.2}{phase_microflakes_02}
    \rendering{\code{stddev}=0.05}{phase_microflakes_005}
}

This plugin implements the anisotropic micro-flake phase function
described in
``A radiative transfer framework for rendering materials with
anisotropic structure'' by Wenzel Jakob, Adam Arbree,
Jonathan T. Moon, Kavita Bala, and Steve Marschner
\cite{Jakob2010Radiative}.

The implementation in this plugin is specific to rough fibers
and uses a Gaussian-type flake distribution. It is much faster
than the spherical harmonics approach proposed in the original
paper. This distribution, as well as the implemented sampling
method, are described in the paper
``Building Volumetric Appearance Models of Fabric using
Micro CT Imaging'' by Shuang Zhao, Wenzel Jakob, Steve Marschner,
and Kavita Bala \cite{Zhao2011Building}.

Note: this phase function must be used with a medium that specifies
the local fiber orientation at different points in space. Please
refer to \pluginref{heterogeneous} for details.

\vspace{1cm}

\renderings{
  \fbox{\includegraphics[width=\textwidth]{images/medium_knitpatterns}}
  \caption{A range of different knit patterns, rendered using the
  \pluginref{heterogeneous} and \pluginref{microflake} plugins. Courtesy
  of Yuksel et al. \cite{Yuksel2012Stitch}.}
}
\plugin{mixturephase}{Mixture phase function}

\parameters{
    \parameter{weights}{\String}{A comma-separated list of phase function weights}
    \parameter{\Unnamed}{\Phase}{Multiple phase function instances that should be
    mixed according to the specified weights}
}

This plugin implements a ``mixture'' scattering model analogous to \pluginref{mixturebsdf},
which represents linear combinations of multiple phase functions. There is no
limit on how many phase function can be mixed, but their combination
weights must be non-negative and sum to a value of one or less to ensure
energy balance.
\input{section_volumes}
\plugin{constvolume}{Constant-valued volume data source}
\parameters{
    \parameter{value}{\Float\Or\Spectrum\Or\Vector}{
      Specifies the value of the volume
    }
}

This plugin provides a volume data source that is
constant throughout its domain. Depending on how it is used,
its value can either be a scalar, a color spectrum,
or a 3D vector.\vspace{4mm}

\begin{xml}[caption={Definition of a heterogeneous medium with homogeneous contents}]
<medium type="heterogeneous">
    <volume type="constvolume" name="density">
        <float name="value" value="1"/>
    </volume>
    <volume type="constvolume" name="albedo">
        <rgb name="value" value="0.9 0.9 0.7"/>
    </volume>
    <volume type="constvolume" name="orientation">
        <vector name="value" x="0" y="1" z="0"/>
    </volume>

    <!-- .... remaining parameters for
         the 'heterogeneous' plugin .... -->
</medium>
\end{xml}
\plugin{gridvolume}{Grid-based volume data source}
\parameters{
    \parameter{filename}{\String}{
      Specifies the filename of the volume data file to be loaded
    }
    \parameter{sendData}{\Boolean}{
      When this parameter is set to \code{true}, the implementation will
      send all volume data to other network render nodes. Otherwise, they
      are expected to have access to an identical volume data file that can be
      mapped into memory. \default{\code{false}}
    }
    \parameter{toWorld}{\Transform}{
        Optional linear transformation that should be applied to the data
    }
    \parameter{min, max}{\Point}{
        Optional parameter that can be used to re-scale the data so that
        it lies in the bounding box between \code{min} and \code{max}.
    }
}

This class implements access to memory-mapped volume data stored on a
3D grid using a simple binary exchange format.
The format uses a little endian encoding and is specified as
follows:\vspace{3mm}

\begin{center}
\begin{tabular}{>{\bfseries}p{2cm}p{11cm}}
\toprule
Position & Content\\
\midrule
Bytes 1-3&   ASCII Bytes '\code{V}', '\code{O}', and '\code{L}' \\
Byte  4&     File format version number (currently 3)\\
Bytes 5-8&   Encoding identifier (32-bit integer). The following
choices are available:
\begin{enumerate}[1.]
\item Dense \code{float32}-based representation
\item Dense \code{float16}-based representation (\emph{currently not supported by this implementation})
\item Dense \code{uint8}-based representation (The range 0..255 will be mapped to 0..1)
\item Dense quantized directions. The directions are stored in spherical
coordinates with a total storage cost of 16 bit per entry.
\end{enumerate}\\
Bytes 9-12 &  Number of cells along the X axis (32 bit integer)\\
Bytes 13-16 &  Number of cells along the Y axis (32 bit integer)\\
Bytes 17-20 &  Number of cells along the Z axis (32 bit integer)\\
Bytes 21-24 &  Number of channels (32 bit integer, supported values: 1 or 3)\\
Bytes 25-48 &  Axis-aligned bounding box of the data stored in single
               precision (order: xmin, ymin, zmin, xmax, ymax, zmax)\\
Bytes 49-*  &  Binary data of the volume stored in the specified encoding.
               The data are ordered so that the following C-style indexing
               operation makes sense after the file has been mapped into memory:\newline
                  \ \ \code{data[((zpos*yres + ypos)*xres + xpos)*channels + chan]}\newline
               where \code{(xpos, ypos, zpos, chan)} denotes the lookup location.\\

\bottomrule
\end{tabular}
\end{center}

Note that Mitsuba expects that entries in direction volumes are either
zero or valid unit vectors.

When using this data source to represent floating point density volumes,
please ensure that the values are all normalized to lie in the
range $[0, 1]$---otherwise, the Woodcock-Tracking integration method in
\pluginref{heterogeneous} will produce incorrect results.
\plugin{volcache}{Caching volume data source}
\parameters{
    \parameter{blockSize}{\Integer}{
        Size of the individual cache blocks
        \default{8, i.e. $8\times8\times 8$}
    }
    \parameter{voxelWidth}{\Float}{
        Width of a voxel (in a cache block) expressed in
        world-space units. \default{set to the ray marching
        step size of the nested medium}
    }
    \parameter{memoryLimit}{\Integer}{
        Maximum allowed memory usage in MiB. \default{1024, i.e. 1 GiB}
    }
    \parameter{toWorld}{\Transform}{
        Optional linear transformation that should be applied
        to the volume data
    }
    \parameter{\Unnamed}{\Volume}{
        A nested volume data source
    }
}

This plugin can be added between the renderer and another
data source, for which it caches all data lookups using a
LRU scheme. This is useful when the nested volume data source
is expensive to evaluate.

The cache works by performing on-demand rasterization of subregions
of the nested volume into blocks ($8\times 8 \times 8$ by default).
These are kept in memory until a user-specifiable threshold is exeeded,
after which point a \emph{least recently used} (LRU) policy removes
records that haven't been accessed in a long time.
\input{section_emitters}
\plugin{point}{Point light source}
\icon{emitter_point}
\order{1}
\parameters{
    \parameter{toWorld}{\Transform\Or\Animation}{
       Specifies an optional sensor-to-world transformation.
       \default{none (i.e. sensor space $=$ world space)}
    }
    \parameter{position}{\Point}{
       Alternative parameter for specifying the light source
       position. Note that only one of the parameters
       \code{toWorld} and \code{position} can be used at a time.
    }
    \parameter{intensity}{\Spectrum}{
        Specifies the radiant intensity in units of
        power per unit steradian.
        \default{1}
    }
    \parameter{samplingWeight}{\Float}{
        Specifies the relative amount of samples
        allocated to this emitter. \default{1}
    }
}

This emitter plugin implements a simple point light source, which
uniformly radiates illumination into all directions.
\plugin{area}{Area light}
\icon{emitter_area}
\order{2}
\parameters{
    \parameter{radiance}{\Spectrum}{
        Specifies the emitted radiance in units of
        power per unit area per unit steradian.
    }
    \parameter{samplingWeight}{\Float}{
        Specifies the relative amount of samples
        allocated to this emitter. \default{1}
    }
}

This plugin implements an area light, i.e. a light source that emits
diffuse illumination from the exterior of an arbitrary shape.
Since the emission profile of an area light is completely diffuse, it
has the same apparent brightness regardless of the observer's viewing
direction. Furthermore, since it occupies a nonzero amount of space, an
area light generally causes scene objects to cast soft shadows.

When modeling scenes involving area lights, it is preferable
to use spheres as the emitter shapes, since they provide a
particularly good direct illumination sampling strategy (see
the \pluginref{sphere} plugin for an example).

To create an area light source, simply instantiate the desired
emitter shape and specify an \code{area} instance as its child:

\vspace{4mm}
\begin{xml}
<!-- Create a spherical light source at the origin -->
<shape type="sphere">
    <emitter type="area">
        <spectrum name="radiance" value="1"/>
    </emitter>
</shape>
\end{xml}
\plugin{spot}{Spot light source}
\icon{emitter_spot}
\order{3}
\parameters{
    \parameter{toWorld}{\Transform\Or\Animation}{
       Specifies an optional sensor-to-world transformation.
       \default{none (i.e. sensor space $=$ world space)}
    }
    \parameter{intensity}{\Spectrum}{
        Specifies the maximum radiant intensity at the center
        in units of power per unit steradian.
        \default{1}
    }
    \parameter{cutoffAngle}{\Float}{Cutoff angle, beyond which the spot light is completely black \default{\code{20} degrees}}
    \parameter{beamWidth}{\Float}{Subtended angle of the central beam portion \default{\code{cutoffAngle}$\ \cdot\ \nicefrac 34$}}
    \parameter{texture}{\Texture}{
        An optional texture to be projected along the spot light
    }
    \parameter{samplingWeight}{\Float}{
        Specifies the relative amount of samples
        allocated to this emitter. \default{1}
    }
}

This plugin provides a spot light with a linear falloff.
In its local coordinate system, the spot light is positioned at the origin
and points along the positive Z direction. It can be conveniently
reoriented using the \code{lookat} tag, e.g.:
\begin{xml}
<emitter type="spot">
    <transform name="toWorld">
        <!-- Orient the light so that points from (1, 1, 1) towards (1, 2, 1) -->
        <lookat origin="1, 1, 1" target="1, 2, 1"/>
    </transform>
</emitter>
\end{xml}

The intensity linearly ramps up from \code{cutoffAngle}
to \code{beamWidth} (both specified in degrees), after which it remains at
the maximum value. A projection texture may optionally be supplied.
\plugin{directional}{Directional emitter}
\icon{emitter_directional}
\order{4}
\parameters{
    \parameter{toWorld}{\Transform\Or\Animation}{
       Specifies an optional emitter-to-world transformation.
       \default{none (i.e. emitter space $=$ world space)}
    }
    \parameter{direction}{\Vector}{
       Alternative to \code{toWorld}: explicitly specifies
       the illumination direction. Note that only one of the
       two parameters can be used.
    }
    \parameter{irradiance}{\Spectrum}{
        Specifies the amount of power per unit area received
        by a hypothetical surface normal to the specified direction
        \default{1}
    }
    \parameter{samplingWeight}{\Float}{
        Specifies the relative amount of samples
        allocated to this emitter. \default{1}
    }
}

This emitter plugin implements a distant directional source, which
radiates a specified power per unit area along a fixed direction.
By default, the emitter radiates in the direction of the postive Z axis.
\plugin{collimated}{Collimated beam emitter}
\icon{emitter_collimated}
\order{5}
\parameters{
    \parameter{toWorld}{\Transform\Or\Animation}{
       Specifies an optional emitter-to-world transformation.
       \default{none (i.e. emitter space $=$ world space)}
    }
    \parameter{power}{\Spectrum}{
        Specifies the amount of power radiated along the beam
        \default{1}
    }
    \parameter{samplingWeight}{\Float}{
        Specifies the relative amount of samples
        allocated to this emitter. \default{1}
    }
}

This emitter plugin implements a collimated beam source, which
radiates a specified amount of power along a fixed ray.
It can be thought of as the limit of a spot light as its field
of view tends to zero.

Such a emitter is useful for conducting virtual experiments and
testing the renderer for correctness.

By default, the emitter is located at the origin and radiates
into the positive Z direction $(0,0,1)$. This can
be changed by providing a custom \code{toWorld} transformation.
\plugin{sky}{Skylight emitter}
\icon{emitter_sky}
\order{6}
\parameters{
    \parameter{turbidity}{\Float}{
        This parameter determines the amount of aerosol present
        in the atmosphere.
        Valid range: 1-10. \default{3, corresponding to a clear sky in a temperate climate}
    }
    \parameter{albedo}{\Spectrum}{Specifies the ground albedo \default{0.15}}
    \parameter{year, month, day}{\Integer}{Denote the date of the
     observation \default{2010, 07, 10}}
    \parameter{hour,minute,\showbreak second}{\Float}{Local time
      at the location of the observer in 24-hour format\default{15, 00, 00,
      i.e. 3PM}}
    \parameter{latitude, longitude, timezone}{\Float}{
      These three parameters specify the oberver's latitude and longitude
      in degrees, and the local timezone offset in hours, which are required
      to compute the sun's position. \default{35.6894, 139.6917, 9 --- Tokyo, Japan}
    }
    \parameter{sunDirection}{\Vector}{Allows to manually
      override the sun direction in world space. When this value
      is provided, parameters pertaining to the computation
      of the sun direction (\code{year, hour, latitude,} etc.
      are unnecessary. \default{none}
    }
    \parameter{stretch}{\Float}{
        Stretch factor to extend emitter below the horizon, must be
        in [1,2] \default{\code{1}, i.e. not used}
    }
    \parameter{resolution}{\Integer}{Specifies the horizontal resolution of the precomputed
        image that is used to represent the sun environment map \default{512, i.e. 512$\times$256}}
    \parameter{scale}{\Float}{
        This parameter can be used to scale the amount of illumination
        emitted by the sky emitter. \default{1}
    }
    \parameter{samplingWeight}{\Float}{
        Specifies the relative amount of samples
        allocated to this emitter. \default{1}
    }
    \parameter{toWorld}{\Transform\Or\Animation}{
       Specifies an optional sensor-to-world transformation.
       \default{none (i.e. sensor space $=$ world space)}
    }
}

\renderings{
    \tinyrendering{5AM}{emitter_sky_small_5}
    \tinyrendering{7AM}{emitter_sky_small_7}
    \tinyrendering{9AM}{emitter_sky_small_9}
    \tinyrendering{11AM}{emitter_sky_small_11}
    \tinyrendering{1PM}{emitter_sky_small_13}
    \tinyrendering{3PM}{emitter_sky_small_15}
    \tinyrendering{5PM}{emitter_sky_small_17}
    \tinyrendering{6:30 PM}{emitter_sky_small_1830}\hfill
    \vspace{-2mm}
    \caption{Time series at the default settings (Equidistant fisheye
    projection of the sky onto a disk. East is left.)}
    \vspace{3mm}
}

This plugin provides the physically-based skylight model by
Ho\v{s}ek and Wilkie \cite{Hosek2012Analytic}. It can be used to
create predictive daylight renderings of scenes under clear skies,
which is useful for architectural and computer vision applications.
The implementation in Mitsuba is based on code that was
generously provided by the authors.

The model has two main parameters: the turbidity of the atmosphere
and the position of the sun.
The position of the sun in turn depends on a number of secondary
parameters, including the \code{latitude}, \code{longitude},
and \code{timezone} at the location of the observer, as well as the
current \code{year}, \code{month}, \code{day}, \code{hour},
\code{minute}, and \code{second}.
Using all of these, the elevation and azimuth of the sun are computed
using the PSA algorithm by Blanco et al. \cite{Blanco2001Computing},
which is accurate to about 0.5 arcminutes (\nicefrac{1}{120} degrees).
Note that this algorithm does not account for daylight
savings time where it is used, hence a manual correction of the
time may be necessary.
For detailed coordinate and timezone information of various cities, see
\url{http://www.esrl.noaa.gov/gmd/grad/solcalc}.

If desired, the world-space solar vector may also be specified
using the \code{sunDirection} parameter, in which case all of the
previously mentioned time and location parameters become irrelevant.

\renderings{
    \tinyrendering{1}{emitter_sky_turb_1}
    \tinyrendering{2}{emitter_sky_turb_2}
    \tinyrendering{3}{emitter_sky_turb_3}
    \tinyrendering{4}{emitter_sky_turb_4}
    \tinyrendering{5}{emitter_sky_turb_5}
    \tinyrendering{6}{emitter_sky_turb_6}
    \tinyrendering{8}{emitter_sky_turb_8}
    \tinyrendering{10}{emitter_sky_turb_10}
    \caption{Sky light for different turbidity values (default configuration at 5PM)}
}

\emph{Turbidity}, the other important parameter, specifies the aerosol
content of the atmosphere. Aerosol particles cause additional scattering that
manifests in a halo around the sun, as well as color fringes near the
horizon.
Smaller turbidity values ($\sim 1-2$) produce an
arctic-like clear blue sky, whereas larger values ($\sim 8-10$)
create an atmosphere that is more typical of a warm, humid day.
Note that this model does not aim to reproduce overcast, cloudy, or foggy
atmospheres with high corresponding turbidity values. An photographic
environment map may be more appropriate in such cases.

The default coordinate system of the emitter associates the up
direction with the $+Y$ axis. The east direction is associated with $+X$
and the north direction is equal to $+Z$. To change this coordinate
system, rotations can be applied using the \code{toWorld} parameter
(see \lstref{sky-up} for an example).

By default, the emitter will not emit any light below the
horizon, which means that these regions are black when
observed directly. By setting the \code{stretch} parameter to values
between $1$ and $2$, the sky can be extended to cover these directions
as well. This is of course a complete kludge and only meant as a quick
workaround for scenes that are not properly set up.

Instead of evaluating the full sky model every on every radiance query,
the implementation precomputes a low resolution environment map
(512$\times$ 256) of the entire sky that is then forwarded to the
\pluginref{envmap} plugin---this dramatically improves rendering
performance. This resolution is generally plenty since the sky radiance
distribution is so smooth, but it can be adjusted manually if
necessary using the \code{resolution} parameter.

Note that while the model encompasses sunrise and sunset configurations,
it does not extend to the night sky, where illumination from stars, galaxies,
and the moon dominate. When started with a sun configuration that lies
below the horizon, the plugin will fail with an error message.

\vspace{5mm}
\begin{xml}[caption={Rotating the sky emitter for scenes that use $Z$ as
the ``up'' direction}, label=lst:sky-up]
<emitter type="sky">
    <transform name="toWorld">
        <rotate x="1" angle="90"/>
    </transform>
</emitter>
\end{xml}

\subsubsection*{Physical units and spectral rendering}
Like the \code{blackbody} emission profile (Page~\pageref{sec:blackbody}),
the sky model introduces physical units into the rendering process.
The radiance values computed by this plugin have units of power ($W$) per
unit area ($m^{-2}$) per steradian ($sr^{-1}$) per unit wavelength ($nm^{-1}$).
If these units are inconsistent with your scene description, you may use the
optional \texttt{scale} parameter to adjust them.

When Mitsuba is compiled for spectral rendering, the plugin switches
from RGB to a spectral variant of the skylight model, which relies on
precomputed data between $320$ and $720 nm$ sampled at $40nm$-increments.

\subsubsection*{Ground albedo}
The albedo of the ground (e.g. due to rock, snow, or vegetation) can have a
noticeable and nonlinear effect on the appearance of the sky.
\figref{sky_groundalbedo} shows an example of this effect. By default,
the ground albedo is set to a 15\% gray.

\renderings{
   \rendering{3 PM}{emitter_sky_mattest_3pm}
   \rendering{6:30 PM}{emitter_sky_mattest_630pm}
   \vspace{-3mm}
   \caption{Renderings with the \pluginref{plastic} material
   under default conditions. Note that these only contain
   skylight illumination. For a model that also includes the sun,
   refer to \pluginref{sunsky}.}
   \vspace{-6mm}
}
\vspace{5mm}
\renderings{
  \medrendering{\code{albedo}=0\%}{emitter_sky_albedo_0}
  \medrendering{\code{albedo}=100\%}{emitter_sky_albedo_1}
  \medrendering{\code{albedo}=20\% green}{emitter_sky_albedo_green}
  \caption{\label{fig:sky_groundalbedo}Influence
  of the ground albedo on the appearance of the sky}
}
\plugin{sun}{Sun emitter}
\icon{emitter_sun}
\order{7}
\parameters{
    \parameter{turbidity}{\Float}{
        This parameter determines the amount of aerosol present in the atmosphere.
        Valid range: 2-10. \default{3, corresponding to a clear sky in a temperate climate}
    }
    \parameter{year, month, day}{\Integer}{Denote the date of the
     observation \default{2010, 07, 10}}
    \parameter{hour,minute,\showbreak second}{\Float}{Local time
      at the location of the observer in 24-hour format\default{15, 00, 00,
      i.e. 3PM}}
    \parameter{latitude, longitude, timezone}{\Float}{
      These three parameters specify the oberver's latitude and longitude
      in degrees, and the local timezone offset in hours, which are required
      to compute the sun's position. \default{35.6894, 139.6917, 9 --- Tokyo, Japan}
    }
    \parameter{sunDirection}{\Vector}{Allows to manually
      override the sun direction in world space. When this value
      is provided, parameters pertaining to the computation
      of the sun direction (\code{year, hour, latitude,} etc.
      are unnecessary. \default{none}
    }
    \parameter{resolution}{\Integer}{Specifies the horizontal resolution of the precomputed
        image that is used to represent the sun environment map \default{512, i.e. 512$\times$256}}
    \parameter{scale}{\Float}{
        This parameter can be used to scale the amount of illumination
        emitted by the sun emitter. \default{1}
    }
    \parameter{sunRadiusScale}{\Float}{
        Scale factor to adjust the radius of the sun, while preserving its power.
        Set to \code{0} to turn it into a directional light source.
    }
    \parameter{samplingWeight}{\Float}{
        Specifies the relative amount of samples
        allocated to this emitter. \default{1}
    }
}
This plugin implements the physically-based sun model proposed by
Preetham et al. \cite{Preetham1999Practical}. Using the provided position
and time information (see \pluginref{sky} for details), it can determine the
position of the sun as seen from the position of the observer.
The radiance arriving at the earth surface is then found based on the spectral
emission profile of the sun and the extinction cross-section of the
atmosphere (which depends on the \code{turbidity} and the zenith angle of the sun).

Like the \code{blackbody} emission profile (Page~\pageref{sec:blackbody}),
the sun model introduces physical units into the rendering process.
The radiance values computed by this plugin have units of power ($W$) per
unit area ($m^{-2}$) per steradian ($sr^{-1}$) per unit wavelength ($nm^{-1}$).
If these units are inconsistent with your scene description, you may use the
optional \texttt{scale} parameter to adjust them.

This plugin supplies proper spectral power distributions when Mitsuba is
compiled in spectral rendering mode. Otherwise, they are simply projected onto
a linear RGB color space.

\remarks{
  \item The sun is an intense light source that subtends a tiny solid angle.
  This can be a problem for certain rendering techniques (e.g. path
  tracing), which produce high variance output (i.e. noise in renderings)
  when the scene also contains specular or glossy or materials.
}
\plugin{sunsky}{Sun and sky emitter}
\icon{emitter_sunsky}
\order{8}
\parameters{
    \parameter{turbidity}{\Float}{
        This parameter determines the amount of aerosol present in the atmosphere.
        Valid range: 1-10. \default{3, corresponding to a clear sky in a temperate climate}
    }
    \parameter{albedo}{\Spectrum}{Specifies the ground albedo \default{0.15}}
    \parameter{year, month, day}{\Integer}{Denote the date of the
     observation \default{2010, 07, 10}}
    \parameter{hour,minute,\showbreak second}{\Float}{Local time
      at the location of the observer in 24-hour format\default{15, 00, 00,
      i.e. 3PM}}
    \parameter{latitude, longitude, timezone}{\Float}{
      These three parameters specify the oberver's latitude and longitude
      in degrees, and the local timezone offset in hours, which are required
      to compute the sun's position. \default{35.6894, 139.6917, 9 --- Tokyo, Japan}
    }
    \parameter{sunDirection}{\Vector}{Allows to manually
      override the sun direction in world space. When this value
      is provided, parameters pertaining to the computation
      of the sun direction (\code{year, hour, latitude,} etc.
      are unnecessary. \default{none}
    }
    \parameter{stretch}{\Float}{
        Stretch factor to extend emitter below the horizon, must be
        in [1,2] \default{\code{1}, i.e. not used}
    }
    \parameter{resolution}{\Integer}{Specifies the horizontal resolution of the precomputed
        image that is used to represent the sun environment map \default{512, i.e. 512$\times$256}}
    \parameter{sunScale}{\Float}{
        This parameter can be used to separately scale the amount of illumination
        emitted by the sun. \default{1}
    }
    \parameter{skyScale}{\Float}{
        This parameter can be used to separately scale the amount of illumination
        emitted by the sky.\default{1}
    }
    \parameter{sunRadiusScale}{\Float}{
        Scale factor to adjust the radius of the sun, while preserving its power.
        Set to \code{0} to turn it into a directional light source.
    }
}
\vspace{-3mm}

\renderings{
  \medrendering{\pluginref{sky} emitter}{emitter_sunsky_sky}
  \medrendering{\pluginref{sun} emitter}{emitter_sunsky_sun}
  \medrendering{\pluginref{sunsky} emitter}{emitter_sunsky_sunsky}
  \vspace{-2mm}
  \caption{A coated rough copper test ball lit with the three
  provided daylight illumination models}
}
\vspace{1mm}
This convenience plugin has the sole purpose of instantiating
\pluginref{sun} and \pluginref{sky} and merging them into a joint
environment map. Please refer to these plugins individually for more
details.
\plugin{envmap}{Environment emitter}
\icon{emitter_envmap}
\order{9}
\parameters{
    \parameter{filename}{\String}{
      Filename of the radiance-valued input image to be loaded;
      must be in latitude-longitude format.
    }
    \parameter{scale}{\Float}{
        A scale factor that is applied to the
        radiance values stored in the input image. \default{1}
    }
    \parameter{toWorld}{\Transform}{
       Specifies an optional linear emitter-to-world space rotation.
       \default{none (i.e. emitter space $=$ world space)}
    }
    \parameter{gamma}{\Float}{
      Optional parameter to override the gamma value of the source bitmap,
      where 1 indicates a linear color space and the special value -1
      corresponds to sRGB. \default{automatically detect based on the
      image type and metadata}
    }
    \parameter{cache}{\Boolean}{
       Preserve generated MIP map data in a cache file? This will cause a file named
       \emph{filename}\code{.mip} to be created.
       \default{automatic---use caching for images larger than 1M pixels.}
    }
    \parameter{samplingWeight}{\Float}{
        Specifies the relative amount of samples
        allocated to this emitter. \default{1}
    }
}
\renderings{
  \rendering{The museum environment map by Bernhard Vogl that is used
  in many example renderings in this document}{emitter_envmap_example}
  \unframedrendering{Coordinate conventions used when mapping the input
  image onto the sphere.}{emitter_envmap_axes}
}

This plugin provides a HDRI (high dynamic range imaging) environment map,
which is a type of light source that is well-suited for representing ``natural''
illumination. Many images in this document are made using the environment map
shown in (a).

The implementation loads a captured illumination environment from a image in
latitude-longitude format and turns it into an infinitely distant emitter.
The image could either be a processed photograph or a rendering made using the
\pluginref{spherical} sensor. The direction conventions of this transformation
are shown in (b).
The plugin can work with all types of images that are natively supported by Mitsuba
(i.e. JPEG, PNG, OpenEXR, RGBE, TGA, and BMP). In practice, a good environment
map will contain high-dynamic range data that can only be represented using the
OpenEXR or RGBE file formats.
High quality free light probes are available on Paul Debevec's website
(\url{http://gl.ict.usc.edu/Data/HighResProbes}) and Bernhard Vogl's website
(\url{http://dativ.at/lightprobes/}).

Like the \pluginref{bitmap} texture, this plugin generates a cache file
named \emph{filename}\code{.mip} when given a large input image. This
significantly accelerates the loading times of subsequent renderings. When this
is not desired, specify \code{cache=false} to the plugin.
\plugin{constant}{Constant environment emitter}
\icon{emitter_constant}
\order{10}
\parameters{
    \parameter{radiance}{\Spectrum}{
        Specifies the emitted radiance in units of
        power per unit area per unit steradian.
    }
    \parameter{samplingWeight}{\Float}{
        Specifies the relative amount of samples
        allocated to this emitter. \default{1}
    }
}

This plugin implements a constant environment emitter, which surrounds
the scene and radiates diffuse illumination towards it. This is often
a good default light source when the goal is to visualize some loaded
geometry that uses basic (e.g. diffuse) materials.
\input{section_sensors}
\plugin{perspective}{Perspective pinhole camera}
\order{1}
\parameters{
    \parameter{toWorld}{\Transform\Or\Animation}{
       Specifies an optional camera-to-world transformation.
       \default{none (i.e. camera space $=$ world space)}
    }
    \parameter{focalLength}{\String}{
        Denotes the camera's focal length specified using
        \code{35mm} film equivalent units. See the main
        description for further details.
        \default{\code{50mm}}
    }
    \parameter{fov}{\Float}{
        An alternative to \code{focalLength}:
        denotes the camera's field of view in degrees---must be
        between 0 and 180, excluding the extremes.
    }
    \parameter{fovAxis}{\String}{
        When the parameter \code{fov} is given (and only then),
        this parameter further specifies the image axis, to
        which it applies.
        \begin{enumerate}[(i)]
            \item \code{\textbf{x}}: \code{fov} maps to the
                \code{x}-axis in screen space.
            \item \code{\textbf{y}}: \code{fov} maps to the
                \code{y}-axis in screen space.
            \item \code{\textbf{diagonal}}: \code{fov}
               maps to the screen diagonal.
            \item \code{\textbf{smaller}}: \code{fov}
               maps to the smaller dimension
               (e.g. \code{x} when \code{width}<\code{height})
            \item \code{\textbf{larger}}: \code{fov}
               maps to the larger dimension
               (e.g. \code{y} when \code{width}<\code{height})
        \end{enumerate}
        The default is \code{\textbf{x}}.
    }
    \parameter{shutterOpen, shutterClose}{\Float}{
        Specifies the time interval of the measurement---this
        is only relevant when the scene is in motion.
        \default{0}
    }
    \parameter{nearClip, farClip}{\Float}{
        Distance to the near/far clip
        planes.\default{\code{near\code}-\code{Clip=1e-2} (i.e.
        \code{0.01}) and {\code{farClip=1e4} (i.e. \code{10000})}}
    }
}
\renderings{
\rendering{The material test ball viewed through a perspective pinhole
camera. Everything is in sharp focus.}{sensor_perspective}
\medrendering{A rendering of the Cornell box}{sensor_perspective_2}
}

This plugin implements a simple idealizied perspective camera model, which
has an infinitely small aperture. This creates an infinite depth of field,
i.e. no optical blurring occurs. The camera is can be specified to move during
an exposure, hence temporal blur is still possible.

By default, the camera's field of view is specified using a 35mm film
equivalent focal length, which is first converted into a diagonal field
of view and subsequently applied to the camera. This assumes that
the film's aspect ratio matches that of 35mm film (1.5:1), though the
parameter still behaves intuitively when this is not the case.
Alternatively, it is also possible to specify a field of view in degrees
along a given axis (see the \code{fov} and \code{fovAxis} parameters).

The exact camera position and orientation is most easily expressed using the
\code{lookat} tag, i.e.:
\begin{xml}
<sensor type="perspective">
    <transform name="toWorld">
        <!-- Move and rotate the camera so that looks from (1, 1, 1) to (1, 2, 1)
             and the direction (0, 0, 1) points "up" in the output image -->
        <lookat origin="1, 1, 1" target="1, 2, 1" up="0, 0, 1"/>
    </transform>
</sensor>
\end{xml}
\plugin{thinlens}{Perspective camera with a thin lens}
\order{2}
\parameters{
    \parameter{toWorld}{\Transform\Or\Animation}{
       Specifies an optional camera-to-world transformation.
       \default{none (i.e. camera space $=$ world space)}
    }
    \parameter{apertureRadius}{\Float}{
        Denotes the radius of the camera's aperture in scene units.
    }
    \parameter{focusDistance}{\Float}{
        Denotes the world-space distance from the camera's aperture to the
        focal plane. \default{\code{0}}
    }
    \parameter{focalLength}{\String}{
        Denotes the camera's focal length specified using
        \code{35mm} film equivalent units. See the main description
        for further details.\default{\code{50mm}}
    }
    \parameter{fov}{\Float}{
        An alternative to \code{focalLength}:
        denotes the camera's field of view in degrees---must be
        between 0 and 180, excluding the extremes.
    }
    \parameter{fovAxis}{\String}{
        When the parameter \code{fov} is given (and only then),
        this parameter further specifies the image axis, to
        which it applies.
        \begin{enumerate}[(i)]
            \item \code{\textbf{x}}: \code{fov} maps to the
                \code{x}-axis in screen space.\vspace{-1mm}
            \item \code{\textbf{y}}: \code{fov} maps to the
                \code{y}-axis in screen space.\vspace{-1mm}
            \item \code{\textbf{diagonal}}: \code{fov}
               maps to the screen diagonal.\vspace{-1mm}
            \item \code{\textbf{smaller}}: \code{fov}
               maps to the smaller dimension
               (e.g. \code{x} when \code{width}<\code{height})\vspace{-1mm}
            \item \code{\textbf{larger}}: \code{fov}
               maps to the larger dimension
               (e.g. \code{y} when \code{width}<\code{height})
               \vspace{-1mm}
        \end{enumerate}
        The default is \code{\textbf{x}}.
    }
    \parameter{shutterOpen, shutterClose}{\Float}{
        Specifies the time interval of the measurement---this
        is only relevant when the scene is in motion.
        \default{0}
    }
    \parameter{nearClip, farClip}{\Float}{
        Distance to the near/far clip
        planes.\default{\code{near\code}-\code{Clip=1e-2} (i.e.
        \code{0.01}) and {\code{farClip=1e4} (i.e. \code{10000})}}
    }
}
\renderings{
\vspace{-4mm}
\rendering{The material test ball viewed through a perspective thin lens
camera. Points away from the focal plane project onto a circle of confusion.}{sensor_thinlens}
\medrendering{A rendering of the Cornell box}{sensor_thinlens_2}
}

This plugin implements a simple perspective camera model with a thin lens
at its circular aperture. It is very similar to the \pluginref{perspective} plugin
except that the extra lens element permits rendering with a
specifiable (i.e. non-infinite) depth of field. To configure this, it has two
extra parameters named \code{apertureRadius} and \code{focusDistance}.

By default, the camera's field of view is specified using a 35mm film
equivalent focal length, which is first converted into a diagonal field
of view and subsequently applied to the camera. This assumes that
the film's aspect ratio matches that of 35mm film (1.5:1), though the
parameter still behaves intuitively when this is not the case.
Alternatively, it is also possible to specify a field of view in degrees
along a given axis (see the \code{fov} and \code{fovAxis} parameters).

The exact camera position and orientation is most easily expressed using the
\code{lookat} tag, i.e.:
\begin{xml}
<sensor type="thinlens">
    <transform name="toWorld">
        <!-- Move and rotate the camera so that looks from (1, 1, 1) to (1, 2, 1)
             and the direction (0, 0, 1) points "up" in the output image -->
        <lookat origin="1, 1, 1" target="1, 2, 1" up="0, 0, 1"/>
    </transform>

    <!-- Focus on the target -->
    <float name="focusDistance" value="1"/>
    <float name="apertureRadius" value="0.1"/>
</sensor>
\end{xml}
\plugin{orthographic}{Orthographic camera}
\order{3}
\parameters{
    \parameter{toWorld}{\Transform\Or\Animation}{
       Specifies an optional camera-to-world transformation.
       \default{none (i.e. camera space $=$ world space)}
    }
    \parameter{shutterOpen, shutterClose}{\Float}{
        Specifies the time interval of the measurement---this
        is only relevant when the scene is in motion.
        \default{0}
    }
    \parameter{nearClip, farClip}{\Float}{
        Distance to the near/far clip
        planes.\default{\code{near\code}-\code{Clip=1e-2} (i.e.
        \code{0.01}) and {\code{farClip=1e4} (i.e. \code{10000})}}
    }
}
\renderings{
\rendering{The material test ball viewed through an orthographic camera.
Note the complete lack of perspective.}{sensor_orthographic}
\medrendering{A rendering of the Cornell box}{sensor_orthographic_2}
}

This plugin implements a simple orthographic camera, i.e. a sensor
based on an orthographic projection without any form of perspective.
It can be thought of as a planar sensor that measures the radiance
along its normal direction. By default, this is the region $[-1, 1]^2$ inside
the XY-plane facing along the positive Z direction. Transformed versions
can be instantiated e.g. as follows:

\begin{xml}
<sensor type="orthographic">
    <transform name="toWorld">
        <!-- Resize the sensor plane to 20x20 world space units -->
        <scale x="10" y="10"/>

        <!-- Move and rotate it so that it contains the point
             (1, 1, 1) and faces direction (0, 1, 0) -->
        <lookat origin="1, 1, 1" target="1, 2, 1" up="0, 0, 1"/>
    </transform>
</sensor>
\end{xml}
\plugin{telecentric}{Telecentric lens camera}
\order{4}
\parameters{
    \parameter{toWorld}{\Transform\Or\Animation}{
       Specifies an optional sensor-to-world transformation.
       \default{none (i.e. camera space $=$ world space)}
    }
    \parameter{apertureRadius}{\Float}{
        Denotes the radius of the camera's aperture in scene units.
        \default{\code{0}}
    }
    \parameter{focusDistance}{\Float}{
        Denotes the world-space distance from the camera's aperture to the
        focal plane. \default{\code{0}}
    }
    \parameter{shutterOpen, shutterClose}{\Float}{
        Specifies the time interval of the measurement---this
        is only relevant when the scene is in motion.
        \default{0}
    }
    \parameter{nearClip, farClip}{\Float}{
        Distance to the near/far clip
        planes.\default{\code{near\code}-\code{Clip=1e-2} (i.e.
        \code{0.01}) and {\code{farClip=1e4} (i.e. \code{10000})}}
    }
}
\renderings{
\rendering{The material test ball viewed through an telecentric camera.
Note the orthographic view together with a narrow depth of field.}{sensor_telecentric}
\medrendering{A rendering of the Cornell box. The red and green walls are partially visible
due to the aperture size.}{sensor_telecentric_2}
}

This plugin implements a simple model of a camera with a \emph{telecentric lens}.
This is a type of lens that produces an in-focus orthographic view on a
plane at some distance from the sensor. Points away from this plane are out of focus
and project onto a circle of confusion. In comparison to idealized orthographic cameras,
telecentric lens cameras exist in the real world and find use in some computer
vision applications where perspective effects cause problems.
This sensor relates to the \pluginref{orthographic} plugin in the same way
that \pluginref{thinlens} does to \pluginref{perspective}.

The configuration is identical to the \pluginref{orthographic} plugin, except that
the additional parameters \code{apertureRadius} and \code{focusDistance} must be provided.
\plugin{spherical}{Spherical camera}
\order{5}
\parameters{
    \parameter{toWorld}{\Transform\Or\Animation}{
       Specifies an optional camera-to-world transformation.
       \default{none (i.e. camera space $=$ world space)}
    }
    \parameter{shutterOpen, shutterClose}{\Float}{
        Specifies the time interval of the measurement---this
        is only relevant when the scene is in motion.
        \default{0}
    }
}

\renderings{
    \rendering{A rendering made using a spherical camera}{sensor_spherical_cbox.jpg}
}

The spherical camera captures the illumination arriving from all
directions and turns it into a latitude-longitude environment map.
It is best used with a high dynamic range film that has 2:1 aspect ratio,
and the resulting output can then be turned into a distant light source
using the \pluginref{envmap} plugin.
By default, the camera is located at the origin, which can
be changed by providing a custom \code{toWorld} transformation.
\plugin{irradiancemeter}{Irradiance meter}
\order{6}
\parameters{
    \parameter{shutterOpen, shutterClose}{\Float}{
        Specifies the time interval of the measurement---this
        is only relevant when the scene is in motion.
        \default{0}
    }
}

This sensor plugin implements a simple irradiance meter, which
measures the average incident power per unit area over a
provided surface.
Such a sensor is useful for conducting virtual experiments and
testing the renderer for correctness.
The result is normalized so that an irradiance sensor inside an
integrating sphere with constant radiance 1 records
an irradiance value of $\pi$.

To create an irradiance meter, instantiate the desired measurement
shape and specify the sensor as its child. Note that when the
sensor's film resolution is larger than $1\times 1$, each pixel
will record the average irradiance over a rectangular part of the
shape's UV parameterization.

\vspace{4mm}
\begin{xml}
<scene version=$\MtsVer$>
    <!-- Measure the average irradiance arriving on
         a unit radius sphere located at the origin -->
    <shape type="sphere">
        <sensor type="irradiancemeter">
            <!-- Write the output to a MATLAB M-file. The output file will
                 contain a 1x1 matrix storing an estimate of the average
                 irradiance over the surface of the sphere. -->
            <film type="mfilm"/>

            <!-- Use 1024 samples for the measurement -->
            <sampler type="independent">
                <integer name="sampleCount" value="1024"/>
            </sampler>
        </sensor>
    </shape>

    <!-- ... other scene declarations ... -->
</scene>
\end{xml}
\plugin{radiancemeter}{Radiance meter}
\order{6}
\parameters{
    \parameter{toWorld}{\Transform\Or\Animation}{
       Specifies an optional sensor-to-world transformation.
       \default{none (i.e. sensor space $=$ world space)}
    }
    \parameter{shutterOpen, shutterClose}{\Float}{
        Specifies the time interval of the measurement---this
        is only relevant when the scene is in motion.
        \default{0}
    }
}

This sensor plugin implements a simple radiance meter, which measures
the incident power per unit area per unit solid angle along a
certain ray. It can be thought of as the limit of a standard
perspective camera as its field of view tends to zero.
Hence, when this sensor is given a film with multiple pixels, all
of them will record the same value.

Such a sensor is useful for conducting virtual experiments and
testing the renderer for correctness.

By default, the sensor is located at the origin and performs
a measurement in the positive Z direction $(0,0,1)$. This can
be changed by providing a custom \code{toWorld} transformation:

\vspace{4mm}
\begin{xml}
<scene version=$\MtsVer$>
    <sensor type="radiancemeter">
        <!-- Measure the amount of radiance traveling
             from the origin to (1,2,3) -->
        <transform name="toWorld">
            <lookat origin="1,2,3"
                    target="0,0,0"/>
        </transform>

        <!-- Write the output to a MATLAB M-file. The output file will
             contain a 1x1 matrix storing an estimate of the incident
             radiance along the specified ray. -->
        <film type="mfilm"/>

        <!-- Use 1024 samples for the measurement -->
        <sampler type="independent">
            <integer name="sampleCount" value="1024"/>
        </sampler>
    </sensor>

    <!-- ... other scene declarations ... -->
</scene>
\end{xml}
\plugin{fluencemeter}{Fluence meter}
\order{7}
\parameters{
    \parameter{toWorld}{\Transform\Or\Animation}{
       Specifies an optional sensor-to-world transformation.
       \default{none (i.e. sensor space $=$ world space)}
    }
    \parameter{shutterOpen, shutterClose}{\Float}{
        Specifies the time interval of the measurement---this
        is only relevant when the scene is in motion.
        \default{0}
    }
}

This sensor plugin implements a simple fluence meter, which measures
the average radiance passing through a specified position.
By default, the sensor is located at the origin.

Such a sensor is useful for conducting virtual experiments and
testing the renderer for correctness.

\vspace{4mm}
\begin{xml}
<scene version=$\MtsVer$>
    <sensor type="fluencemeter">
        <!-- Measure the average radiance traveling
             through the point (1,2,3) -->
        <transform name="toWorld">
            <translate x="1" y="2" z="3"/>
        </transform>

        <!-- Write the output to a MATLAB M-file. The output file will
             contain a 1x1 matrix storing the computed estimate -->
        <film type="mfilm"/>

        <!-- Use 1024 samples for the measurement -->
        <sampler type="independent">
            <integer name="sampleCount" value="1024"/>
        </sampler>
    </sensor>

    <!-- ... other scene declarations ... -->
</scene>
\end{xml}
\plugin[perspectiverdist]{perspective\_rdist}{Perspective pinhole camera with radial distortion}
\order{8}
\parameters{
    \parameter{toWorld}{\Transform\Or\Animation}{
       Specifies an optional camera-to-world transformation.
       \default{none (i.e. camera space $=$ world space)}
    }
    \parameter{kc}{\String}{
        Second and fourth-order coefficient of a polynomial radial distortion model
        specified as a comma-separated list
        The specifics of the model are described in detail on the following page:
        \url{http://www.vision.caltech.edu/bouguetj/calib_doc/htmls/parameters.html}
    }
    \parameter{focalLength}{\String}{
        Denotes the camera's focal length specified using
        \code{35mm} film equivalent units. See the main
        description for further details.
        \default{\code{50mm}}
    }
    \parameter{fov}{\Float}{
        An alternative to \code{focalLength}:
        denotes the camera's field of view in degrees---must be
        between 0 and 180, excluding the extremes.
    }
    \parameter{fovAxis}{\String}{
        When the parameter \code{fov} is given (and only then),
        this parameter further specifies the image axis, to
        which it applies.
        \begin{enumerate}[(i)]
            \item \code{\textbf{x}}: \code{fov} maps to the
                \code{x}-axis in screen space.
            \item \code{\textbf{y}}: \code{fov} maps to the
                \code{y}-axis in screen space.
            \item \code{\textbf{diagonal}}: \code{fov}
               maps to the screen diagonal.
            \item \code{\textbf{smaller}}: \code{fov}
               maps to the smaller dimension
               (e.g. \code{x} when \code{width}<\code{height})
            \item \code{\textbf{larger}}: \code{fov}
               maps to the larger dimension
               (e.g. \code{y} when \code{width}<\code{height})
        \end{enumerate}
        The default is \code{\textbf{x}}.
    }
    \parameter{shutterOpen, shutterClose}{\Float}{
        Specifies the time interval of the measurement---this
        is only relevant when the scene is in motion.
        \default{0}
    }
    \parameter{nearClip, farClip}{\Float}{
        Distance to the near/far clip
        planes.\default{\code{near\code}-\code{Clip=1e-2} (i.e.
        \code{0.01}) and {\code{farClip=1e4} (i.e. \code{10000})}}
    }
}

This plugin extends the \pluginref{perspective} camera with support
for radial distortion. It accepts an additional parameter named \code{kc},
which specifies the second and fourth-order terms in a polynomial
model that accounts for pincushion and barrel distortion. This is
useful when trying to match rendered images to photographs created by a
camera whose distortion is known. When \code{kc=0, 0}, the model
turns into a standard pinhole camera. The reason for creating a separate
plugin for this feature is that distortion involves extra overheads per ray that
users may not be willing to pay for if their scene doesn't use it.
The MATLAB Camera Calibration Toolbox by Jean-Yves Bouguet
(\url{http://www.vision.caltech.edu/bouguetj/calib_doc/}) can be used to
obtain a distortion model, and the first entries of the \code{kc} variable
generated by this tool can directly be passed into this plugin.
\input{section_integrators}
\plugin{ao}{Ambient occlusion integrator}
\order{0}
\parameters{
    \parameter{shadingSamples}{\Integer}{Specifies the number of
        shading samples that should be computed per primary ray
        \default{1}}

    \parameter{rayLength}{\Float}{Specifies the world-space length of the
        ambient occlusion rays that will be cast. \default{\code{-1}, i.e. automatic}}.
}
\renderings{
   \rendering{A view of the scene on page \pageref{fig:rungholt}, rendered using
              the Ambient Occlusion integrator}{integrator_ao}
   \rendering{A corresponding rendering created using the standard \pluginref{path} tracer}{integrator_ao_path}
}
Ambient Occlusion is a simple non-photorealistic rendering technique that simulates the exposure of an object
to uniform illumination incident from all direction. It produces approximate shadowing between closeby
objects, as well as darkening in corners, creases, and cracks. The scattering models associated with objects
in the scene are ignored.
\plugin{direct}{Direct illumination integrator}
\order{1}
\parameters{
    \parameter{shadingSamples}{\Integer}{This convenience parameter can be
        used to set both \code{emitterSamples} and \code{bsdfSamples} at
        the same time.
    }
    \parameter{emitterSamples}{\Integer}{Optional more fine-grained
       parameter: specifies the number of samples that should be generated
       using the direct illumination strategies implemented by the scene's
       emitters\default{set to the value of \code{shadingSamples}}
    }
    \parameter{bsdfSamples}{\Integer}{Optional more fine-grained
       parameter: specifies the number of samples that should be generated
       using the BSDF sampling strategies implemented by the scene's
       surfaces\default{set to the value of \code{shadingSamples}}
    }
    \parameter{strictNormals}{\Boolean}{Be strict about potential
       inconsistencies involving shading normals? See
       page~\pageref{sec:strictnormals} for details.
       \default{no, i.e. \code{false}}
    }
    \parameter{hideEmitters}{\Boolean}{Hide directly visible emitters?
       See page~\pageref{sec:hideemitters} for details.
       \default{no, i.e. \code{false}}
    }
}
\vspace{-1mm}
\renderings{
    \medrendering{Only BSDF sampling}{integrator_direct_bsdf}
    \medrendering{Only emitter sampling}{integrator_direct_lum}
    \medrendering{BSDF and emitter sampling}{integrator_direct_both}
    \caption{
        \label{fig:integrator-direct}
        This plugin implements two different strategies for computing the
        direct illumination on surfaces. Both of them are dynamically
        combined then obtain a robust rendering algorithm.
    }
}

This integrator implements a direct illumination technique that makes use
of \emph{multiple importance sampling}: for each pixel sample, the
integrator generates a user-specifiable number of BSDF and emitter
samples and combines them using the power heuristic. Usually, the BSDF
sampling technique works very well on glossy objects but does badly
everywhere else (\subfigref{integrator-direct}{a}), while the opposite
is true for the emitter sampling technique
(\subfigref{integrator-direct}{b}). By combining these approaches, one
can obtain a rendering technique that works well in both cases
(\subfigref{integrator-direct}{c}).

The number of samples spent on either technique is configurable, hence
it is also possible to turn this plugin into an emitter sampling-only
or BSDF sampling-only integrator.

For best results, combine the direct illumination integrator with the
low-discrepancy sample generator (\code{ldsampler}). Generally, the number
of pixel samples of the sample generator can be kept relatively
low (e.g. \code{sampleCount=4}), whereas the \code{shadingSamples}
parameter of this integrator should be increased until the variance in
the output renderings is acceptable.

\remarks{
   \item This integrator does not handle participating media or
         indirect illumination.
}
\plugin{path}{Path tracer}
\order{2}
\parameters{
    \parameter{maxDepth}{\Integer}{Specifies the longest path depth
        in the generated output image (where \code{-1} corresponds to $\infty$).
        A value of \code{1} will only render directly visible light sources.
        \code{2} will lead to single-bounce (direct-only) illumination,
        and so on. \default{\code{-1}}
    }
    \parameter{rrDepth}{\Integer}{Specifies the minimum path depth, after
       which the implementation will start to use the ``russian roulette''
       path termination criterion. \default{\code{5}}
    }
    \parameter{strictNormals}{\Boolean}{Be strict about potential
       inconsistencies involving shading normals? See the description below
       for details.\default{no, i.e. \code{false}}
    }
    \parameter{hideEmitters}{\Boolean}{Hide directly visible emitters?
       See page~\pageref{sec:hideemitters} for details.
       \default{no, i.e. \code{false}}
    }
}

This integrator implements a basic path tracer and is a \emph{good default choice}
when there is no strong reason to prefer another method.

To use the path tracer appropriately, it is instructive to know roughly how
it works: its main operation is to trace many light paths using \emph{random walks}
starting from the sensor. A single random walk is shown below, which entails
casting a ray associated with a pixel in the output image and searching for
the first visible intersection. A new direction is then chosen at the intersection,
and the ray-casting step repeats over and over again (until one of several
stopping criteria applies).
\begin{center}
\includegraphics[width=.7\textwidth]{images/integrator_path_figure.pdf}
\end{center}
At every intersection, the path tracer tries to create a connection to
the light source in an attempt to find a \emph{complete} path along which
light can flow from the emitter to the sensor. This of course only works
when there is no occluding object between the intersection and the emitter.

This directly translates into a category of scenes where
a path tracer can be expected to produce reasonable results: this is the case
when the emitters are easily ``accessible'' by the contents of the scene. For instance,
an interior scene that is lit by an area light will be considerably harder
to render when this area light is inside a glass enclosure (which
effectively counts as an occluder).

Like the \pluginref{direct} plugin, the path tracer internally relies on multiple importance
sampling to combine BSDF and emitter samples. The main difference in comparison
to the former plugin is that it considers light paths of arbitrary length to compute
both direct and indirect illumination.

For good results, combine the path tracer with one of the
low-discrepancy sample generators (i.e. \pluginref{ldsampler},
\pluginref{halton}, or \pluginref{sobol}).

\paragraph{Strict normals:}\label{sec:strictnormals}
Triangle meshes often rely on interpolated shading normals
to suppress the inherently faceted appearance of the underlying geometry. These
``fake'' normals are not without problems, however. They can lead to paradoxical
situations where a light ray impinges on an object from a direction that is classified as ``outside''
according to the shading normal, and ``inside'' according to the true geometric normal.

The \code{strictNormals}
parameter specifies the intended behavior when such cases arise. The default (\code{false}, i.e. ``carry on'')
gives precedence to information given by the shading normal and considers such light paths to be valid.
This can theoretically cause light ``leaks'' through boundaries, but it is not much of a problem in practice.

When set to \code{true}, the path tracer detects inconsistencies and ignores these paths. When objects
are poorly tesselated, this latter option may cause them to lose a significant amount of the incident
radiation (or, in other words, they will look dark).

The bidirectional integrators in Mitsuba (\pluginref{bdpt}, \pluginref{pssmlt}, \pluginref{mlt} ...)
implicitly have \code{strictNormals} set to \code{true}. Hence, another use of this parameter
is to match renderings created by these methods.

\remarks{
   \item This integrator does not handle participating media
   \item This integrator has poor convergence properties when rendering
   caustics and similar effects. In this case, \pluginref{bdpt} or
   one of the photon mappers may be preferable.
}
\plugin[volpathsimple]{volpath\_simple}{Simple volumetric path tracer}
\order{3}
\parameters{
    \parameter{maxDepth}{\Integer}{Specifies the longest path depth
        in the generated output image (where \code{-1} corresponds to $\infty$).
        A value of \code{1} will only render directly visible light sources.
        \code{2} will lead to single-bounce (direct-only) illumination,
        and so on. \default{\code{-1}}
    }
    \parameter{rrDepth}{\Integer}{Specifies the minimum path depth, after
       which the implementation will start to use the ``russian roulette''
       path termination criterion. \default{\code{5}}
    }
    \parameter{strictNormals}{\Boolean}{Be strict about potential
       inconsistencies involving shading normals? See
       page~\pageref{sec:strictnormals} for details.
       \default{no, i.e. \code{false}}
    }
    \parameter{hideEmitters}{\Boolean}{Hide directly visible emitters?
       See page~\pageref{sec:hideemitters} for details.
       \default{no, i.e. \code{false}}
    }
}

This plugin provides a basic volumetric path tracer that can be used to
compute approximate solutions of the radiative transfer equation. This
particular integrator is named ``simple'' because it does not make use of
multiple importance sampling. This results in a potentially
faster execution time. On the other hand, it also means that this
plugin will likely not perform well when given a scene that contains
highly glossy materials. In this case, please use \pluginref{volpath}
or one of the bidirectional techniques.

This integrator has special support for \emph{index-matched} transmission
events (i.e. surface scattering events that do not change the direction
of light). As a consequence, participating media enclosed by a stencil shape (see
\secref{shapes} for details) are rendered considerably more efficiently when this
shape has \emph{no}\footnote{this is what signals to Mitsuba that the boundary is
index-matched and does not interact with light in any way. Alternatively,
the \pluginref{mask} and \pluginref{thindielectric} BSDF can be used to specify
index-matched boundaries that involve some amount of interaction.} BSDF assigned
to it (as compared to, say, a \pluginref{dielectric} or \pluginref{roughdielectric} BSDF).

\remarks{
   \item This integrator performs poorly when rendering
     participating media that have a different index of refraction compared
     to the surrounding medium.
   \item This integrator has difficulties rendering
     scenes that contain relatively glossy materials (\pluginref{volpath} is preferable in this case).
   \item This integrator has poor convergence properties when rendering
   caustics and similar effects. In this case, \pluginref{bdpt} or
   one of the photon mappers may be preferable.
}
\plugin{volpath}{Extended volumetric path tracer}
\order{4}
\parameters{
    \parameter{maxDepth}{\Integer}{Specifies the longest path depth
        in the generated output image (where \code{-1} corresponds to $\infty$).
        A value of \code{1} will only render directly visible light sources.
        \code{2} will lead to single-bounce (direct-only) illumination,
        and so on. \default{\code{-1}}
    }
    \parameter{rrDepth}{\Integer}{Specifies the minimum path depth, after
       which the implementation will start to use the ``russian roulette''
       path termination criterion. \default{\code{5}}
    }
    \parameter{strictNormals}{\Boolean}{Be strict about potential
       inconsistencies involving shading normals? See
       page~\pageref{sec:strictnormals} for details.
       \default{no, i.e. \code{false}}
    }
    \parameter{hideEmitters}{\Boolean}{Hide directly visible emitters?
       See page~\pageref{sec:hideemitters} for details.
       \default{no, i.e. \code{false}}
    }
}

This plugin provides a volumetric path tracer that can be used to
compute approximate solutions of the radiative transfer equation.
Its implementation makes use of multiple importance sampling to
combine BSDF and phase function sampling with direct illumination
sampling strategies. On surfaces, it behaves exactly
like the standard path tracer.

This integrator has special support for \emph{index-matched} transmission
events (i.e. surface scattering events that do not change the direction
of light). As a consequence, participating media enclosed by a stencil shape (see
\secref{shapes} for details) are rendered considerably more efficiently when this
shape has \emph{no}\footnote{this is what signals to Mitsuba that the boundary is
index-matched and does not interact with light in any way. Alternatively,
the \pluginref{mask} and \pluginref{thindielectric} BSDF can be used to specify
index-matched boundaries that involve some amount of interaction.} BSDF assigned
to it (as compared to, say, a \pluginref{dielectric} or \pluginref{roughdielectric} BSDF).

\remarks{
   \item This integrator will generally perform poorly when rendering
     participating media that have a different index of refraction compared
     to the surrounding medium.
   \item This integrator has poor convergence properties when rendering
     caustics and similar effects. In this case, \pluginref{bdpt} or
     one of the photon mappers may be preferable.
}
\plugin{bdpt}{Bidirectional path tracer}
\order{5}
\parameters{
    \parameter{maxDepth}{\Integer}{Specifies the longest path depth
        in the generated output image (where \code{-1} corresponds to $\infty$).
        A value of \code{1} will only render directly visible light sources.
        \code{2} will lead to single-bounce (direct-only) illumination,
        and so on. \default{\code{-1}}
    }
    \parameter{lightImage}{\Boolean}{Include sampling strategies that connect
       paths traced from emitters directly to the camera? (i.e. what \pluginref{ptracer} does)
       This improves the effectiveness of bidirectional path tracing
       but severely increases the local and remote communication
       overhead, since large \emph{light images} must be transferred between threads
       or over the network. See the text below for a more detailed explanation.
       \default{include these strategies, i.e. \code{true}}
    }
    \parameter{sampleDirect}{\Boolean}{Enable direct sampling strategies? This is a generalization
       of direct illumination sampling that works with both emitters and sensors. Usually a good idea.
       \default{use direct sampling, i.e. \code{true}}}
    \parameter{rrDepth}{\Integer}{Specifies the minimum path depth, after
       which the implementation will start to use the ``russian roulette''
       path termination criterion. \default{\code{5}}
    }
}

\renderings{
    \rendering{Path tracer, 32 samples/pixel}{integrator_bdpt_path}
    \rendering{Bidirectional path tracer, 32 samples/pixel}{integrator_bdpt_bdpt}
    \caption{
     The bidirectional path tracer finds light paths by generating partial
     paths starting at the emitters and the sensor and connecting them in every possible way.
     This works particularly well in closed scenes as the one shown above. Here, the unidirectional
     path tracer has severe difficulties finding some of the indirect illumination paths.
     Modeled after a scene by Eric Veach.
    }
}
\renderings{
   \includegraphics[width=12cm]{images/integrator_bdpt_sketch.pdf}\hfill\,
   \caption{The four different ways in which BDPT can create a direct illumination
   path (matching the first row on the next page): \textbf{(a)} Standard path
   tracing without direct illumination sampling, \textbf{(b)} path tracing with
   direct illumination sampling, \textbf{(c)} Particle tracing with recording of
   scattering events observed by the sensor, \textbf{(d)} Particle tracing with
   recording of particles that hit the sensor.}\vspace{-3mm}
}
\renderings{
    \unframedbigrendering{The individual sampling strategies that comprise BDPT, but
    \emph{without} multiple importance sampling. $s$ denotes the number of steps
    taken from the emitters, and $t$ denotes the number of steps from the sensor.
    Note how almost every strategy has deficiencies of some kind} {integrator_bdpt_unweighted.pdf}
    \unframedbigrendering{The same sampling strategies, but now weighted using multiple importance
    sampling---effectively ``turning off'' each strategy where it does not perform well.
    The final result is computed by summing all of these images.}{integrator_bdpt_weighted.pdf}
}

This plugin implements a bidirectional path tracer (short: BDPT) with support for multiple
importance sampling, as proposed by Veach and Guibas \cite{Veach1994Bidirectional}.

A bidirectional path tracer computes radiance estimates by starting two separate
random walks from an emitter and a sensor. The resulting \emph{subpaths} are
connected at every possible interaction vertex, creating a large number of complete paths
of different lengths. These paths are then used to estimate the amount of
radiance that is transferred from the emitter to a pixel on the sensor.

Generally, some of the created paths will be undesirable, since they lead to
high-variance radiance estimates. To alleviate this situation, BDPT makes
use of \emph{multiple importance sampling} which, roughly speaking, weights
paths based on their predicted utility.

The bidirectional path tracer in Mitsuba is a complete implementation of the
technique that handles all sampling strategies, including those that involve
direct interactions with the sensor. For this purpose, finite-aperture sensors
are explicitly represented by surfaces in the scene so that they can be
intersected by random walks started at emitters.

Bidirectional path tracing is a relatively ``heavy'' rendering technique---for
the same number of samples per pixel, it is easily 3-4 times slower than regular
path tracing. However, it usually makes up for this by producing considerably
lower-variance radiance estimates (i.e. the output images have less noise).

The code parallelizes over multiple cores and machines, but with one caveat:
some of the BDPT path sampling strategies are incompatble with the usual
approach of rendering an image tile by tile, since they can potentially
contribute to \emph{any} pixel on the screen. This means that each
rendering work unit must be associated with a full-sized image!
When network render nodes are involved or the resolution of this \emph{light image}
is very high, a bottleneck can arise where more work is spent  accumulating or
transmitting these images than actual rendering.

There are two possible resorts should this situation arise: the first one
is to reduce the number of work units so that there is approximately one
unit per core (and hence one image to transmit per core). This can be done by
increasing the block size in the GUI preferences or passing the \code{-b}
parameter to the \code{mitsuba} executable. The second option is to simply
disable these sampling strategies at the cost of reducing the
effectiveness of bidirectional path tracing (particularly, when rendering
caustics). For this, set \code{lightImage} to \code{false}.
When rendering an image of a reasonable resolution without network nodes,
this is not a big concern, hence these strategies are enabled by default.

\remarks{
   \item This integrator does not work with dipole-style subsurface
   scattering models.
   \item This integrator does not yet work with certain non-reciprocal BSDFs (i.e.
         \pluginref{bumpmap}, but this will be addressed in the future
}
\plugin{photonmapper}{Photon map integrator}
\order{6}
\parameters{
    \parameter{directSamples}{\Integer}{Number of samples used for the
       direct illumination component \default{16}}
    \parameter{glossySamples}{\Integer}{Number of samples used for the indirect
       illumination component of glossy materials \default{32}}
    \parameter{maxDepth}{\Integer}{Specifies the longest path depth
        in the generated output image (where \code{-1} corresponds to $\infty$).
        A value of \code{1} will only render directly visible light sources.
        \code{2} will lead to single-bounce (direct-only) illumination,
        and so on. \default{\code{-1}}
    }
    \parameter{globalPhotons}{\Integer}{Number of photons that will be collected for the global photon map\default{250000}}
    \parameter{causticPhotons}{\Integer}{Number of photons that will be collected for the caustic photon map\default{250000}}
    \parameter{volumePhotons}{\Integer}{Number of photons that will be collected for the volumetric photon map\default{250000}}
    \parameter{globalLookup\showbreak Radius}{\Float}{Maximum radius of photon lookups in the global photon map (relative to the scene size)\default{0.05}}
    \parameter{causticLookup\showbreak Radius}{\Float}{Maximum radius of photon lookups in the caustic photon map (relative to the scene size)\default{0.0125}}
    \parameter{lookupSize}{\Integer}{Number of photons that should be fetched in photon map queries\default{120}}
    \parameter{granularity}{\Integer}{
     Granularity of photon tracing work units for the purpose
     of parallelization (in \# of shot particles) \default{0, i.e. decide automatically}
    }
    \parameter{hideEmitters}{\Boolean}{Hide directly visible emitters?
       See page~\pageref{sec:hideemitters} for details.
       \default{no, i.e. \code{false}}
    }
    \parameter{rrDepth}{\Integer}{Specifies the minimum path depth, after
       which the implementation will start to use the ``russian roulette''
       path termination criterion. \default{\code{5}}
    }
}
This plugin implements the two-pass photon mapping algorithm as proposed by Jensen \cite{Jensen1996Global}.
The implementation partitions the illumination into three different classes (diffuse, caustic, and volumetric),
and builds a separate photon map for each class.

Following this, a standard recursive ray tracing pass is started which performs kernel density estimation
using these photon maps. Since the photon maps are visualized directly, the result will appear ``blotchy'' (\figref{pmap-blotchy})
unless an extremely large number of photons is used. A simple remedy is to combine the photon mapper with
an irradiance cache, which performs \emph{final gathering} to remove these artifacts. Due to its caching nature,
the rendering process will be faster as well.
\begin{xml}[caption={Instantiation of a photon mapper with irradiance caching}]
<integrator type="irrcache">
    <integrator type="photonmapper"/>
</integrator>
\end{xml}
\renderings{
  \rendering{Rendered using plain photon mapping}{integrator_photonmapper_1}
  \rendering{Rendered using photon mapping together with irradiance caching}{integrator_photonmapper_2}
  \vspace{-2mm}
  \caption{\label{fig:pmap-blotchy}Sponza atrium illuminated by a point light and rendered using 5 million photons.
  Irradiance caching significantly accelerates the rendering time and eliminates the ``blotchy''
  kernel density estimation artifacts. Model courtesy of Marko Dabrovic.}
}

When the scene contains participating media, the Beam Radiance Estimate \cite{Jarosz2008Beam}
by Jarosz et al. is used to estimate the illumination due to volumetric scattering.

\remarks{
    \item Currently, only homogeneous participating media are supported by this implementation
}
\plugin{ppm}{Progressive photon mapping integrator}
\order{7}
\parameters{
    \parameter{maxDepth}{\Integer}{Specifies the longest path depth
        in the generated output image (where \code{-1} corresponds to $\infty$).
        A value of \code{1} will only render directly visible light sources.
        \code{2} will lead to single-bounce (direct-only) illumination,
        and so on. \default{\code{-1}}
    }
    \parameter{photonCount}{\Integer}{Number of photons to be shot per iteration\default{250000}}
    \parameter{initialRadius}{\Float}{Initial radius of gather points in world space units.
        \default{0, i.e. decide automatically}}
    \parameter{alpha}{\Float}{Radius reduction parameter \code{alpha} from the paper\default{0.7}}
    \parameter{granularity}{\Integer}{
Granularity of photon tracing work units for the purpose
of parallelization (in \# of shot particles) \default{0, i.e. decide automatically}
    }
    \parameter{rrDepth}{\Integer}{Specifies the minimum path depth, after
       which the implementation will start to use the ``russian roulette''
       path termination criterion. \default{\code{5}}
    }
    \parameter{maxPasses}{\Integer}{Maximum number of passes to render (where \code{-1}
       corresponds to rendering until stopped manually). \default{\code{-1}}}
}
This plugin implements the progressive photon mapping algorithm by Hachisuka et al.
\cite{Hachisuka2008Progressive}. Progressive photon mapping is a variant of photon
mapping that alternates between photon shooting and gathering passes that involve
a relatively small (e.g. 250K) numbers of photons that are subsequently discarded.

This is done in a way such that the variance and bias of the resulting output
vanish as the number of passes tends to infinity. The progressive nature of this
method enables renderings with an effectively arbitrary number of photons
without exhausting the available system memory.

The desired sample count specified in the sample generator configuration
determines how many photon query points are created per pixel. It should not be
set too high, since the rendering time is approximately proportional to
this number. For good results, use between 2-4 samples along with the
\code{ldsampler}. Once started, the rendering process continues indefinitely
until it is manually stopped.

\remarks{
   \item Due to the data dependencies of this algorithm, the parallelization is
   limited to the local machine (i.e. cluster-wide renderings are not implemented)
   \item This integrator does not handle participating media
   \item This integrator does not currently work with subsurface scattering
   models.
}
\plugin{sppm}{Stochastic progressive photon mapping integrator}
\order{8}
\parameters{
    \parameter{maxDepth}{\Integer}{Specifies the longest path depth
        in the generated output image (where \code{-1} corresponds to $\infty$).
        A value of \code{1} will only render directly visible light sources.
        \code{2} will lead to single-bounce (direct-only) illumination,
        and so on. \default{\code{-1}}
    }
    \parameter{photonCount}{\Integer}{Number of photons to be shot per iteration\default{250000}}
    \parameter{initialRadius}{\Float}{Initial radius of gather points in world space units.
        \default{0, i.e. decide automatically}}
    \parameter{alpha}{\Float}{Radius reduction parameter \code{alpha} from the paper\default{0.7}}
    \parameter{granularity}{\Integer}{
Granularity of photon tracing work units for the purpose
of parallelization (in \# of shot particles) \default{0, i.e. decide automatically}
    }
    \parameter{rrDepth}{\Integer}{Specifies the minimum path depth, after
       which the implementation will start to use the ``russian roulette''
       path termination criterion. \default{\code{5}}
    }
    \parameter{maxPasses}{\Integer}{Maximum number of passes to render (where \code{-1}
       corresponds to rendering until stopped manually). \default{\code{-1}}}
}
This plugin implements stochastic progressive photon mapping by Hachisuka et al.
\cite{Hachisuka2009Stochastic}. This algorithm is an extension of progressive photon
mapping (\pluginref{ppm}) that improves convergence
when rendering scenes involving depth-of-field, motion blur, and glossy reflections.

Note that the implementation of \pluginref{sppm} in Mitsuba ignores the sampler
configuration---hence, the usual steps of choosing a sample generator and a desired
number of samples per pixel are not necessary. As with \pluginref{ppm}, once started,
the rendering process continues indefinitely until it is manually stopped.

\remarks{
   \item Due to the data dependencies of this algorithm, the parallelization is
   limited to the local machine (i.e. cluster-wide renderings are not implemented)
   \item This integrator does not handle participating media
   \item This integrator does not currently work with subsurface scattering
   models.
}
\plugin{pssmlt}{Primary Sample Space Metropolis Light Transport}
\order{9}
\parameters{
    \parameter{bidirectional}{\Boolean}{
    PSSMLT works in conjunction with another rendering
    technique that is endowed with Markov Chain-based sample generation.
    Two choices are available (Default: \code{true}):
     \begin{itemize}
     \item \code{true}: Operate on top of a fully-fleged bidirectional
       path tracer with multiple importance sampling.
     \item \code{false}: Rely on a unidirectional
     volumetric path tracer (i.e. \pluginref{volpath})
     \vspace{-4mm}
     \end{itemize}
    }
    \parameter{maxDepth}{\Integer}{Specifies the longest path depth
        in the generated output image (where \code{-1} corresponds to $\infty$).
        A value of \code{1} will only render directly visible light sources.
        \code{2} will lead to single-bounce (direct-only) illumination,
        and so on. \default{\code{-1}}
    }
    \parameter{directSamples}{\Integer}{
        By default, this plugin renders the direct illumination component
        separately using an optimized direct illumination sampling strategy
        that uses low-discrepancy number sequences for superior performance
        (in other words, it is \emph{not} rendered by PSSMLT). This
        parameter specifies the number of samples allocated to that method. To
        force PSSMLT to be responsible for the direct illumination
        component as well, set this parameter to \code{-1}. \default{16}
    }
    \parameter{rrDepth}{\Integer}{Specifies the minimum path depth, after
       which the implementation will start to use the ``russian roulette''
       path termination criterion. \default{\code{5}}
    }
    \parameter{luminanceSamples}{\Integer}{
       MLT-type algorithms create output images that are only
       \emph{relative}. The algorithm can e.g. determine that a certain pixel
       is approximately twice as bright as another one, but the absolute
       scale is unknown. To recover it, this plugin computes
       the average luminance arriving at the sensor by generating a
       number of samples. \default{\code{100000} samples}
    }
    \parameter{twoStage}{\Boolean}{Use two-stage MLT?
      See below for details. \default{{\footnotesize\code{false}}}}
    \parameter{pLarge}{\Float}{
      Rate at which the implementation tries to replace the current path
      with a completely new one. Usually, there is little need to change
      this. \default{0.3}
    }
}
Primary Sample Space Metropolis Light Transport (PSSMLT) is a rendering
technique developed by Kelemen et al. \cite{Kelemen2002Simple} which is
based on Markov Chain Monte Carlo (MCMC) integration.
\renderings{
   \vspace{-2mm}
   \includegraphics[width=11cm]{images/integrator_pssmlt_sketch.pdf}\hfill\,
   \vspace{-3mm}
   \caption{PSSMLT piggybacks on a rendering method that can turn points
   in the primary sample space (i.e. ``random numbers'') into paths. By
   performing small jumps in primary sample space, it can explore the neighborhood
   of a path\vspace{-5mm}}
}
In contrast to simple methods like path tracing that render
images by performing a na\"ive and memoryless random search for light paths,
PSSMLT actively searches for \emph{relevant} light paths (as is the case
for other MCMC methods). Once such a path is found, the algorithm tries to
explore neighboring paths to amortize the cost of the search. This can
significantly improve the convergence rate of difficult input.
Scenes that were already relatively easy to render usually don't benefit
much from PSSMLT, since the MCMC data management causes additional
computational overheads.

An interesting aspect of PSSMLT is that it performs this exploration
of light paths by perturbing the ``random numbers'' that were initially
used to construct the path. Subsequent regeneration of the path using the
perturbed numbers yields a new path in a slightly different configuration, and
this process repeats over and over again.
The path regeneration step is fairly general and this is what makes
the method powerful: in particular, it is possible to use PSSMLT as a
layer on top of an existing method to create a new ``metropolized''
version of the rendering algorithm that is enhanced with a certain
degree of adaptiveness as described earlier.

The PSSMLT implementation in Mitsuba can operate on top of either a simple
unidirectional volumetric path tracer or a fully-fledged bidirectional path
tracer with  multiple importance sampling, and this choice is controlled by the
\code{bidirectional} flag. The unidirectional path tracer is generally
much faster, but it produces lower-quality samples. Depending on the input, either may be preferable.
\vspace{-7mm}
\paragraph{Caveats:}
There are a few general caveats about MLT-type algorithms that are good
to know. The first one is that they only render ``relative'' output images,
meaning that there is a missing scale factor that must be applied to
obtain proper scene radiance values. The implementation in Mitsuba relies
on an additional Monte Carlo estimator to recover this scale factor. By
default, it uses 100K samples (controlled by the \code{luminanceSamples}
parameter), which should be adequate for most applications.

The second caveat is that the amount of computational expense
associated with a pixel in the output image is roughly proportional to
its intensity. This means that when a bright object (e.g. the sun) is
visible in a rendering, most resources are committed to rendering the
sun disk at the cost of increased variance everywhere else. Since this is
usually not desired, the \code{twoStage} parameter can be used to
enable \emph{Two-stage MLT} in this case.

In this mode of operation, the renderer first creates a low-resolution
version of the output image to determine the approximate distribution of
luminance values. The second stage then performs the actual rendering, while
using the previously collected information to ensure that
the amount of time spent rendering each pixel is uniform.

The third caveat is that, while PSMLT can work with scenes that are extremely
difficult for other methods to handle, it is not particularly efficient
when rendering simple things such as direct illumination (which is more easily
handled by a brute-force type algorithm). By default, the
implementation in Mitsuba therefore delegates this to such a method
(with the desired quality being controlled by the \code{directSamples} parameter).
In very rare cases when direct illumination paths are very difficult to find,
it is preferable to disable this separation so that PSSMLT is responsible
for everything. This can be accomplished by setting
\code{directSamples=-1}.
\plugin{mlt}{Path Space Metropolis Light Transport}
\order{10}
\parameters{
    \parameter{maxDepth}{\Integer}{Specifies the longest path depth
        in the generated output image (where \code{-1} corresponds to $\infty$).
        A value of \code{1} will only render directly visible light sources.
        \code{2} will lead to single-bounce (direct-only) illumination,
        and so on. \default{\code{-1}}
    }
    \parameter{directSamples}{\Integer}{
        By default, the implementation renders direct illumination component
        separately using the \pluginref{direct} plugin, which
        uses low-discrepancy number sequences for superior performance
        (in other words, it is \emph{not} handled by MLT). This
        parameter specifies the number of samples allocated to that method. To
        force MLT to be responsible for the direct illumination
        component as well, set this to \code{-1}. \default{16}
    }
    \parameter{luminanceSamples}{\Integer}{
       MLT-type algorithms create output images that are only
       \emph{relative}. The algorithm can e.g. determine that a certain pixel
       is approximately twice as bright as another one, but the absolute
       scale is unknown. To recover it, this plugin computes
       the average luminance arriving at the sensor by generating a
       number of samples. \default{\code{100000} samples}
    }
    \parameter{twoStage}{\Boolean}{Use two-stage MLT?
      See \pluginref{pssmlt} for details.\!\default{{\footnotesize\code{false}}}\!}
    \parameter{bidirectional\showbreak\newline Mutation,\vspace{1mm}
       [lens,multiChain,\newline caustic,manifold]\showbreak\newline Perturbation}{\Boolean}{
      These parameters can be used to pick the individual mutation and perturbation
      strategies that will be used to explore path space. By default, the original set
      by Veach and Guibas is enabled (i.e. everything except the manifold
      perturbation). It is possible to extend
      this integrator with additional custom perturbations strategies if needed.
    }
    \parameter{lambda}{\Float}{
        Jump size of the manifold perturbation \default{50}}
}
Metropolis Light Transport (MLT) is a seminal rendering technique proposed by Veach and
Guibas \cite{Veach1997Metropolis}, which applies the Metropolis-Hastings
algorithm to the path-space formulation of light transport.
Please refer to the \pluginref{pssmlt} page for a general description of MLT-type
algorithms and a list of caveats that also apply to this plugin.

Like \pluginref{pssmlt}, this integrator explores the space of light paths,
searching with preference for those that carry a significant amount of
energy from an emitter to the sensor. The main difference is that PSSMLT
does this exploration by piggybacking on another rendering technique and
``manipulating'' the random number stream that drives it, whereas MLT does
not use such an indirection: it operates directly on the actual light
paths.

This means that the algorithm has access to considerably more
information about the problem to be solved, which allows it to perform a
directed exploration of certain classes of light paths. The main downside
is that the implementation is rather complex, which may make it more
susceptible to unforeseen problems.
Mitsuba reproduces the full MLT
algorithm except for the lens subpath mutation\footnote{In experiments,
it was not found to produce sigificant convergence improvements and was
subsequently removed.}. In addition, the plugin also provides the
manifold perturbation proposed by Jakob and Marschner \cite{Jakob2012Manifold}.

\renderings{
   \includegraphics[width=\textwidth]{images/integrator_mlt_sketch.pdf}\hfill\,
}

To explore the space of light paths, MLT iteratively makes changes
to a light path, which can either be large-scale \emph{mutations} or small-scale
\emph{perturbations}. Roughly speaking, the \emph{bidirectional mutation} is used
to jump between different classes of light paths, and each one of the perturbations is
responsible for efficiently exploring some of these classes.
All mutation and perturbation strategies can be mixed and matched as
desired, though for the algorithm to work properly, the bidirectional
mutation must be active and perturbations should be selected
as required based on the types of light paths that are present in the
input scene. The following perturbations are available:

\begin{enumerate}[(a)]
\item \emph{Lens perturbation}: this perturbation slightly varies the outgoing
direction at the camera and propagates the resulting ray until it encounters
the first non-specular object. The perturbation then attempts to create a connection to the
(unchanged) remainder of the path.
\item \emph{Caustic perturbation}: essentially a lens perturbation
that proceeds in the opposite direction.
\item \emph{Multi-chain perturbation}: used when there are several chains
of specular interactions, as seen in the swimming pool example above.
After an initial lens perturbation, a cascade of additional perturbations
is required until a connection to the
remainder of the path can finally be established. Depending on the
path type, the entire path may be changed by this.
\item \emph{Manifold perturbation}: this perturbation was designed to
subsume and extend the previous three approaches.
It creates a perturbation at an arbitrary
position along the path, proceeding in either direction. Upon encountering
a chain of specular interactions, it numerically solves for a
connection path (as opposed to the cascading mechanism employed by the
multi-chain perturbation).
\end{enumerate}
\plugin{erpt}{Energy redistribution path tracing}
\order{11}
\parameters{
    \parameter{maxDepth}{\Integer}{Specifies the longest path depth
        in the generated output image (where \code{-1} corresponds to $\infty$).
        A value of \code{1} will only render directly visible light sources.
        \code{2} will lead to single-bounce (direct-only) illumination,
        and so on. \default{\code{-1}}
    }
    \parameter{numChains}{\Float}{
        On average, how many Markov Chains should be started per
        pixel? \default{\code{1}}
    }
    \parameter{maxChains}{\Float}{
        How many Markov Chains should be started \emph{at most} (per
        pixel) \default{\code{0}, i.e. this feature is not used}
    }
    \parameter{chainLength}{\Integer}{
        Specifies the number of perturbation steps that are executed per Markov Chain\default{1}.
    }
    \parameter{directSamples}{\Integer}{
        By default, the implementation renders direct illumination component
        separately using the \pluginref{direct} plugin, which
        uses low-discrepancy number sequences for superior performance
        (in other words, it is \emph{not} handled by ERPT). This
        parameter specifies the number of samples allocated to that method. To
        force MLT to be responsible for the direct illumination
        component as well, set this to \code{-1}. \default{\code{16}}
    }
    \parameter{[lens,multiChain,\!\!\!\!\newline caustic,manifold]\showbreak
      \newline Perturbation}{\Boolean}{
      These parameters can be used to pick the individual perturbation
      strategies that will be used to explore path space. By default, the original set
      by Veach and Guibas is enabled (i.e. everything except the manifold
      perturbation).
    }
    \parameter{lambda}{\Float}{
        Jump size of the manifold perturbation \default{\code{50}}}
}
\renderings{
 \rendering{A brass chandelier with 24 glass-enclosed bulbs}{integrator_mept_luminaire}
 \rendering{Glossy reflective and refractive ableware, lit by the
 chandelier on the left}{integrator_mept_tableware}
 \vspace{-2mm}
 \caption{An interior scene with complex specular and near-specular light paths,
 illuminated entirely through caustics. Rendered by this plugin
 using the manifold perturbation. This scene was designed by Olesya
 Isaenko.\vspace{2mm}}
}
\renderings{
\medrendering{Seed paths generated using bidirectional path tracing.
 Note the high variance of paths that involve reflection of sunlight by
 the torus.}{integrator_erpt_seeds}
\medrendering{Result after running the perturbations of Veach and Guibas
 for 800 steps. Some convergence issues remain.}{integrator_erpt_mlt}
\medrendering{Result after running the manifold
perturbation for the same amount of time}{integrator_erpt_mept}
}

Energy Redistribution Path Tracing (ERPT) by Cline et al. \cite{Cline2005Energy}
combines Path Tracing with the perturbation strategies of Metropolis Light Transport.


An initial set of \emph{seed paths} is generated using a standard bidirectional
path tracer, and for each one, a MLT-style Markov Chain is subsequently started
and executed for some number of steps.
This has the effect of redistributing the energy of the individual samples
over a larger area, hence the name of this method.

\begin{wrapfigure}{R}{0.30\textwidth}
 \fbox{\includegraphics[width=.29\textwidth]{images/integrator_mept_egg.jpg}}
 \caption{Another view, now with exterior lighting.}
 \vspace{-2mm}
\end{wrapfigure}
This is often a good choice when a (bidirectional) path tracer produces mostly reasonable
results except that it finds certain important types of light paths too rarely.
ERPT can
then explore all of the neighborhing paths as well, to prevent the
original sample from showing up as a ``bright pixel'' in the output image.

This plugin shares all the perturbation strategies of the \pluginref{mlt} plugin, and
the same rules for selecting them apply. In contrast to the original
paper by Cline et al., the Mitsuba implementation uses a bidirectional
(rather than an unidirectional) bidirectional path tracer to create seed paths.
Also, since they add bias to the output, this plugin does not use the image
post-processing filters proposed by the authors.

The mechanism for selecting Markov Chain seed paths deserves an
explanation: when commencing work on a pixel in the output image, the
integrator first creates a pool of seed path candidates. The size of this
pool is given by the \code{samplesPerPixel} parameter of the sample
generator. This should be large enough so that the integrator has a
representative set of light paths to work with.

Subsequently, one or more of these candidates are chosen (determined by
\code{numChains} and \code{maxChains} parameter). For each one, a Markov
Chain is created that has an initial configuration matching the seed path.
It is simulated for \code{chainLength} iterations, and each intermediate
state is recorded in the output image.
\plugin{ptracer}{Adjoint particle tracer}
\order{12}
\parameters{
    \parameter{maxDepth}{\Integer}{Specifies the longest path depth
        in the generated output image (where \code{-1} corresponds to $\infty$).
        A value of \code{1} will only render directly visible light sources.
        \code{2} will lead to single-bounce (direct-only) illumination,
        and so on. \default{\code{-1}}
    }
    \parameter{rrDepth}{\Integer}{Specifies the minimum path depth, after
       which the implementation will start to use the ``russian roulette''
       path termination criterion. \default{\code{5}}
    }
    \parameter{granularity}{\Integer}{
       Specifies the work unit granularity used to parallize the particle
       tracing task. This should be set high enough so that accumulating
       partially exposed images (and potentially sending them over the network)
       is not the bottleneck.
       \default{200K particles per work unit, i.e. \code{200000}}
    }
    \parameter{bruteForce}{\Boolean}{
       If set to \code{true}, the integrator does not attempt to create
       connections to the sensor and purely relies on hitting it via ray
       tracing. This is mainly intended for debugging purposes.
       \default{\code{false}}
    }
}

This plugin implements a simple adjoint particle tracer. It does
essentially the exact opposite of the simple volumetric path tracer
(\pluginref[volpathsimple]{volpath\_simple}): instead of tracing rays from
the sensor and attempting to connect them to the light source, this
integrator shoots particles from the light source and attempts to connect
them to the sensor.

Usually, this is a relatively useless rendering technique due to
its high variance, but there are some cases where it excels.
In particular, it does a good job on scenes where most scattering
events are directly visible to the camera.

When rendering with a finite-aperture sensor (e.g. \pluginref{thinlens})
this integrator is able to intersect the actual aperture, which allows
it to handle certain caustic paths that would otherwise not be visible.

It also supports a specialized ``brute force'' mode, where the integrator
does not attempt to create connections to the sensor and purely relies on
hitting it via ray tracing. This is one of the worst conceivable rendering
and not recommended for any applications. It is mainly included for
debugging purposes.

The number of traced particles is given by the number of ``samples per
pixel'' of the sample generator times the pixel count of the output image.
For instance, 16 samples per pixel on a 512$\times$512 image will cause 4M particles
to be generated.

\remarks{
   \item This integrator does not currently work with subsurface scattering
   models.
}
\plugin{adaptive}{Adaptive integrator}
\order{13}
\parameters{
    \parameter{maxError}{\Float}{Maximum relative error
        threshold\default{0.05}}
    \parameter{pValue}{\Float}{
        Required p-value to accept a sample \default{0.05}
    }
    \parameter{maxSampleFactor}{\Integer}{
        Maximum number of samples to be generated \emph{relative} to the
        number of configured pixel samples. The adaptive integrator
        will stop after this many samples, regardless of whether
        or not the error criterion was satisfied.
        A negative value will be interpreted as $\infty$.
        \default{32---for instance, when 64 pixel samples are configured in
        the \code{sampler}, this means that the adaptive integrator
        will give up after 32*64=2048 samples}
    }
}

This ``meta-integrator'' repeatedly invokes a provided sub-integrator
until the computed radiance values satisfy a specified relative error bound
(5\% by default) with a certain probability (95\% by default). Internally,
it uses a Z-test to decide when to stop collecting samples. While repeatedly
applying a Z-test in this manner is not good practice in terms of
a rigorous statistical analysis, it provides a useful mathematically
motivated stopping criterion.

\begin{xml}[caption={An example how to make the \pluginref{path} integrator adaptive}]
<integrator type="adaptive">
    <integrator type="path"/>
</integrator>
\end{xml}

\remarks{
   \item The adaptive integrator needs a variance estimate to work
    correctly. Hence, the underlying sample generator should be set to a reasonably
    large number of pixel samples (e.g. 64 or higher) so that this estimate can be obtained.
   \item This plugin uses a relatively simplistic error heuristic that does not
   share information between pixels and only reasons about variance in image space.
   In the future, it will likely be replaced with something more robust.
}
\plugin{vpl}{Virtual Point Light integrator}
\order{14}
\parameters{
    \parameter{maxDepth}{\Integer}{Specifies the longest path depth
        in the generated output image (where \code{-1} corresponds to $\infty$).
        A value of \code{2} will lead to direct-only illumination.
        \default{\code{5}}
    }
    \parameter{shadowMap\showbreak Resolution}{\Integer}{
      Resolution of the shadow maps that are used
      to compute the point-to-point visibility \default{512}
    }
    \parameter{clamping}{\Float}{
      A relative clamping factor between $[0,1]$ that is
      used to control the rendering artifact discussed below.
      \default{0.1}
    }
}

This integrator implements a hardware-accelerated global illumination
rendering technique based on the Instant Radiosity method by Keller
\cite{Keller1997Instant}. This is the same approach that is also used in
Mitsuba's real-time preview; the reason for providing it as a separate
integrator plugin is to enable automated (e.g. scripted) usage.

The method roughly works as follows: during a pre-process pass, any present direct
and indirect illumination is converted into a set of \emph{virtual point light}
sources (VPLs). The scene is then separately rendered many times, each time using
a different VPL as a source of illumination. All of the renderings created in this
manner are accumulated to create the final output image.

Because the individual rendering steps can be exectuted on a
graphics card, it is possible to render many (i.e. 100-1000) VPLs
per second. The method is not without problems, however. In particular,
it performs poorly when rendering glossy materials, and it produces
artifacts in corners and creases . Mitsuba automatically limits
the ``glossyness'' of materials to reduce the effects of the former
problem. A \code{clamping} parameter is provided to control the latter
(see the figure below). The number of samples per pixel specified to
the sampler is interpreted as the number of VPLs that should be rendered.

\renderings{
    \rendering{\code{clamping=0}: With clamping fully disabled, bright
    blotches appear in corners and creases.}{integrator_vpl_clamping0}
    \rendering{\code{clamping=0.3}: Higher clamping factors remove these
    artifacts, but they lead to visible energy loss (the rendering
    is too dark in certain areas). The default of \code{0.1} is
    usually reasonable.}{integrator_vpl_clamping03}
}
\plugin{irrcache}{Irradiance caching integrator}
\order{15}
\parameters{
    \parameter{resolution}{\Integer}{Elevational resolution of the stratified
     final gather hemisphere. The azimuthal resolution is two times this value. \default{14, i.e. $2\cdot14^2$=392 samples in total}}
    \parameter{quality}{\Float}{Quality factor (the $\kappa$ parameter of
    Tabellion et al. \cite{Tabellion2004Approximate})\default{1.0, which is adequate for most cases}}
    \parameter{gradients}{\Boolean}{Use irradiance gradients \cite{Ward1992Irradiance}?\default{\code{true}}}
    \parameter{clampNeighbor}{\Boolean}{Use neighbor clamping \cite{Krivanek2006Making}?\default{\code{true}}}
    \parameter{clampScreen}{\Boolean}{Use a screen-space clamping criterion \cite{Tabellion2004Approximate}? \default{\code{true}}}
    \parameter{overture}{\Boolean}{Do an overture pass before starting the main rendering process?
     Usually a good idea.\default{\code{true}}}
    \parameter{quality\showbreak Adjustment}{\Float}{When an overture pass is used, Mitsuba subsequently reduces
     the quality parameter by this amount to interpolate amongst more samples, creating a visually
     smoother result. \default{0.5}}
    \parameter{indirectOnly}{\Boolean}{Only show the indirect illumination? This can be useful to check
     the interpolation quality. \default{\code{false}}}
    \parameter{debug}{\Boolean}{Visualize the sample placement? \default{\code{false}}}
}
\renderings{
 \unframedbigrendering{Illustration of the effect of the different optimizatations
  that are provided by this plugin}{integrator_irrcache}
}
This ``meta-integrator'' implements \emph{irradiance caching} by Ward and Heckbert
\cite{Ward1988Ray}. This method computes and caches irradiance information at a sparse
set of scene locations and efficiently determines approximate values at other
locations using interpolation.

This plugin only provides the caching and interpolation part---another plugin
is still needed to do the actual computation of irradiance values at cache points.
This is done using nesting, e.g. as follows:

\begin{xml}[caption={Instantiation of a photon mapper with irradiance caching}]
<integrator type="irrcache">
    <integrator type="photonmapper"/>
</integrator>
\end{xml}

When a radiance query involves a non-diffuse material, all computation
is forwarded to the sub-integrator, i.e. \pluginref{irrcache} is passive.
Otherwise, the existing cache points are interpolated to approximate the
emitted radiance, or a new cache point is created if the resulting accuracy
would be too low.
By default, this integrator also performs a distributed overture pass before
rendering, which is recommended to avoid artifacts resulting from the
addition of samples as rendering proceeds.

Note that wrapping an integrator into \pluginref{irrcache} adds one extra
light bounce. For instance, the method resulting from using \pluginref{direct}
in an irradiance cache renders two-bounce direct illumination.

The generality of this implementation allows it to be used in conjunction
with photon mapping (the most likely application) as well as all other
sampling-based integrators in Mitsuba. Several optimizations are used to
improve the achieved interpolation quality, namely irradiance gradients
\cite{Ward1992Irradiance}, neighbor clamping \cite{Krivanek2006Making}, a screen-space
clamping metric and an improved error function \cite{Tabellion2004Approximate}.
\plugin{multichannel}{Multi-channel integrator}
\order{16}
\parameters{
    \parameter{\Unnamed}{\Integrator}{One or more sub-integrators whose output
    should be rendered into a combined multi-channel image}
}

The multi-channel integrator groups several sub-integrators together
and invokes them at the same time for each pixel; the result from each
integrator is written into a separate channel of the output image.
This could include things like surface normals or the distance
from the camera (via the \pluginref{field} plugin) or ambient occlusion
(via the \pluginref{ao} plugin).
In this way, this integrator can be a powerful tool for unusual applications
of Mitsuba, e.g. to create reference data for computer vision algorithms. Currently, it only
works with a subset of the other plugins---see the red box for details.

Thee \code{multichannel} plugin also disables certain checks for negative or infinite
radiance values during rendering that normally cause warnings to be emitted.
This is simply to process extracted fields for which it is fine
to take on such values.

The following example contains a typical setup for rendering an 7 channel EXR image:
3 for a path traced image (RGB), 3 for surface normals
(encoded as RGB), and 1 channel for the ray distance measured from the camera.

\vspace{2mm}
\begin{xml}
<scene>
    <integrator type="multichannel">
        <integrator type="path"/>
        <integrator type="field">
            <string name="field" value="shNormal"/>
        </integrator>
        <integrator type="field">
            <string name="field" value="distance"/>
        </integrator>
    </integrator>

    <sensor type="perspective">
        <sampler type="halton">
            <integer name="sampleCount" value="32"/>
        </sampler>
        <film type="hdrfilm">
            <string name="pixelFormat" value="rgb, rgb, luminance"/>
            <string name="channelNames" value="color, normal, distance"/>
        </film>
    </sensor>
    <!-- **** scene contents **** -->
</scene>
\end{xml}

\remarks{
\item Requires the \pluginref{hdrfilm} or \pluginref{tiledhdrfilm}.
\item All nested integrators must
conform to Mitsuba's basic \emph{SamplingIntegrator} interface.
Currently, only a few of them do this, including:
\pluginref{field}, \pluginref{ao}, \pluginref{direct}, \pluginref{path},
\pluginref{volpath}, \pluginref[volpathsimple]{volpath\_simple},
and \pluginref{irrcache}.
}
\plugin{field}{Field extraction integrator}
\order{17}
\parameters{
    \parameter{field}{\String}{Denotes the name of the field that should be extracted.
       The following choices are possible:
       \begin{itemize}
           \setlength{\itemsep}{1pt}
           \setlength{\parskip}{1pt}
           \item \code{position}: 3D position in world space
           \item \code{relPosition}: 3D position in camera space
           \item \code{distance}: Ray distance to the shading point
           \item \code{geoNormal}: Geometric surface normal
           \item \code{shNormal}: Shading surface normal
           \item \code{uv}: UV coordinate value
           \item \code{albedo}: Albedo value of the BSDF
           \item \code{shapeIndex}: Integer index of the high-level shape
           \item \code{primIndex}: Integer shape primitive index
       \end{itemize}
    }
    \parameter{undefined}{\Spectrum\Or\Float}{Value that should be returned when
                          there is no intersection \default{0}}
}

This integrator extracts a requested field of from the intersection records of shading
points and converts the resulting data into color values. It is meant to be used in conjunction with
\pluginref{multichannel} to dump auxiliary information (such as depth or surface normals
of surfaces seen by the camera) into extra channels of a rendered image, for instance to
create benchmark data for computer vision applications.
Please refer to the documentation of \pluginref{multichannel} for an example.
\plugin{motion}{Motion and specular motion vector integrator}
\parameters{
    \parameter{time}{\Float}{
      Denotes the time stamp of the target frame of the motion vectors.
      The current frame is specified via the sensor's \code{shutterOpen}
      and \code{shutterClose} parameters, which should both be set to
      the same value. \default{0}
    }
    \parameter{time}{\String}{
       Path configuration of the desired motion vectors:
       \begin{enumerate}[(i)]
          \item \textbf{d}: Primary (non-specular) hit points\vspace{-1mm}
          \item \textbf{rd}: A non-specular surface seen through a specular reflection\vspace{-1mm}
          \item \textbf{ttd}: A non-specular surface seen through a pair of specular refractions\vspace{-1mm}
          \item \textbf{trtd}: A non-specular surface seen through a sequence of refraction, reflection, and refraction events.\vspace{-1mm}
       \end{enumerate}
       etc.
    }
    \parameter{derivativesOnly}{\Boolean}{
      By default, the Manifold Exploration technique is used to accurately
      solve the underlying specular flow problem. When this parameter is set to \code{true},
      the nonlinear solver is deactivated, and only first-order extrapolations are provided.
      \default{\code{false}}
    }
    \parameter{glossyThreshold}{\Float}{
       Threshold on the roughness parameter of reflectance models
       to be classified as specular.
       \default{\texttt{0}, i.e.~only perfectly specular materials are classified as specular}
    }
    \parameter{maxTimeSteps}{\Integer}{
       Maximum number of temporal sub-steps \default{5}
    }
    \parameter{maxSpaceSteps}{\Integer}{
       Maximum number of spatial sub-steps \default{10}
    }
}
This integrator extracts motion vectors for animated input scenes, alternatively
at primary hit points or at hit points observed through sequences of reflective
and refractive objects. The first two color components (R and G) of the resulting
rendering specify the screen-space motion in 2D pixel coordinates, and the last
component (B) denotes the change in distance of the observed 3D point to the
camera position. Sometimes a specular path cannot be tracked from one frame to the other,
e.g. because it does not exist, or because the solver did not converge. In this case,
the pixel color is set to infinity. The images on the following page show
motion vectors obtained for a sphere that is moving from the left to the right.

\renderings{
    \rendering{Input scene at time $t=0$}{integrator_motion_sphere_1}
    \rendering{Input scene at time $t=1$}{integrator_motion_sphere_2}
}
\renderings{
    \medrendering{\code{config="d"}}{integrator_motion_path_d}
    \medrendering{\code{config="rd"}}{integrator_motion_path_rd}
    \medrendering{\code{config="ttd"}}{integrator_motion_path_ttd}
}
\renderings{
    \rendering{\code{config="trtd"}}{integrator_motion_path_trtd}
    \rendering{\code{config="trrtd"}}{integrator_motion_path_trrtd}
}
\clearpage

\begin{xml}[caption={Exemplary scene configuration for computing specular motion vectors}]
<scene>
    <integrator type="motion">
        <string name="config" value="ttd"/>
        <float name="time" value="1"/>
    </integrator>

    <shape type="serialized">
        <string name="filename" value="..."/>
        <animation name="toWorld">
            <transform time="0">
                <!-- Transformation at time zero -->
            </transform>
            <transform time="1">
                <!-- Transformation at time one -->
            </transform>
        </animation>
         <bsdf type="dielectric"/>
    </shape>

    <sensor type="perspective">
        <float name="shutterOpen" value="0"/>
        <float name="shutterClose" value="0"/>

        <sampler type="ldsampler">
            <integer name="sampleCount" value="1"/>
            <boolean name="pixelCenters" value="true"/>
        </sampler>

        <film type="hdrfilm" id="film">
            <string name="pixelFormat" value="rgb"/>
            <boolean name="banner" value="false"/>
            <rfilter type="box"/>
        </film>
    </sensor>
</scene>
\end{xml}
\input{section_samplers}
\plugin{independent}{Independent sampler}
\order{1}
\parameters{
    \parameter{sampleCount}{\Integer}{
      Number of samples per pixel \default{4}
    }
}

\renderings{
    \unframedrendering{A projection of the first 1024 points
    onto the first two dimensions. Note the sample clumping.}{sampler_independent}
}

The independent sampler produces a stream of independent and uniformly
distributed pseudorandom numbers. Internally, it relies on a fast SIMD version
of the Mersenne Twister random number generator \cite{Saito2008SIMD}.

This is the most basic sample generator; because no precautions are taken to avoid
sample clumping, images produced using this plugin will usually take longer to converge.
In theory, this sampler is initialized using a deterministic procedure, which means
that subsequent runs of Mitsuba should create the same image. In practice, when
rendering with multiple threads and/or machines, this is not true anymore, since the
ordering of samples is influenced by the operating system scheduler.

Note that the Metropolis-type integrators implemented in Mitsuba are incompatible with
the more sophisticated sample generators shown in this section. They \emph{require} this
specific sampler and refuse to work otherwise.
\plugin{stratified}{Stratified sampler}
\order{2}
\parameters{
    \parameter{sampleCount}{\Integer}{
      Number of samples per pixel; should be a perfect square
      (e.g. 1, 4, 9, 16, 25, etc.), or it will be rounded up to the
      next one \default{4}
    }
    \parameter{dimension}{\Integer}{
      Effective dimension, up to which stratified samples are provided. The
      number here is to be interpreted as the number of subsequent 1D or 2D sample
      requests that can be satisfied using ``good'' samples. Higher high values
      increase both storage and computational costs.
      \default{4}
    }
}
\renderings{
    \unframedrendering{A projection of the first 1024 points
    onto the first two dimensions.}{sampler_stratified}
    \unframedrendering{The same samples shown together with the
    underlying strata for illustrative purposes}{sampler_stratified_strata}
}

The stratified sample generator divides the domain into a discrete number
of strata and produces a sample within each one of them. This generally leads to less
sample clumping when compared to the independent sampler, as well as better
convergence. Due to internal storage costs, stratified samples are only provided up to a
certain dimension, after which independent sampling takes over.

Like the \pluginref{independent} sampler, multicore and network renderings
will generally produce different images in subsequent runs due to the nondeterminism
introduced by the operating system scheduler.
\plugin{ldsampler}{Low discrepancy sampler}
\order{3}
\parameters{
    \parameter{sampleCount}{\Integer}{
      Number of samples per pixel; should be a power of two
      (e.g. 1, 2, 4, 8, 16, etc.), or it will be rounded up to the next one
      \default{4}
    }
    \parameter{dimension}{\Integer}{
      Effective dimension, up to which low discrepancy samples are provided. The
      number here is to be interpreted as the number of subsequent 1D or 2D sample
      requests that can be satisfied using ``good'' samples. Higher high values
      increase both storage and computational costs.
      \default{4}
    }
}
\vspace{-2mm}
\renderings{
    \unframedrendering{A projection of the first 1024 points
    onto the first two dimensions.}{sampler_ldsampler_0}
    \unframedrendering{A projection of the first 1024 points
    onto the 32 and 33th dimension, which look almost identical. However,
    note that the points have been scrambled to reduce
    correlations between dimensions.}{sampler_ldsampler_32}
}
This plugin implements a simple hybrid sampler that combines aspects of a Quasi-Monte
Carlo sequence with a pseudorandom number generator based on a technique proposed
by Kollig and Keller \cite{Kollig2002Efficient}.
It is a good and fast general-purpose sample generator and therefore chosen as
the default option in Mitsuba. Some of the QMC samplers in the following pages can generate
even better distributed samples, but this comes at a higher cost in terms of performance.

Roughly, the idea of this sampler is that all of the individual 2D sample dimensions are
first filled using the same (0, 2)-sequence, which is then randomly scrambled and permuted
using numbers generated by a Mersenne Twister pseudorandom number generator \cite{Saito2008SIMD}.
Note that due to internal storage costs, low discrepancy samples are only provided
up to a certain dimension, after which independent sampling takes over.
The name of this plugin stems from the fact that (0, 2) sequences minimize the so-called
\emph{star disrepancy}, which is a quality criterion on their spatial distribution. By
now, the name has become slightly misleading since there are other samplers in Mitsuba
that just as much try to minimize discrepancy, namely the \pluginref{sobol} and
\pluginref{halton} plugins.

Like the \pluginref{independent} sampler, multicore and network renderings
will generally produce different images in subsequent runs due to the nondeterminism
introduced by the operating system scheduler.
\plugin{halton}{Halton QMC sampler}
\order{4}
\parameters{
    \parameter{sampleCount}{\Integer}{
       Number of samples per pixel \default{4}
    }
    \parameter{scramble}{\Integer}{
       This plugin can operate in one of three scrambling modes:
       \begin{enumerate}[(i)]
       \item When set to \code{0}, the implementation will provide the standard Halton sequence.

       \item When set to \code{-1}, the implementation will compute
       a scrambled variant of the Halton sequence based on permutations by
       Faure \cite{Faure1992Good}, which has better equidistribution properties
       in high dimensions.

       \item When set to a value greater than one, a random permutation is chosen based
       on this number. This is useful to break up temporally coherent noise when rendering
       the frames of an animation --- in this case, simply set the parameter to the current frame index.
       \end{enumerate}
       Default: \code{-1}, i.e. use the Faure permutations. Note that permutations rely on a
       precomputed table that consumes approximately 7 MiB of additional memory at run time.
    }
}
\renderings{
    \unframedrendering{Projection of the first 1024 points
    of the Faure-scrambled Halton seq. onto the first two dimensions.}{sampler_halton_scrambled_0}
    \unframedrendering{Projection of the first 1024 points
    of the Faure-scrambled Halton seq. onto the 32th and 33th dim.}{sampler_halton_scrambled_32}
}
This plugin implements a Quasi-Monte Carlo (QMC) sample generator based on the
Halton sequence. QMC number sequences are designed to reduce sample clumping
across integration dimensions, which can lead to a higher order of
convergence in renderings. Because of the deterministic character of the samples,
errors will manifest as grid or moir\'e patterns rather than random noise, but
these diminish as the number of samples is increased.

The Halton sequence in particular provides a very high quality point set that unfortunately
becomes increasingly correlated in higher dimensions. To ameliorate this problem, the Halton
points are usually combined with a scrambling permutation, and this is also the default.
Because everything that happens inside this sampler is completely deterministic and
independent of operating system scheduling behavior, subsequent runs of Mitsuba will always
compute the same image, and this even holds when rendering with multiple threads
and/or machines.
\renderings{
    \unframedrendering{A projection of the first 1024 points
    of the \emph{original} Halton sequence onto the first two dimensions, obtained by
    setting \code{scramble=0}}{sampler_halton_nonscrambled_0}
    \unframedrendering{A projection of the first 1024 points
    of the \emph{original} Halton sequence onto the 32th and 33th dimensions.
    Note the strong correlation -- a scrambled sequence is usually preferred to avoid this problem.}{sampler_halton_nonscrambled_32}
}
\renderings{
    \unframedrendering{A projection of the first 1024 points
    of a randomly scrambled Halton sequence onto the first two dimensions (\code{scramble=1}).}{sampler_halton_rscrambled_0}
    \unframedrendering{A projection of the first 1024 points
    of a randomly scrambled Halton sequence onto the 32th and 33th dimensions.}{sampler_halton_rscrambled_32}
}

By default, the implementation provides a scrambled variant of the Halton sequence based
on permutations by Faure \cite{Faure1992Good} that has better equidistribution properties
in high dimensions, but this can be changed using the \code{scramble} parameter.
Internally, the plugin uses a table of prime numbers to provide elements
of the Halton sequence up to a dimension of 1024. Because of this upper bound,
the maximum path depth of the integrator must be limited (e.g. to 100), or
rendering might fail with the following error message: \emph{Lookup dimension
exceeds the prime number table size! You may have to reduce the 'maxDepth'
parameter of your integrator}.

To support bucket-based renderings, the Halton sequence is internally enumerated
using a scheme proposed by Gr\"unschlo\ss\ et al. \cite{Grunschloss2010Enumerating};
the implementation in Mitsuba is based on a Python script by the authors of
this paper.

\remarks{
  \item This sampler is incompatible with Metropolis Light Transport (all variants).
  It interoperates poorly with Bidirectional Path Tracing and Energy Redistribution
  Path Tracing, hence these should not be used together. The \pluginref{sobol} QMC
  sequence is an alternative for the latter two cases, and \pluginref{ldsampler}
  works as well.
}
\plugin{hammersley}{Hammersley QMC sampler}
\order{5}
\parameters{
    \parameter{sampleCount}{\Integer}{
       Number of samples per pixel \default{4}
    }
    \parameter{scramble}{\Integer}{
       This plugin can operate in one of three scrambling modes:
       \begin{enumerate}[(i)]
       \item When set to \code{0}, the implementation will provide the standard Hammersley sequence.

       \item When set to \code{-1}, the implementation will compute
       a scrambled variant of the Hammersley sequence based on permutations by
       Faure \cite{Faure1992Good}, which has better equidistribution properties
       in high dimensions.

       \item When set to a value greater than one, a random permutation is chosen based
       on this number. This is useful to break up temporally coherent noise when rendering
       the frames of an animation --- in this case, simply set the parameter to the current frame index.
       \end{enumerate}
       Default: \code{-1}, i.e. use the Faure permutations. Note that permutations rely on a
       precomputed table that consumes approximately 7 MiB of additional memory at run time.
    }
}
\renderings{
    \unframedrendering{Projection of the first 1024 points
    of the Faure-scrambled sequence onto the first two dimensions.}{sampler_hammersley_0}
    \unframedrendering{Projection of the first 1024 points
    of the Faure-scrambled sequence onto the 32th and 33th dim.}{sampler_hammersley_32}
}
This plugin implements a Quasi-Monte Carlo (QMC) sample generator based on the
Hammersley sequence. QMC number sequences are designed to reduce sample clumping
across integration dimensions, which can lead to a higher order of
convergence in renderings. Because of the deterministic character of the samples,
errors will manifest as grid or moir\'e patterns rather than random noise, but
these diminish as the number of samples is increased.

The Hammerlsey sequence is closely related to the Halton sequence and yields a very
high quality point set that is slightly more regular (and has lower discrepancy),
especially in the first few dimensions. As is the case with the Halton sequence,
the points should be scrambled to reduce patterns that manifest due to correlations
in higher dimensions. Please refer to the \pluginref{halton} page for more information
on how this works.

Note that this sampler will cause odd-looking intermediate results when combined with rendering
techniques that trace paths starting at light source (e.g. \pluginref{ptracer})---these vanish
by the time the rendering process finishes.

\remarks{
  \item This sampler is incompatible with Metropolis Light Transport (all variants).
  It interoperates poorly with Bidirectional Path Tracing and Energy Redistribution
  Path Tracing, hence these should not be used together. The \pluginref{sobol} QMC
  sequence is an alternative for the latter two cases, and \pluginref{ldsampler}
  works as well.
}
\plugin{sobol}{Sobol QMC sampler}
\order{6}
\parameters{
    \parameter{sampleCount}{\Integer}{
      Number of samples per pixel \default{4}
    }
    \parameter{scramble}{\Integer}{
      This parameter can be used to set a scramble value to break up temporally coherent
      noise patterns. For stills, this is irrelevant. When rendering
      an animation, simply set it to the current frame index. \default{0}
    }
}

\renderings{
    \unframedrendering{A projection of the first 1024 points
    onto the first two dimensions.}{sampler_sobol_nonscrambled_0}
    \unframedrendering{A projection of the first 1024 points
    onto the 32 and 33th dimension.}{sampler_sobol_nonscrambled_32}
}

This plugin implements a Quasi-Monte Carlo (QMC) sample generator based on the
Sobol sequence. QMC number sequences are designed to reduce sample clumping
across integration dimensions, which can lead to a higher order of
convergence in renderings. Because of the deterministic character of the samples,
errors will manifest as grid or moir\'e patterns rather than random noise, but
these diminish as the number of samples is increased.

The Sobol sequence in particular provides a relatively good point set that can
be computed extremely efficiently. One downside is the susceptibility to pattern artifacts
in the generated image. To minimize these artifacts, it is advisable to use
a number of samples per pixel that is a power of two.

Because everything that happens inside this sampler is completely
deterministic and independent of operating system scheduling behavior, subsequent
runs of Mitsuba will always compute the same image, and this even holds when rendering
with multiple threads and/or machines.

The plugin relies on a fast implementation of the Sobol sequence by Leonhard
Gr\"unschlo\ss\ using direction numbers provided by Joe and Kuo
\cite{Joe2008Constructing}.
These direction numbers are given up to a dimension of 1024. Because of this
upper bound, the maximum path depth of the integrator must be limited (e.g. to 100), or
rendering might fail with the following error message: \emph{Lookup dimension
exceeds the direction number table size! You may have to reduce the 'maxDepth'
parameter of your integrator}.

Note that this sampler generates a $(0,2)$-sequence in the first two dimensions,
and therefore the point plot shown in (a) happens to match the corresponding plots of
\pluginref{ldsampler}. In higher dimensions, however, they behave rather differently.

When this sampler is used to perform parallel block-based renderings,
the sequence is internally enumerated using a scheme proposed and implemented
by Gr\"unschlo\ss\ et al. \cite{Grunschloss2010Enumerating}.
\remarks{
  \item This sampler is incompatible with Metropolis Light Transport (all variants).
}
\input{section_films}
\plugin{hdrfilm}{High dynamic range film}
\order{1}
\parameters{
    \parameter{width, height}{\Integer}{
      Width and height of the camera sensor in pixels
      \default{768, 576}
    }
    \parameter{fileFormat}{\String}{
      Denotes the desired output file format. The options
      are \code{openexr} (for ILM's OpenEXR format),
      \code{rgbe} (for Greg Ward's RGBE format),
      or \code{pfm} (for the Portable Float Map format)
      \default{\code{openexr}}
    }
    \parameter{pixelFormat}{\String}{Specifies the desired pixel format
        of output images. The options are \code{luminance},
        \code{luminanceAlpha}, \code{rgb}, \code{rgba}, \code{xyz},
        \code{xyza}, \code{spectrum}, and \code{spectrumAlpha}.
        For the \code{spectrum*} options, the number of written channels depends on
        the value assigned to \code{SPECTRUM\_SAMPLES} during compilation
        (see \secref{compiling} for details)
        \default{\code{rgb}}
    }
    \parameter{componentFormat}{\String}{Specifies the desired floating
        point component format of output images. The options are
        \code{float16}, \code{float32}, or \code{uint32}.
        \default{\code{float16}}
    }
    \parameter{cropOffsetX, cropOffsetY, cropWidth, cropHeight}{\Integer}{
      These parameters can optionally be provided to select a sub-rectangle
      of the output. In this case, Mitsuba will only render the requested
      regions. \default{Unused}
    }
    \parameter{attachLog}{\Boolean}{Mitsuba can optionally attach
        the entire rendering log file as a metadata field so that this
        information is permanently saved.
        \default{\code{true}, i.e. attach it}
    }
    \parameter{banner}{\Boolean}{Include a small Mitsuba banner in the
        output image? \default{\code{true}}
    }
    \parameter{highQualityEdges}{\Boolean}{
       If set to \code{true}, regions slightly outside of the film
       plane will also be sampled. This may improve the image
       quality at the edges, especially when using very large
       reconstruction filters. In general, this is not needed though.
       \default{\code{false}, i.e. disabled}
    }
    \parameter{\Unnamed}{\RFilter}{Reconstruction filter that should
    be used by the film. \default{\code{gaussian}, a windowed Gaussian filter}}
}

This is the default film plugin that is used when none is explicitly
specified. It stores the captured image as a high dynamic range OpenEXR file
and tries to preserve the rendering as much as possible by not performing any
kind of post processing, such as gamma correction---the output file
will record linear radiance values.

When writing OpenEXR files, the film will either produce a luminance, luminance/alpha,
RGB(A), XYZ(A) tristimulus, or spectrum/spectrum-alpha-based bitmap having a
\code{float16}, \code{float32}, or \code{uint32}-based internal representation
based on the chosen parameters.
The default configuration is RGB with a \code{float16} component format,
which is appropriate for most purposes. Note that the spectral output options
only make sense when using a custom build of Mitsuba that has spectral
rendering enabled (this is not the case for the downloadable release builds).
For OpenEXR files, Mitsuba also supports fully general multi-channel output;
refer to the \pluginref{multichannel} plugin for details on how this works.

The plugin can also write RLE-compressed files in the Radiance RGBE format
pioneered by Greg Ward (set \code{fileFormat=rgbe}), as well as the
Portable Float Map format (set \code{fileFormat=pfm}).
In the former case,
the \code{componentFormat} and \code{pixelFormat} parameters are ignored,
and the output is ``\code{float8}''-compressed RGB data.
PFM output is restricted to \code{float32}-valued images using the
\code{rgb} or \code{luminance} pixel formats.
Due to the superior accuracy and adoption of OpenEXR, the use of these
two alternative formats is discouraged however.

When RGB(A) output is selected, the measured spectral power distributions are
converted to linear RGB based on the CIE 1931 XYZ color matching curves and
the ITU-R Rec. BT.709-3 primaries with a D65 white point.

\begin{xml}[caption=Instantiation of a film that writes a full-HD RGBA OpenEXR file without the Mitsuba banner]
<film type="hdrfilm">
    <string name="pixelFormat" value="rgba"/>
    <integer name="width" value="1920"/>
    <integer name="height" value="1080"/>
    <boolean name="banner" value="false"/>
</film>
\end{xml}

\subsubsection*{Render-time annotations:}
\label{sec:film-annotations}
The \pluginref{ldrfilm} and \pluginref{hdrfilm} plugins support a
feature referred to as \emph{render-time annotations} to facilitate
record keeping.
Annotations are used to embed useful information inside a rendered image so
that this information is later available to anyone viewing the image.
Exemplary uses of this feature might be to store the frame or take number,
rendering time, memory usage, camera parameters, or other relevant scene
information.

Currently, two different types are supported: a \code{metadata} annotation
creates an entry in the metadata table of the image, which is preferable
when the image contents should not be touched. Alternatively, a \code{label}
annotation creates a line of text that is overlaid on top of the image. Note
that this is only visible when opening the output file (i.e. the line is not
shown in the interactive viewer).
The syntax of this looks as follows:

\begin{xml}
<film type="hdrfilm">
 <!-- Create a new metadata entry 'my_tag_name' and set it to the
      value 'my_tag_value' -->
 <string name="metadata['key_name']" value="Hello!"/>

 <!-- Add the label 'Hello' at the image position X=50, Y=80 -->
 <string name="label[50, 80]" value="Hello!"/>
</film>
\end{xml}

The \code{value="..."} argument may also include certain keywords that will be
evaluated and substituted when the rendered image is written to disk. A list all available
keywords is provided in Table~\ref{tbl:film-keywords}.

Apart from querying the render time,
memory usage, and other scene-related information, it is also possible
to `paste' an existing parameter that was provided to another plugin---for instance,
the camera transform matrix would be obtained as \code{\$sensor['toWorld']}. The name of
the active integrator plugin is given by \code{\$integrator['type']}, and so on.
All of these can be mixed to build larger fragments, as following example demonstrates.
The result of this annotation is shown in Figure~\ref{fig:annotation-example}.
\begin{xml}[mathescape=false]
<string name="label[10, 10]" value="Integrator: $integrator['type'],
  $film['width']x$film['height'], $sampler['sampleCount'] spp,
  render time: $scene['renderTime'], memory: $scene['memUsage']"/>
\end{xml}
\vspace{1cm}
\renderings{
\fbox{\includegraphics[width=.8\textwidth]{images/annotation_example}}\hfill\,
\caption{\label{fig:annotation-example}A demonstration of the label annotation feature
 given the example string shown above.}
}
\vspace{2cm}
\begin{table}[htb]
\centering
\begin{savenotes}
\begin{tabular}{ll}
\toprule
\code{\$scene['renderTime']}& Image render time, use \code{renderTimePrecise} for more digits.\\
\code{\$scene['memUsage']}& Mitsuba memory usage\footnote{The definition of this quantity unfortunately
varies a bit from platform to platform. On Linux and Windows, it denotes the total
amount of allocated RAM and disk-based memory that is private to the process (i.e. not
shared or shareable), which most intuitively captures the amount of memory required for
rendering. On OSX, it denotes the working set size---roughly speaking, this is the
amount of RAM apportioned to the process (i.e. excluding disk-based memory).}.
Use \code{memUsagePrecise} for more digits.\\
\code{\$scene['coreCount']}& Number of local and remote cores working on the rendering job\\
\code{\$scene['blockSize']}& Block size used to parallelize up the rendering workload\\
\code{\$scene['sourceFile']}& Source file name\\
\code{\$scene['destFile']}& Destination file name\\
\code{\$integrator['..']}& Copy a named integrator parameter\\
\code{\$sensor['..']}& Copy a named sensor parameter\\
\code{\$sampler['..']}& Copy a named sampler parameter\\
\code{\$film['..']}& Copy a named film parameter\\
\bottomrule
\end{tabular}
\end{savenotes}
\caption{\label{tbl:film-keywords}A list of all special
keywords supported by the annotation feature}
\end{table}

\plugin{tiledhdrfilm}{Tiled high dynamic range film}
\order{2}
\parameters{
    \parameter{width, height}{\Integer}{
      Width and height of the camera sensor in pixels
      \default{768, 576}
    }
    \parameter{cropOffsetX, cropOffsetY, cropWidth, cropHeight}{\Integer}{
      These parameters can optionally be provided to select a sub-rectangle
      of the output. In this case, Mitsuba will only render the requested
      regions. \default{Unused}
    }
    \parameter{pixelFormat}{\String}{Specifies the desired pixel format
        for OpenEXR output images. The options are \code{luminance},
        \code{luminanceAlpha}, \code{rgb}, \code{rgba}, \code{xyz},
        \code{xyza}, \code{spectrum}, and \code{spectrumAlpha}. In the latter two cases,
        the number of written channels depends on the value assigned to
        \code{SPECTRUM\_SAMPLES} during compilation (see Section~\ref{sec:compiling}
        for details)
        \default{\code{rgb}}
    }
    \parameter{componentFormat}{\String}{Specifies the desired floating
        point component format used for the output. The options are
        \code{float16}, \code{float32}, or \code{uint32}
        \default{\code{float16}}
    }

    \parameter{\Unnamed}{\RFilter}{Reconstruction filter that should
    be used by the film. \default{\code{gaussian}, a windowed Gaussian filter}}
}

This plugin implements a camera film that stores the captured image
as a \emph{tiled} high dynamic-range OpenEXR file. It is very similar to
\pluginref{hdrfilm}, the main difference being that it does not keep
the rendered image in memory. Instead, image tiles are directly written
to disk as they are being rendered, which enables renderings of extremely
large output images that would otherwise not fit into memory (e.g.
100K$\times$100K).

When the image can fit into memory, usage of this plugin is discouraged:
due to the extra overhead of tracking image tiles, the rendering process
will be slower, and the output files also generally do not compress as
well as those produced by \pluginref{hdrfilm}.

Based on the provided parameter values, the film will either write a luminance,
luminance/alpha, RGB(A), XYZ(A) tristimulus, or spectrum/spectrum-alpha-based
bitmap having a \code{float16}, \code{float32}, or \code{uint32}-based
internal representation. The default is RGB and \code{float16}.
Note that the spectral output options only make sense when using a
custom compiled Mitsuba distribution that has spectral rendering
enabled. This is not the case for the downloadable release builds.

When RGB output is selected, the measured spectral power distributions are
converted to linear RGB based on the CIE 1931 XYZ color matching curves and
the ITU-R Rec. BT.709 primaries with a D65 white point.

\remarks{
   \item This film is only meant for command line-based rendering. When
   used with \texttt{mtsgui}, the preview image will be black.
   \item This plugin is slower than \pluginref{hdrfilm}, and therefore should
   only be used when the output image is too large to fit into system memory.
}
\plugin{ldrfilm}{Low dynamic range film}
\order{3}
\parameters{
    \parameter{width, height}{\Integer}{
      Camera sensor resolution in pixels
      \default{768, 576}
    }
    \parameter{fileFormat}{\Integer}{
      The desired output file format:
      \code{png} or \code{jpeg}. \default{\code{png}}
    }
    \parameter{pixelFormat}{\String}{Specifies the pixel format
        of the generated image. The options are \code{luminance},
        \code{luminanceAlpha}, \code{rgb} or \code{rgba} for PNG output
        and \code{rgb} or \code{luminance} for JPEG output.
    }
    \parameter{tonemapMethod}{\String}{
      Method used to tonemap recorded radiance values
      \vspace{-2mm}
      \begin{enumerate}[(i)]
         \item \code{gamma}: Exposure and gamma correction (default)
         \vspace{-1mm}
         \item \code{reinhard}: Apply the
         tonemapping technique by Reinhard et al. \cite{Reinhard2002Photographic}
         followd by gamma correction.
      \vspace{-4mm}
      \end{enumerate}
    }
    \parameter{gamma}{\Float}{
      The gamma curve applied to correct the output image,
      where the special value -1 indicates sRGB. \default{-1}
    }
    \parameter{exposure}{\Float}{
      When \code{gamma} tonemapping is active, this parameter specifies
      an exposure factor in f-stops that is applied to the image before
      gamma correction (scaling the radiance values by $2^{\,\text{exposure}}$).
      \default{0, i.e. do not change the exposure}
    }
    \parameter{key}{\Float}{
      When \code{reinhard} tonemapping is active, this parameter in $(0,1]$ specifies
      whether a low-key or high-key image is desired.
      \default{0.18, corresponding to a middle-grey}
    }
    \parameter{burn}{\Float}{
      When \code{reinhard} tonemapping is active, this parameter in $[0,1]$ specifies how much
      highlights can burn out. \default{0, i.e. map all luminance values into the displayable range}
    }
    \parameter{banner}{\Boolean}{Include a banner in the
        output image?\default{\code{true}}
    }
    \parameter{cropOffsetX, cropOffsetY, cropWidth, cropHeight}{\Integer}{
      These parameters can optionally be provided to select a sub-rectangle
      of the output. In this case, Mitsuba will only render the requested
      regions. \default{Unused}
    }
    \parameter{highQualityEdges}{\Boolean}{
       If set to \code{true}, regions slightly outside of the film
       plane will also be sampled. This may improve image
       quality at the edges, but is not needed in general.
       \default{\code{false}}
    }
    \parameter{\Unnamed}{\RFilter}{Reconstruction filter that should
    be used by the film. \default{\code{gaussian}, a windowed Gaussian filter}}
}
This plugin implements a low dynamic range film that can write out 8-bit PNG
and JPEG images in various configurations. It provides basic tonemapping techniques
to map recorded radiance values into a reasonable displayable range. An alpha (opacity)
channel can be written if desired. By default, the plugin writes gamma-corrected
PNG files using the sRGB color space and no alpha channel.

This film is a good choice when low dynamic range output is desired
and the rendering setup can be configured to capture the relevant portion
of the dynamic range reliably enough so that the original HDR data can safely
be discarded. When this is not the case, it may be easier to use \pluginref{hdrfilm}
along with the batch tonemapper (\secref{tonemapper}).

By default, the plugin assumes that no special tonemapping needs to be done and simply
applies an exposure multiplier and sRGB gamma curve to the recorded radiance values
before converting them to 8 bit. When the dynamic range varies greatly, it may be
preferable to use the photographic tonemapping technique by Reinhard et al.
\cite{Reinhard2002Photographic}, which can be activated by setting \code{tonemapMethod=reinhard}.

Note that the interactive tonemapper that is available in the graphical user interface \code{mtsgui}
interoperates with this plugin. In particular, when saving the scene
(\emph{File}$\to$\emph{Save}), the currently active tonemapper
settings are automatically exported into the updated scene file.

The RGB values exported by this plugin correspond to the ITU-R Rec. BT. 709-3
primaries with a D65 white point. When $\texttt{gamma}$ is set to $\code{-1}$ (the default),
the output is in the sRGB color space and will display as intended on compatible devices.

Note that this plugin supports render-time \emph{annotations}, which
are described on page~\pageref{sec:film-annotations}.
\plugin{mfilm}{MATLAB / Mathematica / NumPy film}
\order{4}
\parameters{
    \parameter{width, height}{\Integer}{
      Width and height of the sensor in pixels
      \default{1, 1}
    }
    \parameter{cropOffsetX, cropOffsetY, cropWidth, cropHeight}{\Integer}{
      These parameters can optionally be provided to select a sub-rectangle
      of the output. In this case, Mitsuba will only render the requested
      regions. \default{Unused}
    }
    \parameter{fileFormat}{\String}{
      Specifies the desired output format; must be one of
      \code{matlab}, \code{mathematica}, or \code{numpy}. \default{\code{matlab}}
    }
    \parameter{digits}{\Integer}{
      Number of significant digits to be written \default{4}
    }
    \parameter{variable}{\String}{
      Name of the target variable \default{\code{"data"}}
    }
    \parameter{pixelFormat}{\String}{Specifies the desired pixel format
        of the generated image. The options are \code{luminance},
        \code{luminanceAlpha}, \code{rgb}, \code{rgba}, \code{spectrum},
        and \code{spectrumAlpha}. In the latter two cases,
        the number of written channels depends on the value assigned to
        \code{SPECTRUM\_SAMPLES} during compilation (see Section~\ref{sec:compiling}
        for details) \default{\code{luminance}}
    }
    \parameter{highQualityEdges}{\Boolean}{
       If set to \code{true}, regions slightly outside of the film
       plane will also be sampled. This may improve the image
       quality at the edges, especially when using very large
       reconstruction filters. In general (and particularly using the
       default box filter), this is not needed though.
       \default{\code{false}, i.e. disabled}
    }
    \parameter{\Unnamed}{\RFilter}{Reconstruction filter that should
    be used by the film. \default{\code{box}, a simple box filter}}
}

\renderings{
    \rendering{Importing and tonemapping an image in Mathematica}{film_mfilm_mathematica.jpg}
}

This plugin provides a camera film that exports spectrum, RGB, XYZ, or
luminance values as a matrix to a MATLAB or Mathematica ASCII file or a NumPy binary file.
This is useful when running Mitsuba as simulation step as part of a
larger virtual experiment. It can also come in handy when
verifying parts of the renderer using an automated test suite.
\input{section_rfilters}
